\documentclass[11pt,a4paper]{article}
\usepackage[T1]{fontenc}
\usepackage{lmodern}
\usepackage[margin=27mm]{geometry}
\usepackage{amsmath,amssymb,amsthm,mathtools,booktabs,array,microtype}
\usepackage[hidelinks]{hyperref}
\usepackage{enumitem}
\usepackage{algorithm,algorithmic}
\renewcommand{\algorithmiccomment}[1]{{\footnotesize\itshape // #1}}
\setlist{itemsep=3pt,topsep=5pt}
\newtheorem{theorem}{Theorem}[section]
\newtheorem{lemma}[theorem]{Lemma}
\newtheorem{proposition}[theorem]{Proposition}
\newtheorem{corollary}[theorem]{Corollary}
\theoremstyle{definition}
\newtheorem{definition}[theorem]{Definition}
\newtheorem{question}[theorem]{Open question}
\newtheorem{target}[theorem]{Desired theorem}
\theoremstyle{remark}
\newtheorem{remark}[theorem]{Remark}
\newcommand{\Ostar}{O^{*}}
\newcommand{\E}{\mathbb E}
\newcommand{\Prb}{\mathbb P}
\newcommand{\OPT}{\operatorname{OPT}}
\newcommand{\supp}{\operatorname{supp}}
\newcommand{\ind}{\mathbf 1}
\newcommand{\bits}{\{0,1\}}
\newcommand{\sign}{\operatorname{sgn}}
\DeclareMathOperator{\rank}{rank}
\emergencystretch=2em
\clubpenalty=10000
\widowpenalty=10000
\hypersetup{pdftitle={New Upper Bounds for Dense Max 2-SAT by Accumulating Pruning Constraints}}
\title{New Upper Bounds for Dense Max 2-SAT\\[3pt]
by Accumulating Pruning Constraints}
\author{}
\date{12 October 2026\\\small First research draft}

\begin{document}
\maketitle

\begin{abstract}
We study exact algorithms for classes of Max 2-SAT instances whose clause
density can grow with the number of variables. The objective is a running
time $\Ostar(2^{(1-\gamma)n})$, for a constant $\gamma>0$ independent of
density, using polynomial space. We organize the analysis around comparisons
between assignments: a branch needs to be explored only where it improves
on an alternative that has already been covered. Such comparisons produce
hard constraints whose accumulation can reduce the search space.
We give three principal families: formulas with a persistent relative literal
bias; independently signed clauses on any fixed almost-regular interaction
graph; and formulas with repeated interactions across a large cut, including
literal bicliques with a linear number of defective cross pairs. We also
prove a compression bound for many bounded-size blocks that admit two
assignments with identical, or nearly identical, outside effects.
Conditional comparisons connect these results to bounded caching and
generalized clause learning. We distinguish these proved mechanisms from
the open structural question: whether failure of constraint accumulation
necessarily exposes redundancy that can be exploited at comparable cost.
\end{abstract}

\paragraph{Status of this draft.}
This is a paper draft with an explicit research agenda, not a completed
novelty or independent proof audit. The core statements below include proofs.
The main presentation starts from literal gaps, random formulas and
Max-Cut on regular graph families and the stochastic block model.
Short combinatorial routines precede the spectral analysis; hidden
communities additionally use explicit spectral preprocessing. General
accumulation and discovery lemmas, technical membership estimates, and
sharper special-case constants are placed in appendices. The final sections record proposed algorithms and open
implications without treating them as established theorems. Numerical
constants are conservative.

\section{Introduction}

Max 2-SAT asks for a Boolean assignment satisfying the largest number of
clauses of a formula with at most two literals per clause. Exhaustive search
takes $\Ostar(2^n)$ time and polynomial space. Our interest is in structural
conditions that give a fixed exponential saving even when the number of
binary clauses is superlinear in $n$.

The space requirement is central. Williams gave an algorithm with running
time $\Ostar(2^{\omega n/3})$ and exponential space, where $\omega$ is the
matrix multiplication exponent~\cite{Williams}. Thus the question here is
not whether unrestricted-space exact algorithms can beat $2^n$.
Kulikov and Kutzkov obtained polynomial-space algorithms with a fixed
exponential saving for every fixed clause density, using information from
one branch to constrain another~\cite{KK}. Degree-dependent algorithms and
separator methods give further improvements~\cite{GK}. A bound whose saving
vanishes as the degree grows does not, by itself, give a fixed base below two
on a family of growing density.

The central observation is elementary. If an alternative assignment is at
least as good for a particular completion, there is no need to explore both
assignments on that completion. The difficulty is to make enough such
comparisons simultaneously useful, discover them at an affordable cost,
and enumerate the surviving search space efficiently.

We pursue two ways in which comparisons become useful. First, comparisons
can restrict different aspects of a completion, so their combined effect is
strong even when each individual constraint is weak. Second, many assignments
can have the same effect on the remaining variables. In this case, comparing
the assignments directly can replace repeated search by a smaller collection
of representatives. These mechanisms can coexist.

\subsection{Input families and the common idea}

We first present three kinds of input families, each with a direct
description. Their precise qualifications are part of the corresponding
theorems.
\begin{enumerate}
\item \textbf{Literal gaps.} A linear fraction of variables have a fixed
relative literal bias and the degrees have bounded spread. The dense
specialization includes a linear number of variables with linear gaps.
Neighborhood correlations are arbitrary.
\item \textbf{Random formulas.} Independently random literal signs work
on every fixed graph, including several clauses per pair and independently
fixed units. Further models allow vanishing clause probabilities, parity
bias, and units chosen after the binary clauses, above a logarithmic degree
threshold.
\item \textbf{Max-Cut families.} Regular graphs with small nontrivial
eigenvalues admit a combinatorial count-and-compare algorithm. In the
stochastic block model, a fixed number of hidden communities and comparable
edge variances give a fixed saving whenever $n v_{\min}\ge C\log n$.
The community labels and edge probabilities need not be supplied.
\end{enumerate}
The common operation is to compare a branch with the opposite value and
accumulate hard constraints on the remaining assignment. Literal bias,
random signs and regular-graph Max-Cut use combinatorial routines with
explicit pseudocode. Spectral estimates belong to their analysis. Hidden
communities also require spectral preprocessing in our current algorithm;
when their partition and probabilities are supplied, group counts suffice.

We then study a complementary mechanism: repeated effects on the remaining
variables permit compression into representatives. This yields bounds for
literal bicliques, a linear number of defective cross pairs, repeated
interaction profiles, and bounded blocks. These mechanisms can coexist;
they are not an exhaustive classification. Compression across a cut does
not subsume literal bias, and an exact literal biclique can coexist with
almost balanced literal degrees; see Section~\ref{sec:separation}.
Section~\ref{sec:applications} gives a modeling interpretation of literal
bias in terms of asymmetric pairwise conflicts, together with examples and
the qualifications needed before claiming practical relevance.

\begin{table}[t]
\centering\small
\begin{tabular}{@{}p{0.26\textwidth}p{0.35\textwidth}p{0.30\textwidth}@{}}
\toprule
Family & Assumptions to keep visible & Exponential bound \\
\midrule
Relative literal bias & Canonical pair form; units allowed; fixed
$\varepsilon,\eta,R$ & $2^{(1-\varepsilon^2\eta^2/(3000R))n}$ \\
Weighted random signs & Any fixed graph; bounded independent pair laws;
units fixed independently & $(2e^{-c_0/2})^n$ with exponentially high probability \\
Regular-graph Max-Cut & Regular; $\lambda(G)\le\lambda_0\sqrt d$;
$d\le(1-\varepsilon)n$; fields arbitrary & $c^n$, $c(\lambda_0,\varepsilon)<2$ \\
Stochastic block model & Fixed $k,\chi$; comparable edge variances;
$n v_{\min}\ge C\log n$; groups hidden & $c^n$, $c(k,\chi)<2$, with probability $1-o(1)$ \\
Literal biclique & Sides of sizes $s,t$; exact prescribed cross clauses &
$2^{n-\min(s,t)}$ when supplied \\
Bounded blocks & $q$ disjoint blocks, size at most $b$; at most $r$
exceptional outside variables each &
$2^n(1-2^{-b})^{q/(2r+1)}$ \\
\bottomrule
\end{tabular}
\caption{Core bounds proved in the draft. All suppress polynomial factors
and use polynomial space. Structural detection costs are treated separately;
the random-instance probabilities are over the input, not the algorithm.}
\label{tab:main}
\end{table}

\subsection{Related work and scope}

The comparison interpretation is closely related to the clause-learning
approach of Kulikov and Kutzkov~\cite{KK}; it should not be presented as a
new general principle. Local maxima are also a classical object in exact
Max 2-SAT algorithms. Hastings relates their number to exact algorithms and
obtains further degree-dependent improvements under additional
assumptions~\cite{Hastings}. Formula caching has a theoretical literature,
including the proof-complexity analysis of Beame, Impagliazzo, Pitassi and
Segerlind~\cite{BIPS}. Our question is which concrete dense classes admit
density-independent savings from these mechanisms with polynomial space.
A full comparison of the class statements, rather than only the algorithmic
ideas, remains part of the publication audit.

\section{Notation, simplification, and pruning constraints}\label{sec:framework}

Let $F$ be an unweighted Max 2-SAT formula on a set $V$ of $n$ variables.
In this section and in the literal-gap theorem, the input is \emph{simple}:
no clause occurs twice, and no variable occurs twice in a clause.
We write $\mathbf v$ for a variable and $v,\neg v$ for its two literals.
The numbers $d_1^F(v)$ and $d_2^F(v)$ count the unit and binary clauses
containing $v$, respectively. We omit the superscript when the formula is
clear. Unit clauses created by simplification retain their multiplicities.
The number of binary clauses is $m_2$, and their density is $\rho=m_2/n$.
The notation $\Ostar$ suppresses factors polynomial in the input size.

A \emph{partial assignment} $\alpha$ assigns values to a subset of $V$.
A variable outside that subset is \emph{unassigned}; a complete assignment
extending $\alpha$ is a \emph{completion} of $\alpha$.
For a complete assignment $x$, let $F(x)$ be its number of satisfied clauses,
and let $\OPT(F)=\max_x F(x)$.
Restricting $F$ by $\alpha$ removes satisfied clauses and false literals.
We record the removed satisfied clauses in a separate score $s$.
Thus, for every completion, its original score equals $s$ plus its score in
the residual formula. An empty soft clause contributes zero and is discarded.

The input clauses are soft: an optimal assignment may leave some of them
unsatisfied. A pruning constraint is an additional restriction used to
reduce the search space. We write $C$ for one such constraint and
$\mathcal C$ for a collection of them. Our constraints have the form
$\sum_i\ell_i\le b$, where every $\ell_i$ is a Boolean literal and repeated
occurrences are counted separately.

\subsection{Simplification rules and preprocessing}

\begin{lemma}[Unit clause cancellation]\label{lem:unit_canc}
Removing one copy of $(v)$ and one copy of $(\neg v)$ decreases the score
of every complete assignment by exactly one.
\end{lemma}
\begin{proof}
Exactly one of the two unit clauses is satisfied by every assignment.
\end{proof}

\begin{lemma}[Simplifying two clauses on the same variable pair]\label{lem:simplify}
Let $a,b$ be literals on distinct variables. Replacing
$(a\vee b),(a\vee\neg b)$ by the unit clause $(a)$ decreases the score of
every complete assignment by exactly one.
\end{lemma}
\begin{proof}
If $a=1$, the two original clauses contribute two and the new unit contributes
one. If $a=0$, exactly one original clause is satisfied and the new unit
contributes zero. In both cases the difference is one.
\end{proof}

\begin{lemma}[Unit clause domination]\label{lem:unit_dom}
Suppose $d_1(\neg v)=0$ and $d_1(v)\ge d_2(\neg v)$. For every fixed
assignment of the other variables, setting $v=1$ satisfies at least as many
clauses as setting $v=0$.
\end{lemma}
\begin{proof}
Changing $v$ from zero to one gains all $d_1(v)$ positive units. Clauses
containing $v$ cannot lose satisfaction. The only possible losses are among
the $d_2(\neg v)$ binary clauses containing $\neg v$, since there are no
units $(\neg v)$. The net gain is therefore at least
$d_1(v)-d_2(\neg v)\ge0$.
\end{proof}

In particular, a pure literal, whose opposite has no occurrences, can be
set to one. When pruning constraints are present, such an assignment must
also respect them. The procedure below uses domination only when the
literal being set to one occurs in no current constraint.

\begin{algorithm}[H]
\caption{\textnormal{\textsc{Preprocess}}$(F)$}\label{alg:preprocess}
\begin{algorithmic}[1]
\STATE $s\gets0$
\WHILE{$F$ contains $(a\vee b),(a\vee\neg b)$ on distinct variables}
  \STATE Replace these clauses by $(a)$ using Lemma~\ref{lem:simplify}.
  \STATE $s\gets s+1$
\ENDWHILE
\FOR{each variable $\mathbf v$}
  \STATE $t\gets\min\{d_1(v),d_1(\neg v)\}$
  \STATE Remove $t$ copies of each opposite unit; set $s\gets s+t$.
\ENDFOR
\STATE \textbf{return} $(F,s)$
\end{algorithmic}
\end{algorithm}

The score identities justify \textnormal{\textsc{Preprocess}} for every assignment,
not only for an optimum. After it terminates, a variable pair occurs in
at most two binary clauses. If there are two, their literals are
complementary in both positions: otherwise the two clauses would share a
literal and the loop could apply again. We call this form \emph{canonical}.
Assigning variables creates only units, so the binary clauses do not need
to be preprocessed again at later nodes.

\subsection{Constructing and enforcing a pruning constraint}

\begin{lemma}[Pruning constraints from comparison]\label{lem:pruning}
In a residual formula $F$, let $v$ be either literal of an unassigned
variable. List the partner literals in clauses containing $v$ as
$x_1,\ldots,x_{d_2(v)}$, and those in clauses containing $\neg v$ as
$y_1,\ldots,y_{d_2(\neg v)}$. Define
\begin{equation}
 C:\quad \sum_{i=1}^{d_2(v)}\neg x_i+
             \sum_{j=1}^{d_2(\neg v)}y_j
       \le d_2(\neg v)-d_1(v)+d_1(\neg v)-1.
 \label{eq:gap_pruning}
\end{equation}
For every complete assignment of the other variables, $C$ holds if and
only if setting $v=0$ has strictly larger score than setting $v=1$.
\end{lemma}
\begin{proof}
Keep all other variables fixed. A unit $(v)$ contributes one to the gain
of $v=1$ over $v=0$, and a unit $(\neg v)$ contributes minus one.
A clause $(v\vee x_i)$ contributes $1-x_i=\neg x_i$ to this gain.
A clause $(\neg v\vee y_j)$ contributes $-(1-y_j)=y_j-1$.
All clauses not containing the variable of $v$ have the same score in both
assignments. Their contributions cancel. The exact gain is therefore
\[
 d_1(v)-d_1(\neg v)-d_2(\neg v)
       +\sum_i\neg x_i+\sum_j y_j.
\]
Setting $v=0$ is strictly better precisely when this integer is at most
$-1$, which is the stated constraint.
\end{proof}

To restrict a constraint after an assignment, remove every false literal
occurrence. Remove every true literal occurrence and decrease its bound by
one. Also, a complementary pair $w,\neg w$ can be removed while decreasing
the bound by one, since $w+\neg w=1$. A negative bound proves infeasibility.
After cancellation, if a literal occurs more often than the bound, it must
be false in every satisfying completion.

\begin{lemma}[Infeasibility by summation]\label{lem:constraint_sum}
Let $\mathcal C$ consist of constraints $C:\sum_{i=1}^{r_C}\ell_i\le b_C$,
with literal occurrences counted separately. Add the constraints, and
repeatedly remove complementary literal pairs from the sum, decreasing
its bound by one per pair. If the resulting bound is negative, no
assignment satisfies all constraints in $\mathcal C$.
\end{lemma}
\begin{proof}
An assignment satisfying every constraint also satisfies their sum.
Cancelling a pair preserves the inequality because its contribution is
exactly one for every assignment. The remaining left-hand side is
nonnegative, so it cannot satisfy a negative bound.
\end{proof}

This is a sufficient infeasibility test. Passing it does not imply that
the collection has a common satisfying assignment. The summed constraint
is temporary; it is not added to $\mathcal C$.

In \textnormal{\textsc{Simplify}}, the instruction \emph{assign $v=p$} records the
corresponding variable value in $\alpha$, updates the residual soft formula
and score $s$, and substitutes the value in every constraint. All choices
within preprocessing and simplification follow fixed orders.

\begin{algorithm}[H]
\caption{\textnormal{\textsc{Simplify}}$(F,\mathcal C)$}\label{alg:simplify}
\begin{algorithmic}[1]
\STATE $s\gets0$; $\alpha\gets\emptyset$
\WHILE{true}
  \STATE Cancel opposite units by Lemma~\ref{lem:unit_canc}, adding their constant to $s$.
  \STATE Cancel complementary literal pairs within each $C\in\mathcal C$.
  \STATE Let $r_C$ be its remaining number of occurrences; remove $C$ if $b_C\ge r_C$.
  \IF{some $b_C<0$, or Lemma~\ref{lem:constraint_sum} proves infeasibility}
    \STATE \textbf{return} \textsc{Infeasible}
  \ENDIF
  \IF{some literal $v$ occurs more than $b_C$ times in some $C\in\mathcal C$}
    \STATE Assign $v=0$, updating $F,\mathcal C,s,\alpha$.
  \ELSIF{an unassigned literal $v$ satisfies $d_1(\neg v)=0$, $d_1(v)\ge d_2(\neg v)$,
          and occurs in no $C\in\mathcal C$}
    \STATE Assign $v=1$, updating $F,\mathcal C,s,\alpha$.
  \ELSE
    \STATE \textbf{return} $(F,\mathcal C,s,\alpha)$
  \ENDIF
\ENDWHILE
\end{algorithmic}
\end{algorithm}

The last assignment rule preserves feasibility of the current constraints:
setting $v=1$ cannot increase any left-hand side, because only $\neg v$ can
occur there. Lemma~\ref{lem:unit_dom} shows that the soft score does not
decrease. Thus this rule preserves the best score subject to the current
constraints. An empty collection has sum $0\le0$.

\section{The literal gap algorithm and its analysis}\label{sec:bias}

\subsection{Intuition}\label{sec:bias_intuition}

Suppose a linear number of variables have a substantial difference between
the frequencies of their two literals. For the initial intuition, ignore
unit clauses. If one literal occurs much more often than its negation,
setting it to one satisfies more clauses on average over assignments to
the other variables. The other value may still be better for particular
assignments, so we cannot simply discard it. Instead, the pruning constraint
of Lemma~\ref{lem:pruning} describes exactly when that other value is strictly
better. A large gap makes this a substantial requirement on the neighboring
literals. Here a large gap means a constant fraction of the variable's
degree; the degree itself need not be linear in $n$.

Consider a path in the branching tree on which each branching value is
chosen with probability $1/2$. This random path is only a way of examining
the search tree; the algorithm is deterministic. Whenever we branch on a
variable that still has a substantial gap, the nonpreferred value adds a
constraint whose bound is well below the expected number of true literal
occurrences. For example, $\sum_{i=1}^r\ell_i\le r/4$ allows only half the
expected number $r/2$ of true occurrences under a uniform assignment.
Each occurrence contributes either zero or one; its expectation is one
half. Taking the nonpreferred value at many such steps therefore accumulates
requirements that the remaining assignment must meet.

The constraints need not have disjoint neighborhoods, and may even be
identical. The question is whether a few shared variables can satisfy many
requirements at once. Each branching variable creates at most one
constraint, using its still-present binary clauses. Consequently, a fixed
variable can occur across the accumulated constraints at most as many times
as it occurs in binary clauses after preprocessing. Small degrees limit
overlap in this way. Larger degrees allow more overlap, but a relative gap
at a branching variable also gives a stronger constraint. The theorem
controls the concentration of these occurrence counts so that a small set
of variables cannot account for all the accumulated requirements.

This suggests the source of the savings: on most paths, enough restrictive
constraints accumulate for their sum to become infeasible while many
variables remain unassigned. The infeasibility test of
Lemma~\ref{lem:constraint_sum} then rejects the whole remaining subtree.
Simplification provides a second source of savings: when it assigns a
variable without branching, only one value needs to be explored. Both
effects are useful. The analysis must also account for gaps changing as
neighboring variables are assigned; it does not assume that every initial
gap persists along every path.

\subsection{The algorithm}\label{sec:bias_algorithm}

The algorithm first applies \textnormal{\textsc{Preprocess}}, then searches
recursively. At each node, \textnormal{\textsc{Simplify}} applies reduction
rules and checks the accumulated constraints. If branching is still
necessary, the algorithm explores both values of the next variable and
adds a pruning constraint only to its nonpreferred child. The procedures
from Section~\ref{sec:framework} perform the simplification and propagation;
we only need to specify the variable order and the preferred value.

To include unit clauses in the preference calculation, define the
\emph{effective gap} of a variable $\mathbf v$ in a residual formula $H$ by
\begin{equation}
 g_{\mathbf v}(H)=d_2^H(v)-d_2^H(\neg v)
                    +2\bigl(d_1^H(v)-d_1^H(\neg v)\bigr).
 \label{eq:effectivegap}
\end{equation}
Without units this is the difference of the two literal frequencies.
A unit contributes its full preference irrespective of the other variables,
whereas a binary clause contributes its preference only when its partner
literal is false. Under uniformly assigned remaining variables, the latter
event has probability one half. Thus $g_{\mathbf v}(H)/2$ is the expected
score gain from setting $v=1$ rather than $v=0$.

Fix an order on the input variables. At every branching node, choose the
first unassigned variable $\mathbf u$ in this order. Its \emph{preferred
literal} is
\begin{equation}
 v=\begin{cases}
 u,&d_2(u)+2d_1(u)\ge d_2(\neg u)+2d_1(\neg u),\\
 \neg u,&\text{otherwise}.
 \end{cases}
 \label{eq:bias_preference}
\end{equation}
All counts here are taken in the current formula, after
\textnormal{\textsc{Simplify}}. Setting $v=1$ is the preferred choice;
ties prefer the positive literal. The algorithm calculates this preference
directly from clause counts, without sampling assignments. The fixed order
does not have to put variables with large gaps first.

The global variables $M$ and $\alpha^\star$ store the best complete score
and a corresponding assignment. The parameters $F,\mathcal C,\alpha,s$
belong to the current node: its residual formula, pruning constraints,
partial assignment, and accumulated satisfied-clause contribution.
Variables remain in the unassigned set even if they disappear from $F$.

\begin{algorithm}[H]
\caption{\textnormal{\textsc{LiteralGap}}$(F)$}\label{alg:literal_gap}
\begin{algorithmic}[1]
\STATE Fix an order on the input variables.
\STATE $M\gets-\infty$; $\alpha^\star\gets\emptyset$
\COMMENT{Global best score and complete assignment.}
\STATE $(F,s)\gets\textnormal{\textsc{Preprocess}}(F)$
\STATE $\textnormal{\textsc{GapSearch}}(F,\emptyset,\emptyset,s)$
\STATE \textbf{return} $(M,\alpha^\star)$
\end{algorithmic}
\end{algorithm}

\begin{algorithm}[H]
\caption{\textnormal{\textsc{GapSearch}}$(F,\mathcal C,\alpha,s)$}
\label{alg:gap_search}\label{alg:bias}
\begin{algorithmic}[1]
\STATE $R\gets\textnormal{\textsc{Simplify}}(F,\mathcal C)$
\IF{$R=\textsc{Infeasible}$}
  \STATE \textbf{return}
\ENDIF
\STATE $(F,\mathcal C,t,\beta)\gets R$; $s\gets s+t$; $\alpha\gets\alpha\cup\beta$
\COMMENT{Include the assignments and score obtained by simplification.}
\IF{all input variables are assigned}
  \STATE If $s>M$, set $M\gets s$ and $\alpha^\star\gets\alpha$.
  \STATE \textbf{return}
\ENDIF
\STATE Let $\mathbf u$ be the first unassigned variable in the fixed order.
\STATE Choose its preferred literal $v$ by~\eqref{eq:bias_preference}.
\STATE Construct $C$ from $F$ and $v$ by Lemma~\ref{lem:pruning}.
\FOR{$p\in[1,0]$}
  \STATE Make separate copies $(F_p,\mathcal C_p,\alpha_p,s_p)$ of the current state.
  \IF{$p=0$}
    \STATE $\mathcal C_p\gets\mathcal C_p\cup\{C\}$
  \ENDIF
  \STATE Assign $v=p$, updating the copied formula, constraints, assignment, and score.
  \STATE $\textnormal{\textsc{GapSearch}}(F_p,\mathcal C_p,\alpha_p,s_p)$
\ENDFOR
\end{algorithmic}
\end{algorithm}

The algorithm does not need to identify a set of variables with large gaps
or know the parameters in the running-time theorem. It applies the same
rules to every simple input formula. The theorem's assumptions determine
when these rules give an exponential saving.

\subsection{Correctness}

A single pruning constraint follows from comparing the two values of one
variable. We must also show that many such comparisons, together with
simplification, preserve at least one optimum. The following lemma does
this by choosing the earliest optimum in a fixed ordering of the full
decision tree and showing that none of the rules removes its path. This
establishes correctness independently of the assumptions used to bound
the running time.

\begin{lemma}[Correctness of the literal gap algorithm]\label{lem:gap_correctness}
For every simple input formula, \textnormal{\textsc{LiteralGap}} returns its optimum
score and an assignment attaining that score.
\end{lemma}
\begin{proof}
Preprocessing and unit cancellation preserve every assignment's score
when their recorded constants are included. It remains to show that the
search retains an optimum despite the pruning constraints and assignments
made by \textnormal{\textsc{Simplify}}.

For this purpose, represent each assignment made by \textnormal{\textsc{Simplify}} as a
decision with one retained child. Add the other child temporarily, and also
restore subtrees rejected by infeasibility tests. Complete these restored
subtrees by arbitrary binary decisions until every variable is assigned.
This gives a full decision tree whose leaves are all complete assignments.
At an ordinary branching node, place the preferred child first. At an
assignment made by \textnormal{\textsc{Simplify}}, place the retained value first.
This defines an order on all complete assignments: list the leaves from
left to right. Choose the earliest optimum $x^*$ in this order.

We show that the actual algorithm retains the path of $x^*$. Suppose it
has reached a node on that path. At an ordinary branch where $x^*$ takes
the nonpreferred value $v=0$, change only this variable to $v=1$. The
resulting assignment lies in the earlier child. Its score must be strictly
smaller than that of $x^*$: a larger score contradicts optimality, and an
equal score contradicts the choice of the earliest optimum. By
Lemma~\ref{lem:pruning}, $x^*$ therefore satisfies the newly added constraint.
Inductively, it satisfies every constraint generated on its path.

Consequently, neither an individual constraint nor their sum can certify
infeasibility on this path. An assignment forced by a constraint agrees
with $x^*$. If a domination rule instead assigns $v=1$ and $x^*$ had $v=0$,
Lemma~\ref{lem:unit_dom} would give an assignment with at least the same score
in the earlier child, again contradicting the choice of $x^*$.
Removing an already satisfied constraint and cancelling complementary
literals preserve satisfaction. Thus every operation retains this optimum.

The algorithm eventually visits its complete assignment and records its
original score. Every other score recorded also belongs to a complete
assignment of the original formula, so none exceeds the optimum.
\end{proof}

\subsection{Running-time guarantee}

The following theorem makes the two requirements from the intuition
precise: many variables have a gap proportional to their binary degree,
and the total occurrence counts are not too concentrated on a few
variables. The first supplies strong pruning constraints. The second
bounds how much shared variables can affect their sum and how much bias
can remain in a small unassigned set. All these conditions concern the
formula immediately after preprocessing.

\begin{theorem}[Relative literal bias]\label{thm:bias}
Fix constants $\varepsilon,\eta\in(0,1]$ and $R\ge1$.
Let $F$ be a simple unweighted Max 2-SAT formula on a set $V$ of $n$ variables, and let
$(F_0,s_0)=\textnormal{\textsc{Preprocess}}(F)$. For each variable $\mathbf v$, define
\[
 D_{\mathbf v}=d_2^{F_0}(v)+d_2^{F_0}(\neg v),\qquad
 u_{\mathbf v}=d_1^{F_0}(v)+d_1^{F_0}(\neg v),
\]
\[
 g_{\mathbf v}=d_2^{F_0}(v)-d_2^{F_0}(\neg v)
                    +2\bigl(d_1^{F_0}(v)-d_1^{F_0}(\neg v)\bigr).
\]
Suppose there is a set $S\subseteq V$ with $|S|\ge\eta n$ such that
$|g_{\mathbf v}|\ge\varepsilon D_{\mathbf v}$ for every $\mathbf v\in S$.
Writing $d=|S|^{-1}\sum_{\mathbf v\in S}D_{\mathbf v}$, suppose also that
\begin{equation}
 d\ge1,\qquad
 \sum_{\mathbf v\in V}(D_{\mathbf v}+u_{\mathbf v})^2\le Rnd^2.
 \label{eq:biasregularity}
\end{equation}
Then an optimal assignment for $F$ can be computed deterministically in
polynomial space and time $\Ostar(2^{(1-\gamma)n})$, where
$\gamma=\varepsilon^2\eta^2/(3000R)$.
\end{theorem}

The average degree $d$ may grow with $n$. For fixed
$\varepsilon,\eta,R$, the saving $\gamma$ remains a positive constant
independent of $d$. Condition~\eqref{eq:biasregularity} allows unequal
degrees; it controls their squared sum rather than imposing a common
degree or requiring disjoint neighborhoods.

\subsection{Proof of the running-time bound}

The proof measures the fraction of paths that survive to a depth at which
a small constant fraction of the variables is still unassigned. A path
includes assignments made by \textnormal{\textsc{Simplify}}, as well as
ordinary branching decisions. The argument proceeds as follows.
\begin{enumerate}
\item Paths with many assignments made by simplification are already few:
at each such step only one of the two values is retained. This is the
role of Lemma~\ref{lem:forced_path} in Step~2 below.
\item On the other paths, a substantial part of the initial total gap
must be encountered at ordinary branching steps. A small set of variables
cannot contain most of it, and random changes caused by assigning
neighbors are unlikely to remove it. Step~3 combines the degree bound
with Lemmas~\ref{lem:gap_changes} and~\ref{lem:bias_probabilities}.
\item These branching steps generate constraints whose combined
requirements are unlikely to be met by the assigned values and the few
remaining variables. Step~4 uses the exact sum identity of
Lemma~\ref{lem:bias_sum} and the remaining estimates of
Lemma~\ref{lem:bias_probabilities} to prove rejection by the sum test.
\item An exponentially small surviving fraction at the chosen depth
gives the search-tree bound in Step~5. Shallower levels have few nodes
because they are still a linear distance from the leaves.
\end{enumerate}
The main proof gives these steps and their quantitative bounds. The
supporting lemmas, including the martingale calculations, follow it.

\begin{proof}[Proof of Theorem~\ref{thm:bias}]
We prove the bound by analyzing \textnormal{\textsc{LiteralGap}}
(Algorithms~\ref{alg:literal_gap} and~\ref{alg:gap_search}).
Its correctness follows from Lemma~\ref{lem:gap_correctness}.
Depth-first search stores a polynomial-size state at each of at most $n$
assignment levels, so its space use is polynomial. Each call to
\textnormal{\textsc{Simplify}} performs at most $n$ assignments and polynomial work
between assignments. We therefore need to bound the size of the search tree.

\par\medskip\noindent\textbf{1. Choose an analysis depth.} 
Put $\theta=\varepsilon\eta$, $B_{\mathbf v}=D_{\mathbf v}+u_{\mathbf v}$,
and
\[
 \delta=\frac{\theta^2}{256R},\qquad
 T=\lceil(1-\delta)n\rceil.
\]
Expose assignments one variable at a time, including those made within
\textnormal{\textsc{Simplify}}. Call the latter \emph{forced steps}; call decisions with
two recursive children \emph{branching steps}. Restore missing children
and pruned subtrees as in the correctness proof, obtaining a full binary
tree. Follow a random path in this tree, choosing each value by an
independent fair coin. After a path is pruned, continue it arbitrarily by
selecting unassigned variables, using the same preferred-value and
constraint-generation rules but ignoring rejection tests. This continuation
is used only in the analysis.

Let $\pi_t$ be the probability that the random path is retained through
$t$ assignments, including all infeasibility checks at that level. Each
length-$t$ path in the full tree has probability $2^{-t}$, so the number of
retained states at that level is $2^t\pi_t$. Rejected children and the work
inside simplification contribute only polynomial factors. It suffices to
bound $\max_{0\le t\le n}2^t\pi_t$.

\par\medskip\noindent\textbf{2. Set aside paths with many forced steps.} 
Let $R_T$ be the number of forced steps among the first $T$ assignments.
By Lemma~\ref{lem:forced_path},
\[
 \Prb(\text{the path survives to }T\text{ and }R_T>\delta n)
 \le 2^{-\delta n}.
\]
On the remaining paths, at most $\delta n$ variables have been assigned
by simplification. We show that ordinary branching then creates enough
constraints to reject almost all such paths by depth $T$.

\par\medskip\noindent\textbf{3. Account for the initial gaps.} 
We first bound how much gap can be confined to a small set of variables.
We will use this bound both for variables still unassigned at depth $T$
and for those assigned by simplification.
For every set $W\subseteq V$ of size at most $\delta n$, the
Cauchy--Schwarz inequality gives
\begin{align*}
 \sum_{\mathbf v\in W}B_{\mathbf v}
 &\le \left(\sum_{\mathbf v\in W}1^2\right)^{1/2}
       \left(\sum_{\mathbf v\in W}B_{\mathbf v}^2\right)^{1/2}\\
 &\le\sqrt{\delta n}\sqrt{Rnd^2}
   =\sqrt{\delta R}\,nd=\frac{\theta nd}{16}.
\end{align*}
The total initial absolute gap of variables in $S$ is at least
\[
 \Phi_0:=\sum_{\mathbf v\in S}|g_{\mathbf v}|
 \ge\varepsilon\sum_{\mathbf v\in S}D_{\mathbf v}
 =\varepsilon|S|d\ge\theta nd.
\]
Let $\Phi_T$ be their total absolute gap among the variables still
unassigned at depth $T$. There are at most $\delta n$ unassigned variables.
Lemma~\ref{lem:gap_changes} bounds every residual absolute gap by
$2B_{\mathbf v}$, so $\Phi_T\le\theta nd/8$.
The same argument bounds the sum of the gaps of variables in $S$ removed
at forced steps by $\theta nd/8$ whenever $R_T\le\delta n$.

We must also account for the changes in a variable's gap before it is
selected. Let $H_T$ sum the absolute gaps of variables in $S$ at the moments they
are assigned, including forced steps. Let $G_T$ be the corresponding sum
over branching steps only. The first estimate of
Lemma~\ref{lem:bias_probabilities}, with $a=\theta nd/4$, gives
\[
 \Prb\bigl(H_T<\Phi_0-\Phi_T-\theta nd/4\bigr)
 \le e^{-\theta^2n/(128R)}.
\]
Outside this exceptional event, on every path with $R_T\le\delta n$,
\[
 G_T\ge \theta nd-\frac{\theta nd}{8}
                    -\frac{\theta nd}{4}-\frac{\theta nd}{8}
      =\frac{\theta nd}{2}.
\]

\par\medskip\noindent\textbf{4. Show that the summed constraints are infeasible.} 
We separate the total requirement at constraint creation, the effect of
values assigned later, and the best contribution still available from
unassigned variables. This expresses precisely what the sum test must
compare.
For a constraint $C$, let $r_C^0$ and $b_C^0$ refer to its form when it is
created. Its initial restriction relative to the expected number of true
literal occurrences is $r_C^0/2-b_C^0$.
Let $A_T$ be the sum of these quantities over constraints created by depth
$T$. Let $Q_T$ be the total number of true occurrences subsequently assigned
in these constraints, minus one half the number of occurrences assigned.
Finally, let
\[
 U_T=\frac12\sum_{\mathbf w\text{ unassigned}}
             |m_{\mathbf w}^+-m_{\mathbf w}^-|,
\]
where $m_{\mathbf w}^+$ and $m_{\mathbf w}^-$ count the occurrences of $w$
and $\neg w$ across those constraints. These quantities count occurrences,
not distinct literals.

Lemma~\ref{lem:bias_sum} shows that the minimum left-hand side minus
right-hand side of their sum is
\begin{equation}
 Z_T=A_T+Q_T-U_T.\label{eq:race}
\end{equation}
In particular, a surviving path must satisfy $Z_T\le0$.
The occurrence bound $m_{\mathbf w}^++m_{\mathbf w}^-\le D_{\mathbf w}$
and the Cauchy--Schwarz estimate above give
\[
 U_T\le\frac12\sum_{\mathbf w\text{ unassigned}}D_{\mathbf w}
 \le\frac12\sum_{\mathbf w\text{ unassigned}}B_{\mathbf w}
 \le\frac{\theta nd}{32}.
\]
The remaining estimates of Lemma~\ref{lem:bias_probabilities}, with
$b=c=\theta nd/32$, show that
\[
 A_T\ge G_T/4-\theta nd/32,\qquad
 Q_T\ge-\theta nd/32
\]
except on events with probabilities at most
$e^{-\theta^2n/(4096R)}$ and $e^{-\theta^2n/(512R)}$, respectively.
If none of the three exceptional events occurs and $R_T\le\delta n$,
then
\[
 Z_T\ge\frac{\theta nd}{8}
        -\frac{\theta nd}{32}-\frac{\theta nd}{32}
        -\frac{\theta nd}{32}
      =\frac{\theta nd}{32}>0.
\]
The sum test therefore rejects this path.

\par\medskip\noindent\textbf{5. Convert the survival probability to a time bound.} 
Combining the forced-step bound and the three exceptional probabilities,
\[
 \pi_T\le2^{-\delta n}+3e^{-\theta^2n/(4096R)}
       \le4e^{-\theta^2n/(4096R)}.
\]
For $t\le T$, $2^t\pi_t\le2^T$. For $t\ge T$, survival to depth $t$
requires survival to depth $T$, so $2^t\pi_t\le2^n\pi_T$. Consequently,
\[
 \max_t2^t\pi_t\le
 \max\left\{2^{(1-\delta)n+1},
            4\cdot2^{(1-\theta^2/(4096R\ln2))n}\right\}.
\]
Both exponents give at least the saving
$\theta^2/(3000R)$, since $1/256>1/3000$ and
$1/(4096\ln2)>1/3000$. Polynomial work per state completes the proof.
\end{proof}

\subsection{Supporting lemmas for the random-path analysis}

Throughout this subsection, use the full random decision tree and notation
from the proof. Let $\mathcal U_t$ be the unassigned variables after $t$
steps, and let $\mathcal H_t$ denote the information revealed by those steps.
The next variable and whether the step is forced are determined before its
value is revealed. Write $\mathbf x_t$ for the variable assigned at step $t$.
Let $g_{\mathbf v}(t)$ be the effective gap in the residual soft formula
after $t$ assignments. Opposite-unit cancellation does not change this gap.

The first estimate handles paths with many assignments made by
\textnormal{\textsc{Simplify}}. Each such step requires the random path to
follow a specified value, even though the position and the required value
of the step may depend on the earlier choices. This is the estimate used
to set those paths aside in Step~2 of the main proof.

\begin{lemma}[Forced assignments on a random path]\label{lem:forced_path}
For every real $r\ge0$, the probability that the path survives to depth
$T$ and has more than $r$ forced steps is at most $2^{-r}$.
\end{lemma}
\begin{proof}
Let $I_t$ indicate that all forced assignments encountered up to step $t$
were followed, and let $R_t$ count those steps. Put $W_t=2^{R_t}I_t$.
Initially $W_0=1$. At a branching step, $W_t=W_{t-1}$. At a forced step
with $I_{t-1}=1$, one fair choice doubles $W_{t-1}$ and the other gives zero.
If $I_{t-1}=0$, both give zero. In all cases
$\E[W_t\mid\mathcal H_{t-1}]=W_{t-1}$, and thus $\E W_T=1$.
Survival with $R_T>r$ implies $I_T=1$ and $W_T>2^r$.
Markov's inequality gives the claimed probability bound.
\end{proof}

For the remaining paths we need to relate the initial gaps to the gaps
seen when variables are selected. The next lemma supplies two elementary
facts for Step~3: an occurrence bound on every residual gap, and a
description of how assigning one neighbor changes it. The second fact
will also control gap fluctuations in Lemma~\ref{lem:bias_probabilities}.

\begin{lemma}[Changes in the effective gap]\label{lem:gap_changes}
For every unassigned variable $\mathbf v$,
$|g_{\mathbf v}(t)|\le2(D_{\mathbf v}+u_{\mathbf v})$.
Let $c_{\mathbf x\mathbf v}$ be the number of binary clauses on the pair
$\{\mathbf x,\mathbf v\}$ in $F_0$. When an unassigned neighbor
$\mathbf x$ is assigned a fair value, the change in $g_{\mathbf v}$ is
$+c_{\mathbf x\mathbf v}$ or $-c_{\mathbf x\mathbf v}$ with equal probability.
\end{lemma}
\begin{proof}
Every binary clause incident to $\mathbf v$ either remains binary,
disappears, or becomes a unit. Its absolute contribution to the effective
gap is therefore at most two. Each original uncancelled unit also contributes
at most two. This proves the first assertion.

For a clause $(v\vee x)$, the contribution to $g_{\mathbf v}$ changes from
one to zero when $x=1$, and from one to two when $x=0$. Thus the two
changes are minus one and plus one. Replacing either literal by its
negation reverses the corresponding signs. If two clauses survive on the
same pair after \textnormal{\textsc{Preprocess}}, they are complementary in both
positions, and their changes have the same sign for either assignment of
$\mathbf x$. They therefore add to $\pm c_{\mathbf x\mathbf v}$.
Cancelling opposite units changes neither effective gap.
\end{proof}

Step~4 needs a numerical condition that guarantees rejection by the
algorithm's sum test. The following identity separates the requirements
introduced by the constraints from the contributions of assigned and
unassigned variables. It also justifies retaining all created constraints
in the analysis, even when \textnormal{\textsc{Simplify}} has already
discarded some as satisfied.

\begin{lemma}[The minimum of the summed constraints]\label{lem:bias_sum}
With $A_T,Q_T,U_T$ as defined in the main proof, the minimum difference
between the left- and right-hand sides of the sum is $A_T+Q_T-U_T$.
For each constraint $C$, let $g_C$ be the absolute effective gap of its
branching variable when the constraint is created. Then
\begin{equation}
 r_C^0/2-b_C^0=g_C/2+1,\qquad
 C:\ \sum_{i=1}^{r_C^0}\ell_i\le(r_C^0-g_C)/2-1,
 \label{eq:cardcut}
\end{equation}
Also
$m_{\mathbf w}^++m_{\mathbf w}^-\le D_{\mathbf w}$.
These facts imply that survival to depth $T$ requires $A_T+Q_T-U_T\le0$,
even when \textnormal{\textsc{Simplify}} has removed already satisfied constraints.
\end{lemma}
\begin{proof}
At creation, $r_C^0=d_2(v)+d_2(\neg v)$, and Lemma~\ref{lem:pruning} gives
$b_C^0=d_2(\neg v)-d_1(v)+d_1(\neg v)-1$. Subtraction yields
\[
 r_C^0/2-b_C^0
 =\frac{d_2(v)-d_2(\neg v)+2(d_1(v)-d_1(\neg v))}{2}+1.
\]
The expression in the numerator is nonnegative for the preferred literal
and equals the absolute effective gap of its variable.

For the analysis, keep a copy of every constraint in its creation form.
Let $a_C$ be the number of its occurrences assigned by depth $T$, and let
$j_C$ be the number of those occurrences made true. Its remaining bound
is $b_C^0-j_C$. Summing over constraints, the remaining left-hand side is
\[
 \sum_{\mathbf w\in\mathcal U_T}
       \bigl(m_{\mathbf w}^+w+m_{\mathbf w}^-\neg w\bigr).
\]
For each variable its minimum contribution is
$\min(m_{\mathbf w}^+,m_{\mathbf w}^-)$, obtained by choosing the better
of its two values. Since
$\min(p,q)=(p+q-|p-q|)/2$ and
$\sum_{\mathbf w}(m_{\mathbf w}^++m_{\mathbf w}^-)=\sum_C(r_C^0-a_C)$,
the minimum of the sum minus its bound is
\begin{align*}
 \sum_{\mathbf w}\min(m_{\mathbf w}^+,m_{\mathbf w}^-)-\sum_C(b_C^0-j_C)
 &=\sum_C(r_C^0/2-b_C^0)+\sum_C(j_C-a_C/2)
       -\tfrac12\sum_{\mathbf w}|m_{\mathbf w}^+-m_{\mathbf w}^-|\\
 &=A_T+Q_T-U_T.
\end{align*}
Cancellation in Lemma~\ref{lem:constraint_sum} computes this same minimum.

Each other variable creates at most one constraint, containing exactly the
still-present binary occurrences incident to it. Thus a fixed
$\mathbf w$ contributes at most $\sum_{\mathbf v}c_{\mathbf v\mathbf w}
=D_{\mathbf w}$ occurrences across the constraints. This proves the
occurrence bound without any assumption that neighborhoods are disjoint.

Finally, a constraint discarded by \textnormal{\textsc{Simplify}} has nonpositive
left-hand side minus right-hand side for every later completion of the
current partial assignment. Removing it can only increase the value of
the remaining summed expression, pointwise. A positive minimum for the
sum of all the recorded constraints therefore implies a positive minimum
for the live sum tested by the algorithm. Such a node is rejected.
\end{proof}

The probabilistic estimates must allow each constraint to depend on earlier
branching choices. The next lemma bounds cumulative deviations without
requiring independence between steps: it uses conditional mean zero and
bounds on the size of each change. We apply it in
Lemma~\ref{lem:bias_probabilities}; the main proof only needs the three
resulting estimates.

\begin{lemma}[A bounded-increment martingale estimate]\label{lem:martingale_tail}
Let $M_0=0,M_1,\ldots,M_T$ be random variables, with $M_t$ determined
by the information $\mathcal H_t$ revealed through step $t$, such that
$\E[M_t\mid\mathcal H_{t-1}]=M_{t-1}$ for every $t$.
Suppose $|M_t-M_{t-1}|\le b_t$, where $b_t$ is known from
$\mathcal H_{t-1}$, and $\sum_{t=1}^T b_t^2\le B$ on every path for a
constant $B>0$. Then, for every $a>0$,
\[
 \Prb(M_T\le-a)\le\exp(-a^2/(2B)).
\]
\end{lemma}
\begin{proof}
The displayed conditional-expectation identity is the martingale property:
the expected next value, given everything revealed so far, is the current
value. Hence $Y_t=M_t-M_{t-1}$ has conditional mean zero.
Convexity of the exponential on $[-b_t,b_t]$ gives, for $\lambda>0$,
\[
 \E[e^{-\lambda Y_t}\mid\mathcal H_{t-1}]
 \le\tfrac12(e^{\lambda b_t}+e^{-\lambda b_t})
 \le e^{\lambda^2b_t^2/2}.
\]
For the first inequality, bound the exponential by the line joining its
endpoint values and use $\E[Y_t\mid\mathcal H_{t-1}]=0$.
The second follows by comparing the power series of $\cosh z$ and
$e^{z^2/2}$, using $(2k)!\ge2^k k!$.
It follows by successive conditional expectations that
\[
 \E\exp\left(-\lambda M_T-
                 \frac{\lambda^2}{2}\sum_{t=1}^T b_t^2\right)\le1.
\]
On $M_T\le-a$, the random variable inside this expectation is at least
$\exp(\lambda a-\lambda^2B/2)$. Markov's inequality bounds the probability
by $\exp(-\lambda a+\lambda^2B/2)$. Taking $\lambda=a/B$ proves the claim.
\end{proof}

There are three ways the accumulated constraints could fail to give the
required pruning. Too much initial gap could disappear before its variable
is selected; branching could introduce too little restriction relative
to the gaps it encounters; or later assigned values could be unusually
favorable to the constraints. The following estimates bound these three
events in that order. Step~3 uses the first estimate, and Step~4 uses the
other two together with Lemma~\ref{lem:bias_sum}.

\begin{lemma}[Three estimates for a random path]\label{lem:bias_probabilities}
Assume $d\ge1$ and $\sum_{\mathbf v}B_{\mathbf v}^2\le Rnd^2$.
For every $a,b,c>0$, the quantities in the main proof satisfy
\begin{align*}
 \Prb(H_T<\Phi_0-\Phi_T-a)&\le e^{-a^2/(8Rnd^2)},\\
 \Prb(A_T<G_T/4-b)&\le e^{-b^2/(4Rnd^2)},\\
 \Prb(Q_T<-c)&\le e^{-2c^2/(Rnd^2)}.
\end{align*}
\end{lemma}
\begin{proof}
\par\medskip\noindent\textbf{The total gap encountered.} 
At step $t$, let $h_t=|g_{\mathbf x_t}(t-1)|$ if $\mathbf x_t\in S$ and
$h_t=0$ otherwise. Removing the selected variable reduces the potential
$\Phi$ by $h_t$. Let $Y_t$ be the change in the total absolute gap of the
other unassigned variables of $S$ caused by assigning its value. Thus
\[
 \Phi_t=\Phi_{t-1}-h_t+Y_t.
\]
For a remaining variable with current signed gap $g$ and pair multiplicity
$q$, Lemma~\ref{lem:gap_changes} shows that its expected new absolute gap is
$(|g+q|+|g-q|)/2\ge|g|$, by the triangle inequality.
Adding these expectations does not require the changes for different
neighbors to be independent. Consequently
$z_t:=\E[Y_t\mid\mathcal H_{t-1}]\ge0$.

Define
\[
 M_t=\sum_{j=1}^t(Y_j-z_j),\qquad M_0=0.
\]
The quantities $z_t$ are determined before the next value is revealed, and
$\E[Y_t-z_t\mid\mathcal H_{t-1}]=0$. Therefore
$\E[M_t\mid\mathcal H_{t-1}]=M_{t-1}$, which proves that $M$ is a martingale.
Each absolute gap changes by at most its pair multiplicity, so
$|Y_t|\le\sum_{\mathbf v}c_{\mathbf x_t\mathbf v}=D_{\mathbf x_t}$ and
$|Y_t-z_t|\le2D_{\mathbf x_t}$. Each variable is assigned once; hence
the sum of the squared increment bounds is at most
$4\sum_{\mathbf v}D_{\mathbf v}^2\le4Rnd^2$.
Summing the potential identity gives
\[
 H_T=\Phi_0-\Phi_T+\sum_{t=1}^Tz_t+M_T.
\]
Since $z_t\ge0$, the first claimed event implies $M_T<-a$.
Lemma~\ref{lem:martingale_tail} proves the first estimate.

\par\medskip\noindent\textbf{The restrictions introduced at branching steps.} 
At a branching step let $k_t=|g_{\mathbf x_t}(t-1)|/2+1$.
The increment $X_t$ of $A_t$ is $k_t$ in the nonpreferred child and zero
in the preferred child. Each has conditional probability one half, so
$\E[X_t\mid\mathcal H_{t-1}]=k_t/2$.
At a forced step put $X_t=0$ and its conditional mean equal to zero.
The centered sum of these increments is a martingale. At a branching step,
its absolute increment is $k_t/2$, and
\[
 (k_t/2)^2\le(B_{\mathbf x_t}+1)^2/4
       \le(B_{\mathbf x_t}^2+1)/2
       \le B_{\mathbf x_t}^2+1.
\]
The sum of these squared bounds is at most
$\sum_{\mathbf v}(B_{\mathbf v}^2+1)\le Rnd^2+n\le2Rnd^2$.
The sum of the conditional means is at least $G_T/4$, because $G_T$
counts only the branching steps whose variables belong to $S$.
Thus $A_T<G_T/4-b$ implies a centered deviation below $-b$.
Lemma~\ref{lem:martingale_tail} gives the second estimate.

\par\medskip\noindent\textbf{The effect of subsequent assignments.} 
Immediately before assigning $\mathbf x_t$, let $p$ and $q$ be the
numbers of positive and negative occurrences of this variable in previously
created constraints. Its contribution to $Q_t-Q_{t-1}$ is $(p-q)/2$ when
its value is one, and $(q-p)/2$ when its value is zero. These outcomes have
equal conditional probabilities, so the conditional mean is zero.
A newly created constraint contains only variables still unassigned and
makes no additional contribution to $Q_t$ at this step. Hence $Q$ is a
martingale starting at zero. Its increment is at most
$(p+q)/2\le D_{\mathbf x_t}/2$ in absolute value, by
Lemma~\ref{lem:bias_sum}. The squared bounds sum to at most $Rnd^2/4$.
Applying Lemma~\ref{lem:martingale_tail} proves the third estimate.
\end{proof}

\subsection{Dense gaps and the scope of the analysis}

When the initial gaps are linear in $n$, the same analysis gives a simpler
input condition. Simplicity already bounds each variable's total number
of occurrences by $O(n)$, so no additional assumption on the degree
sequence is needed. The following corollary states this condition directly
for the original formula; its proof explains why preprocessing preserves
the quantities needed by the argument.

\begin{corollary}[Linear effective gaps in a simple formula]\label{cor:lineargap}
Fix $\varepsilon,\eta\in(0,1]$. Let $F$ be a simple unweighted Max 2-SAT
formula on $n$ variables. Suppose at least $\eta n$ variables satisfy
\[
 |d_2^F(v)-d_2^F(\neg v)+2(d_1^F(v)-d_1^F(\neg v))|
 \ge\varepsilon n.
\]
Then an optimal assignment for $F$ can be computed deterministically in
polynomial space and time
$\Ostar(2^{(1-\varepsilon^2\eta^2/48000)n})$.
\end{corollary}
\begin{proof}
Before preprocessing, each variable has at most $4(n-1)$ binary occurrences
and at most two unit occurrences. The replacement in \textnormal{\textsc{Preprocess}}
preserves its effective gap: at the shared literal, two signed binary
occurrences are replaced by a unit with coefficient two; at the other
variable, opposite signed occurrences cancel. Unit cancellation also
preserves the gap. Neither operation increases the sum of binary and unit
occurrences at any variable. Thus, in $F_0$,
\[
 B_{\mathbf v}=D_{\mathbf v}+u_{\mathbf v}\le4(n-1)+2\le4n,
 \qquad \sum_{\mathbf v}B_{\mathbf v}^2\le16n^3.
\]
For the set $S$ of biased variables, the initial total absolute gap is at
least $\varepsilon\eta n^2$.

The preceding proof uses the relative-bias hypotheses only through the
two inequalities $\Phi_0\ge\theta nd$ and
$\sum_{\mathbf v}B_{\mathbf v}^2\le Rnd^2$, with $d\ge1$.
Use that proof with scale $d=n$, $\theta=\varepsilon\eta$, and $R=16$.
Here $d$ is a scale, rather than an average residual degree.
All estimates, including those for forced assignments, then apply and give
the stated exponent. In particular, the argument remains valid if
preprocessing converts most of the original binary bias into unit clauses.
\end{proof}

Without units, the hypothesis is exactly a gap between literal frequencies.
For a simple input with units, the gap counting units and binary clauses
equally differs from the effective gap by at most one. A linear gap in
that convention therefore gives a linear effective gap with any fixed
slightly smaller coefficient for sufficiently large $n$.
No restriction on the degree sequence is needed for this dense corollary.

\paragraph{Variable selection and multiplicities.}
The algorithm uses a fixed branching order to make its specification
unambiguous. The probability estimates also apply to a deterministic
selection rule based on the current state: the selected variable is known
before its fair value is revealed, and each variable is still assigned
once. This observation is a separate extension of the stated algorithm.
Similarly, the analysis after preprocessing uses the occurrence counts and
their squared-sum bound, rather than simplicity itself. It therefore
extends to explicitly repeated clauses whenever those quantitative bounds
hold. Simplicity supplies the automatic bound used in the dense corollary.

\paragraph{Order of the two children.}
The preferred value specifies which child receives the pruning constraint.
The correctness proof uses a preferred-first order only to select an
optimum that survives. Actual exploration can visit the two children in
either order, provided the state-based choices and pruning rules remain
the same. No conclusion here relies on a random choice by the algorithm.


\subsection{Interpretation and possible applications}\label{sec:applications}

Literal bias has a direct interpretation in a model of binary decisions:
one alternative participates in more potential conflicts than the other.
This gives a natural way to look for applications of the theorem. It does
not establish that a particular application produces the required density
or gap sizes. We distinguish exact encodings and class membership from
candidate sources of practical instances.

\paragraph{Two alternatives and asymmetric conflicts.}
Suppose object $i$ has two alternatives, represented by $x_i=0$ and $x_i=1$.
These could be two placements, routes, or configurations. A conflict between
the choices $x_i=a$ and $x_j=b$ is represented by the soft clause
\[
 (x_i\ne a)\vee(x_j\ne b),
\]
where $x_i\ne0$ denotes the literal $x_i$ and $x_i\ne1$ denotes $\neg x_i$.
The clause is violated exactly when the conflicting choices are both made.
Thus minimizing the number of pairwise conflicts is precisely unweighted
Max 2-SAT when every conflict has unit cost. With at most one copy of each
conflicting pair of alternatives, this produces a simple formula.

Let $c_i^a$ count the conflicts involving alternative $a$ of object $i$.
Before pair reduction there are no units, and
\[
 d_2(\neg x_i)-d_2(x_i)=c_i^1-c_i^0.
\]
Consequently, if at least $\eta n$ objects satisfy
$|c_i^1-c_i^0|\ge\varepsilon n$, Corollary~\ref{cor:lineargap} applies.
One possible source is a placement problem where one candidate position
lies in a congested region and the other lies in a less congested region.
This is a modeling example, not an assertion about measured placement
instances. The more conflict-prone alternative need not be globally
inferior: its actual cost depends on the choices made for the other objects.
The pruning constraint describes the completions on which that alternative
can still be useful. For smaller relative gaps, the degree-spread condition
in~\eqref{eq:biasregularity} must also be checked.

\paragraph{Same-side and different-side preferences.}
A related model asks for a partition into two labeled sides, with soft
preferences that selected pairs lie on the same or on different sides.
This is the two-cluster version of correlation clustering. Signed-network
models and correlation clustering also arise in the study of social
balance and link classification~\cite{CBGV12}. The exact unit-cost encodings
are
\[
\begin{array}{c|c|c}
\text{preference}&\text{clauses}&\text{number satisfied}\\ \hline
 x_u\ne x_v &(x_u\vee x_v),\ (\neg x_u\vee\neg x_v)
     &1+[x_u\ne x_v]\\
 x_u=x_v &(x_u\vee\neg x_v),\ (\neg x_u\vee x_v)
     &1+[x_u=x_v].
\end{array}
\]
Both encodings contribute one occurrence of each polarity at each endpoint.
Thus adding same-side preferences to Max-Cut does not create literal bias:
every variable remains balanced, regardless of the mixture of preferences.
Complementing every variable preserves the score. These instances may have
other useful structure, but the gap theorem does not follow from the
presence of agreement preferences alone.

Fixed labels or preferences for a particular side can break this symmetry.
For example, substitute the known labels of a set of anchor vertices.
Each preference between an anchor and an unassigned vertex becomes a unit
preference plus a constant. If $h_v$ is the number of resulting units
favoring $x_v=1$ minus the number favoring $x_v=0$, then the effective gap
of the residual formula is $g_{\mathbf v}=2h_v$: interactions between unassigned
vertices still contribute balanced complementary pairs. A few unweighted
anchors cannot generate linear gaps; a linear net preference requires
linearly many such contributions.

There is a precise sufficient condition for this anchored model. Let $N$
vertices remain unassigned, let there be at most $aN$ fixed anchors for a
constant $a$, and allow at most one unit-cost preference per vertex pair.
If at least $\eta N$ unassigned vertices satisfy $|h_v|\ge\beta N$ for
constants $\eta,\beta>0$, the same bias argument gives $\Ostar(c^N)$ time
for some $c<2$. Indeed, $D_{\mathbf v}\le2(N-1)$, $u_{\mathbf v}\le aN$, and hence
\[
 \sum_v(D_{\mathbf v}+u_{\mathbf v})^2\le(2+a)^2N^3,
 \qquad
 \sum_{v\in S}|g_{\mathbf v}|\ge2\beta\eta N^2.
\]
Apply the proof of Corollary~\ref{cor:lineargap} with scale $d=N$,
$R=(2+a)^2$, and $\theta=\min\{1,2\beta\}\eta$.
The original labels are supplied data, rather than values chosen by the
algorithm. Repeated units created by substituting them are retained.
This is a concrete family covered by the proof; whether a given labeled
network has this much net anchor information is an empirical question.

\paragraph{Selecting compatible items.}
Another exact encoding starts from a simple conflict graph $G=(V,E)$.
Introduce a unit clause $(x_v)$ for each item and a binary clause
$(\neg x_u\vee\neg x_v)$ for each conflict $uv\in E$. If $S$ is the set
of selected items, the score is
\[
 F_G(S)=|E|+|S|-|E(G[S])|.
\]
If $S$ contains an edge, removing either endpoint loses one unit reward but
eliminates at least one violated conflict clause. The score therefore does
not decrease. Repeating this operation produces an independent set with
at least the original score. Conversely, an independent set $S$ has score
$|E|+|S|$. It follows that
\[
 \OPT(F_G)=|E|+\alpha(G),
\]
where $\alpha(G)$ is the maximum independent-set size. An optimal assignment
can be converted to an optimal compatible selection by the same procedure.
All clauses have unit weight; there are no additional hard constraints.

Here $d_2(\neg x_v)=\deg_G(v)$, $d_2(x_v)=0$, and there is one positive
unit, so $g_{\mathbf v}=2-\deg_G(v)$. If a linear number of vertices have degree at
least $\kappa n$ for a constant $\kappa>0$, their absolute effective gaps
are at least $\kappa n/2$ for sufficiently large $n$. Thus the encoding
belongs to Corollary~\ref{cor:lineargap}. This example illustrates the
meaning of the class; it is not offered as a new improvement for maximum
independent set.

\paragraph{Documented applications and the remaining empirical question.}
Miyashiro and Matsui~\cite{MM03} formulate sports scheduling with minimum
home/away breaks as Max 2-SAT. One eliminates a variable from each game to
enforce that exactly one team plays at home; the remaining objective is a
collection of soft binary clauses. This establishes an application of Max
2-SAT, but the formulation does not by itself establish our dense-gap
assumptions.

Erlebach, Hall and Schank~\cite{EHS02} use Max 2-SAT when inferring Internet
customer--provider relationships from routing paths. Variables describe
edge orientations, and two-literal clauses forbid incompatible orientations
along consecutive edges of an observed path. For length-two paths there
is one clause per path, so maximizing satisfied clauses exactly maximizes
the number of valid paths. For longer paths this correspondence is not
generally one-to-one; their work analyzes the resulting approximation.
The asymmetric interpretation of the variables makes these formulas a
candidate for examining literal bias, not an established member of our
class.

A useful empirical study would measure the gaps and degree distribution
before and after canonical reduction in such encodings. For the dense
corollary, the relevant evidence is a linear number of linear effective
gaps in a simple input. For the relative-bias theorem, one must also measure
the sum of squared degrees and unit multiplicities in
\eqref{eq:biasregularity}. Repeated observations and unequal conflict costs
must be accounted for explicitly; they cannot be discarded merely to make
an instance simple. Nor can general application-specific hard constraints
be treated as the auxiliary pruning constraints of our algorithm, since
replacement assignments must respect genuine hard constraints. These checks
would establish whether the proved classes explain practical instances,
beyond the natural interpretations given here.

\section{Random formulas}\label{sec:random}

\subsection{Intuition}\label{sec:random_intuition}

The literal-gap result in Section~\ref{sec:bias} uses a substantial initial
preference for one value of many variables. Random formulas need not have
such a preference. For intuition, consider a fixed graph in which every
vertex has degree $d$, and put one clause on each edge by choosing its two
literal signs independently and uniformly. Assume for this explanation
that there are no initial unit clauses. The difference between the two
literal frequencies of a variable is typically of order $\sqrt d$, rather
than a constant fraction of $d$. Thus, for growing $d$, we cannot directly
apply the relative literal-bias theorem. Nevertheless, these smaller gaps
still provide useful pruning constraints.

To see why, follow a path in the branching tree and consider the constraints
created along it. As in the previous section, the preferred value of the
next variable is the one that satisfies more clauses on average when the
remaining variables are assigned uniformly. When the path takes the other
value, we add a constraint requiring that this choice be strictly better
than the preferred value for the eventual completion. This constraint
requires the remaining assignment to reverse the current preference.
In the random model, a typical preference has size of order $\sqrt d$.
The resulting individual constraint is therefore weaker than one arising
from a gap proportional to $d$.

It is useful to relate this to ordinary SAT clauses. A clause on $k$
distinct variables forbids exactly one of their $2^k$ assignments, so a
shorter clause excludes a larger fraction of the possibilities. Our
pruning constraints can forbid many assignments, and their strength
depends on both their length and their bound. For illustration, consider
\[
 z_1+\cdots+z_k\le \frac{k}{2}-r,
\]
where the $z_i$ are distinct Boolean variables and $r>0$. Under a uniform
assignment, the sum has mean $k/2$ and standard deviation $\sqrt{k}/2$.
Satisfying the constraint therefore requires a downward deviation of at
least $r$ from the mean. We call $r$ its \emph{deficit}.

For a fixed deficit, reducing the number of variables makes the required
deviation larger relative to the standard deviation. This explains why
shorter constraints can prune more. If instead the deficit grows
proportionally to $\sqrt{k}$, the required number of standard deviations
stays constant, and constraints of different lengths have comparable
pruning power. For example,
\[
 \sum_{i=1}^{100}z_i\le40,
 \qquad
 \sum_{i=1}^{400}z_i\le180
\]
both require a deviation of two standard deviations below the mean. The
second constraint is four times longer, but the fractions of assignments
satisfying them are comparable. If the deficit is a fixed positive
fraction of $k$, however, only an exponentially small fraction of
assignments satisfy the constraint. In that case, a longer constraint
can be substantially stronger. Thus length alone does not determine
pruning power: the deficit must be compared with the fluctuations that
the remaining variables can provide.

Now consider a path obtained by choosing each branching value with
probability $1/2$. About half of its branching steps take the nonpreferred
value, so a linear number of branching steps creates a linear number of
constraints. Their requirements accumulate. This alone would not suffice:
in an adversarial formula, many constraints might ask for essentially the
same assignment to their shared variables. One choice of a shared variable
could then help satisfy all of them. Random literal signs prevent this
alignment from persisting throughout the accumulated constraints.

The role of opposing signs can be seen in a small example. Suppose two
constraints contain an unassigned variable $y$ as follows:
\[
 y+A\le b,\qquad \neg y+B\le c,
\]
where $A$ and $B$ are sums of other literal occurrences. Adding the
constraints and using $y+\neg y=1$ gives
\[
 A+B\le b+c-1.
\]
Whichever value we give $y$, it helps one constraint and hurts the other.
Its contribution to their sum is fixed. If both constraints instead
contained $y$, setting $y=0$ would help both. A single cancellation need
not yield substantial pruning, but many such cancellations limit how much
the remaining assignments can help satisfy the combined requirement.

This is the useful effect of overlap in random formulas. Shared variables
occur with both polarities, so positive and negative occurrences largely
cancel when we add the constraints. The relevant diversity is therefore
in the literal signs, not necessarily in the sets of neighboring variables:
even constraints with the same variable set can impose different demands.
Cancellation preserves the accumulated shortfall from the average while
reducing the variation that the remaining variables can provide. The
summed constraint can consequently admit very few completions even before
it becomes infeasible. The algorithm exploits this by enumerating the
completions of that sum, as described below.

For the regular-graph example, the requirements accumulated over a linear
number of branching steps have scale $n\sqrt d$, whereas the ordinary
fluctuations in the summed contributions have scale $\sqrt{nd}$. The proof
makes this comparison precise and bounds the exceptional assignments that
provide enough compensation by an exponentially small fraction. It does
not assume that the learned constraints are independent, nor that every
path has the stated behavior. Instead, it bounds the total enumeration
work over all paths and proves an exponential saving for most random
inputs. The algorithm itself is deterministic.

For an irregular graph, the degree-dependent weights introduced below
put these requirements on a comparable scale before adding them. A larger
constraint does not necessarily prune a larger fraction of assignments:
in the random setting, its larger typical imbalance is accompanied by
larger fluctuations. Dividing by approximately the square root of the
degree accounts for this change of scale. Multiplying a constraint by a
positive weight does not change its satisfying assignments; the weight
determines its contribution to the summed constraint. An unusually large
imbalance still gives a stronger requirement after normalization. These
weights allow the argument to handle arbitrary degree sequences and
several independently sampled clauses on a variable pair.

\subsection{Notation and enumeration}\label{sec:enumeration}

We use the same residual formulas, literal counts, and pruning constraints
as in Section~\ref{sec:bias}. The new operation is to stop branching after
a specified number of variables and enumerate completions satisfying the
summed constraint $C_\Sigma$ defined in~\eqref{eq:weighted} below.

Let $n$ denote the number of input variables. The algorithm separately
optimizes variables occurring in no binary clause and searches on the
remaining $r\le n$ variables. List these as
$\mathbf u_1,\ldots,\mathbf u_r$ in their original input order, inherited
from the fixed input indices. This list
specifies the branching order; the pruning constraints do not determine
or change it. Write $x_t$ for the Boolean value of $\mathbf u_t$.
A \emph{prefix assignment of length $T$} specifies
$x_1,\ldots,x_T$; a \emph{suffix assignment} specifies the remaining
$r-T$ values. These terms refer to this fixed order.
As before, $\alpha$ denotes the current partial assignment and $\mathcal C$
the collection of pruning constraints. We retain the individual constraints,
substituting each assigned value in them as the search proceeds.

To define $C_\Sigma$, give each constraint $C:\sum_i\ell_i\le b_C$ a
positive rational weight $\omega_C$, fixed when it is created. Multiply
both sides of each constraint by its weight and add the resulting
inequalities:
\begin{equation}
 C_\Sigma:\qquad
 \sum_{C\in\mathcal C}\omega_C\sum_i\ell_i
 \le \sum_{C\in\mathcal C}\omega_C b_C .
\label{eq:weighted}
\end{equation}
These weights are assigned to whole constraints; they are distinct from
the coefficients obtained by collecting repeated literal occurrences.
When every $\omega_C=1$, this is the ordinary sum used in
Lemma~\ref{lem:constraint_sum}, and each literal's coefficient is simply
its total number of occurrences. The random-formula algorithm specifies
the weights in~\eqref{eq:random_weights}.
Every assignment satisfying all individual constraints satisfies this sum.
The converse need not hold. In an empty collection the sum is $0\le0$.

The sum can be simplified just as in Lemma~\ref{lem:constraint_sum}.
If a variable contributes $aw+b\neg w$ with $a,b\ge0$, replace this by
$(a-b)w+b$ when $a\ge b$, or by $(b-a)\neg w+a$ otherwise, and move the
constant to the right-hand side. After collecting all occurrences, we obtain
\begin{equation}
 \sum_{i=1}^k w_i z_i\le B,\qquad w_i\ge0,\quad z_i\in\bits,
 \label{eq:budget}
\end{equation}
where each $z_i$ is one remaining variable or its negation, and each
remaining variable appears once. Variables with zero coefficient are
included as well, so an enumerated assignment is complete.

To enumerate this inequality, assign the $z_i$ one at a time and reject
a partial assignment as soon as its assigned contribution exceeds $B$.
Every retained partial assignment has a completion: set all remaining
$z_i$ to zero. Lemma~\ref{lem:enumeration}, proved after the main running-time
argument, uses this observation to bound enumeration time by the number
of satisfying completions, up to polynomial factors.

\subsection{The algorithm}\label{sec:random_algorithm}

The search has two stages. First, it branches on a fixed prefix of the
variable order, using the preferred literal and pruning constraint from
the literal-gap algorithm. For each prefix assignment that survives,
it then enumerates assignments to the remaining variables satisfying
$C_\Sigma$. These completions may violate an individual constraint;
we evaluate their actual scores, so considering them does not affect
correctness. Enumeration under $C_\Sigma$, in the simplified
form~\eqref{eq:budget}, has the running-time guarantee just described.

The algorithm receives only $F$. It constructs the graph
$G_F=(V_F,E_F)$ directly from the binary clauses of $F$, before applying
\textnormal{\textsc{Preprocess}}:
\begin{equation}
 E_F=\bigl\{\{u,v\}:F\text{ contains a binary clause on }u,v\bigr\}.
\label{eq:observed_support}
\end{equation}
The vertices are exactly the endpoints, $V_F=\bigcup_{e\in E_F}e$, so
$G_F$ has no isolated vertices. Literal signs do not affect whether a
pair is an edge, and unit clauses create no edges. This graph remains
fixed throughout the search. The graph used to generate $F$ belongs to
the distributional assumptions below and is not supplied to the algorithm.

A variable outside $V_F$ occurs only in unit clauses. Set it to the value
satisfying more of its units, breaking ties by choosing zero. Let
$\alpha_0$ collect these assignments, let $s_0$ be their total score,
and let $F_{\rm core}$ contain the clauses on the remaining variables.
An optimum for $F_{\rm core}$, extended by $\alpha_0$, is an optimum for
$F$, with $s_0$ added to its score. If $V_F$ is empty, this already solves
the instance. Otherwise write $r=|V_F|$ and search on these $r$ variables.

Before preprocessing, record the weights
\begin{equation}
 \omega_t=\frac{1}{\lceil\sqrt{\deg_{G_F}(\mathbf u_t)}\rceil}.
 \label{eq:random_weights}
\end{equation}
A constraint created when branching on variable $\mathbf u_t$ receives weight
$\omega_t$. As explained in Section~\ref{sec:random_intuition}, this
normalizes the degree-dependent scale of the random requirements and
their fluctuations. The actual imbalance remains in the bound of each
constraint. If $G_F$ is regular, all weights are equal,
and the weighted sum is equivalent to the ordinary sum.
The weights remain fixed throughout the search.

Set $T=\lceil r/2\rceil$. This is the number of variables assigned by
branching before switching to enumeration; it is not the number of
constraints created. At a node where $t<T$ variables have been assigned,
the next branching variable is $\mathbf u_{t+1}$. Once $t=T$, the search
creates no further constraints on that path and enumerates assignments
to $\mathbf u_{T+1},\ldots,\mathbf u_r$ satisfying the current $C_\Sigma$.
For example, with eight remaining variables we branch on $\mathbf u_1,\ldots,\mathbf u_4$,
then complete each surviving prefix by enumerating the last four variables.
The choice of half the variables gives at most $2^T$ prefixes while
allowing a linear number of branching steps to create constraints; the
proof below bounds the total work of completing those prefixes.

The two procedures use the following data. Shared variables are initialized
once by \textnormal{\textsc{RandomFormula}} and are available to every call
of \textnormal{\textsc{RandomSearch}}.
\begin{description}
\item[Fixed shared data.]
$U=(\mathbf u_1,\ldots,\mathbf u_r)$ is the variable order, and
$(\omega_1,\ldots,\omega_r)$ contains the constraint weights.
Neither changes during the search. All variables remain in $U$, including
variables that disappear from the residual formula.
\item[Shared best solution.]
$M$ is the largest score in $F_{\rm core}$ found so far, and
$\alpha^\star$ assigns all variables of $V_F$ and attains that score.
The outer procedure adds $s_0$ and $\alpha_0$ when returning the solution
for the original formula. Initially
$M=-\infty$ and $\alpha^\star=\emptyset$, indicating that no complete
assignment has yet been evaluated. Every recursive call may update this
same pair. An improvement is retained when the call returns and when
the search backtracks to explore another branch.
\item[Parameters of a recursive call.]
$F$ is the residual soft formula, $\mathcal C$ is the collection of
pruning constraints with their weights, and $\alpha$ is the current
partial assignment. The value $s$ is the score offset accumulated through
preprocessing and assignments: for any completion $\beta$, its score in
the formula $F_{\rm core}$ is $s$ plus the number of clauses it satisfies in
$F$. These four parameters belong to the current node and are copied
separately for its children. The parameter $T$ is the stopping depth;
it is passed unchanged to every recursive call.
\end{description}
At each call, $\alpha$ assigns exactly the first $|\alpha|$ variables of
$U$. In particular, $M$ and $\alpha^\star$ are shared across branches,
whereas $s$ and $\alpha$ describe only the current branch.
The identities used by \textnormal{\textsc{Preprocess}} apply to copies
of clauses as well, so the procedure also preserves scores for the
bounded multisets allowed by the random model.

\begin{algorithm}[H]
\caption{\textnormal{\textsc{RandomFormula}}$(F)$}\label{alg:weighted}
\begin{algorithmic}[1]
\STATE Construct $G_F=(V_F,E_F)$ from $F$ by~\eqref{eq:observed_support}, before preprocessing.
\STATE Optimize variables outside $V_F$, obtaining $F_{\rm core},s_0,\alpha_0$ as above.
\IF{$V_F=\emptyset$}
  \STATE \textbf{return} $(s_0,\alpha_0)$.
\ENDIF
\STATE Initialize shared $U=(\mathbf u_1,\ldots,\mathbf u_r)$ in the original input order, where $r=|V_F|$.
\STATE $T\gets\lceil r/2\rceil$.
\COMMENT{Branch on the first $T$ variables, then enumerate completions.}
\STATE Initialize the shared weights $\omega_t$ from $G_F$ by~\eqref{eq:random_weights}.
\STATE Initialize shared $M\gets-\infty$ and $\alpha^\star\gets\emptyset$.
\STATE $(F,s)\gets\textnormal{\textsc{Preprocess}}(F_{\rm core})$.
\STATE $\textnormal{\textsc{RandomSearch}}(F,\emptyset,\emptyset,s,T)$.
\STATE \textbf{return} $(M+s_0,\alpha^\star\cup\alpha_0)$.
\end{algorithmic}
\end{algorithm}

\begin{algorithm}[H]
\caption{\textnormal{\textsc{RandomSearch}}$(F,\mathcal C,\alpha,s,T)$}
\label{alg:random_search}
\begin{algorithmic}[1]
\REQUIRE Fixed shared data $U,(\omega_t)_{t=1}^r$ and shared best solution $M,\alpha^\star$, initialized above.
\STATE $t\gets|\alpha|$.
\COMMENT{The variables $\mathbf u_1,\ldots,\mathbf u_t$ have been assigned.}
\STATE Form $C_\Sigma$ by~\eqref{eq:weighted} and simplify it to~\eqref{eq:budget}.
\IF{its bound $B$ is negative}
  \STATE \textbf{return}
\ENDIF
\IF{$t=T$}
  \FOR{each suffix assignment $\beta$ satisfying $C_\Sigma$, by Lemma~\ref{lem:enumeration}}
    \STATE Let $q$ be $s$ plus the number of clauses of $F$ satisfied by $\beta$.
    \STATE If $q>M$, update shared $M\gets q$ and $\alpha^\star\gets\alpha\cup\beta$.
  \ENDFOR
  \STATE \textbf{return}
\ENDIF
\STATE Let $\mathbf u\gets\mathbf u_{t+1}$.
\COMMENT{The next variable in the fixed list $U$.}
\STATE Choose its preferred literal $v$ by~\eqref{eq:bias_preference}.
\STATE Construct $C$ from $F$ and $v$ by Lemma~\ref{lem:pruning}; give it weight $\omega_{t+1}$.
\FOR{$p\in[1,0]$}
  \STATE Make separate copies $(F_p,\mathcal C_p,\alpha_p,s_p)$ of the current state.
  \IF{$p=0$}
    \STATE $\mathcal C_p\gets\mathcal C_p\cup\{C\}$.
  \ENDIF
  \STATE Assign $v=p$, updating the copied formula, constraints, assignment, and score.
  \STATE $\textnormal{\textsc{RandomSearch}}(F_p,\mathcal C_p,\alpha_p,s_p,T)$.
\ENDFOR
\end{algorithmic}
\end{algorithm}

This algorithm uses substitution, cancellation, and rejection during the
prefix search, but does not call the full \textnormal{\textsc{Simplify}}
procedure. In particular, it makes no forced assignments before reaching
depth $T$. This keeps the first $T$ variables fixed, as required by the
analysis below. Created constraints remain in $\mathcal C$ with their
updated bounds, even if they have become automatically satisfied.

\subsection{Correctness}

The comparisons are the same as in the literal-gap search. The only new
point is that suffix enumeration enforces $C_\Sigma$ from~\eqref{eq:weighted}, which is
weaker than enforcing every constraint separately. The next lemma shows
that this still retains an optimum; its proof does not use randomness.

\begin{lemma}[Correctness of weighted enumeration]\label{lem:exactA}
For every formula $F$, including formulas with repeated
clauses, \textnormal{\textsc{RandomFormula}} returns an optimum score and
an assignment attaining it.
\end{lemma}
\begin{proof}
Variables outside $V_F$ can be optimized independently from their units.
It therefore suffices to prove that the recursive search finds an optimum
for $F_{\rm core}$. Preprocessing changes the score of every assignment
by the same recorded constant. It therefore preserves both optimal assignments and all
comparisons between the two values of a variable.

Consider the full decision tree in the fixed variable order, restoring
any prefixes rejected by the algorithm. At each of the first $T$ levels,
place the preferred child first. Complete the rest of the tree in any
fixed order, and choose the earliest optimal complete assignment $x^*$.
If its path takes a nonpreferred child at variable $t$, changing only
that value to the preferred one produces an earlier assignment. Its
score must be strictly smaller: a larger score contradicts optimality,
and an equal score contradicts the choice of $x^*$.
Lemma~\ref{lem:pruning} therefore implies that $x^*$ satisfies the
constraint added at that step.

Thus $x^*$ satisfies every constraint created on its prefix, and hence
their positively weighted sum at every earlier depth. The sum test never
rejects that prefix, and enumeration at depth $T$ visits its suffix.
Every enumerated candidate is scored in $F_{\rm core}$, including
the preprocessing and assignment offsets. Adding the separately optimized
score $s_0$ and extending by $\alpha_0$ therefore returns an optimum for $F$.
\end{proof}

\subsection{Random inputs and the running-time guarantee}

Two familiar models motivate the assumptions below. Fix any set $E$
of pairs of distinct input variables, of any cardinality. There is no
density, regularity, or connectivity requirement. The generating graph
is $G=(V,E)$, where $V=\bigcup_{e\in E}e$: an edge means that its pair
is selected for clause sampling. The graph is fixed before sampling and
is not an input to the optimization algorithm. On each
selected pair we may choose exactly one of the four possible clauses,
uniformly at random and independently across pairs.

Alternatively, include each of the four possible clauses independently
with probability $1/4$, with all choices independent across pairs as well.
The number $X_e$ of clauses on a selected pair $e\in E$ then satisfies
\[
 X_e\sim\operatorname{Binomial}(4,1/4),\qquad \E X_e=1.
\]
Thus a selected pair can contribute zero, one, two, three, or four clauses.
There is one clause per selected pair in expectation, and $|E|$ binary
clauses in total in expectation. Pairs outside $E$ contribute no binary
clauses. More generally, any fixed inclusion probability $q\in(0,1)$ is
allowed: then $X_e\sim\operatorname{Binomial}(4,q)$ and $\E X_e=4q$.
The observed graph $G_F$ contains only pairs that actually receive a
clause. Thus $E_F\subseteq E$, with equality in the one-clause-per-pair
model but not necessarily under independent inclusion.

Both models treat the two polarities of each endpoint symmetrically.
The next definition allows both models, while excluding pair distributions
whose endpoint literal counts are always exactly balanced.

\begin{definition}[Complementation-invariant random signs]\label{def:invariant}
Fix constants $m_0\ge1$ and $0<\kappa\le m_0^2$, a set of $n$ input variables,
and any set $E$ of unordered pairs of distinct input variables.
Let $V=\bigcup_{e\in E}e$ and $G=(V,E)$.
For each edge $e=\{u,v\}$ draw a multiset $M_e$ of binary clauses on
$u,v$, independently over the edges. The binary part of $F$ is the
multiset union of these $M_e$; pairs outside $E$ contribute no binary
clauses. A sampled $M_e$ may be empty, so the observed graph $G_F$
in~\eqref{eq:observed_support} may have fewer edges than $G$. Assume that
\begin{enumerate}
\item $|M_e|\le m_0$;
\item complementing the literal of either endpoint in every clause of
$M_e$ leaves its distribution unchanged;
\item for each endpoint $u$,
\[
 \E\left[\left(\sum_{c\in M_e}s_u(c)\right)^2\right]\ge\kappa,
\]
where $s_u(c)=1$ if the literal of $u$ in $c$ is positive and
$s_u(c)=-1$ otherwise.
\end{enumerate}
Add any multiset of unit clauses on the input variables, fixed
independently of these draws.
\end{definition}

The last condition requires some variation in the difference between
positive and negative occurrences at an endpoint. One uniformly signed
clause per edge has $m_0=\kappa=1$. Including each of the four clauses
independently with probability $q\in(0,1)$ has $m_0=4$ and
$\kappa=4q(1-q)$. In particular, $q=1/4$ gives the independent-clause
model discussed in the introduction.

The theorem bounds the complexity of solving these random inputs.
Its constants depend on the pair distribution but not on the density,
degree sequence, or choice of the fixed support graph.

\begin{theorem}[Random signs on every graph]\label{thm:weighted}
Fix $m_0\ge1$ and $0<\kappa\le m_0^2$. For a formula $F$ on $n$ variables drawn
according to Definition~\ref{def:invariant} with any fixed generating
graph $G$, an optimal assignment can be computed from $F$ alone by a
deterministic polynomial-space algorithm with expected running time
\[
 \Ostar(2^{n/2}+2^n e^{-c_0n}),
 \qquad
 c_0=\frac{c_1^2}{10816m_0^2},\qquad
 c_1=\frac{\sqrt\kappa}{2\sqrt{1+4m_0^4/\kappa^2}}.
\]
With probability at least $1-3e^{-c_0n/2}$, its running time is
$\Ostar(2^ne^{-c_0n/2})$.
The algorithm returns an optimum on every input; the expectations and
probabilities refer only to the random input formula.
\end{theorem}

For one uniformly signed clause per edge, $c_0=1/216320$.
Appendix~\ref{app:regular} gives a sharper constant on almost-regular
graphs without units. The theorem above also covers irregular supports
and independently fixed unit clauses. The bound holds for any variable
order fixed independently of the sampled signs.

\subsection{Proof of the running-time bound}

We first condition on the observed graph $G_F$. Lemma~\ref{lem:observed_pairs}
shows that the surviving pairs still satisfy the random-model assumptions,
with the same constants. Thus the graph used by the algorithm can be
treated as fixed for the conditional analysis, even though it was
constructed from $F$. Write $r=|V_F|$ for the number of variables searched.
We count assignments to these variables that satisfy the summed constraint
at their own prefix. The enumeration lemma then turns that count into a
time bound. The proof has three ingredients.
\begin{enumerate}
\item Along a uniform path, each nonpreferred choice adds a positive
requirement measured relative to the average contribution of the
remaining literals. Their weighted total is at least of order $r$ on
most paths. Lemma~\ref{lem:random_requirements} bounds the probability
that it is much smaller.
\item The actual remaining values can help satisfy these requirements,
but their effects have random signs. Lemma~\ref{lem:random_compensation}
bounds the chance of an unusually favorable total contribution.
Its proof uses the randomness of unrevealed endpoint signs, not
independence of the learned constraints.
\item Combining these estimates shows that an exponentially small
fraction of conditional formula--assignment pairs pass the sum. Averaging
over assignments and then over the observed graph bounds the expected
enumeration work; a final probability bound gives the running time for
most formulas.
\end{enumerate}
The supporting lemmas and their elementary probability estimates follow
the main proof.

\begin{proof}[Proof of Theorem~\ref{thm:weighted}]
Use \textnormal{\textsc{RandomFormula}} on $F$ alone.
Correctness is Lemma~\ref{lem:exactA}. Depth-first
search and suffix enumeration use polynomial space. Constructing $G_F$
and optimizing variables outside $V_F$ take polynomial time.

Condition on $G_F=H$, where $H$ is any possible nonempty observed graph,
and write $r=|V(H)|$. By Lemma~\ref{lem:observed_pairs}, its pair laws
satisfy the assumptions of Definition~\ref{def:invariant} with the same
$m_0,\kappa$. Its vertex order, inherited from the input indices, and its
degrees are now fixed. All probabilities in Steps~1 and~2 refer to this
conditional model. The weights are rational with polynomial encoding
length, and satisfy
\begin{equation}
 \frac{1}{2\sqrt{\deg_H(\mathbf u_t)}}\le\omega_t
 \le\frac{1}{\sqrt{\deg_H(\mathbf u_t)}}.
 \label{eq:random_weight_bounds}
\end{equation}
Thus all arithmetic per search state has polynomial cost.

\par\medskip\noindent\textbf{1. Express the summed requirement.}
Choose a uniform assignment $x\in\bits^r$, independently of the remaining
random choices in $F$. For the analysis only, construct its first $T=\lceil r/2\rceil$
comparisons even if an earlier sum test would have rejected its prefix.
Let $\mathrm{gain}_t(x)$ be the score with $x_t=1$ minus the score
with $x_t=0$, keeping all other values fixed. It does not depend on $x_t$.
Let $\Delta_t$ be its average over uniform later values, with
$x_1,\ldots,x_{t-1}$ fixed. The preferred value of the original variable
is $p_t=1$ when $\Delta_t\ge0$ and $p_t=0$ otherwise.
Equivalently, the algorithm's preferred literal is $x_t$ or $\neg x_t$,
respectively, by~\eqref{eq:bias_preference}.

Put $I=\{t\le T:x_t\ne p_t\}$, the indices at which a constraint is
created. Lemma~\ref{lem:pruning} expresses each such constraint as
$(2p_t-1)\mathrm{gain}_t(x)+1\le0$. Hence its weighted sum is
\begin{equation}
 A_T+Q_T\le0,
 \label{eq:random_sum}
\end{equation}
where
\begin{align*}
 A_T&=\sum_{t\in I}\omega_t(|\Delta_t|+1),\\
 Q_T&=\sum_{t\in I}\omega_t(2p_t-1)
                       \bigl(\mathrm{gain}_t(x)-\Delta_t\bigr).
\end{align*}
Here $A_T$ records the requirements at creation, while $Q_T$ records the
deviation supplied by the subsequent values. The identity follows from
$(2p_t-1)\Delta_t=|\Delta_t|$.
Preprocessing preserves each score difference, so these expressions
computed from the original clauses are exactly the ones used by the search.

\par\medskip\noindent\textbf{2. Bound the fraction that satisfies the sum.}
Lemma~\ref{lem:random_requirements} gives, for any $a,b>0$,
\[
 \Prb\left(A_T<\frac{c_1T}{26}-a-b\right)
 \le e^{-2a^2/(Tm_0^2)}+e^{-169b^2/(m_0^2r)}.
\]
Lemma~\ref{lem:random_compensation} gives, for any $c>0$,
\[
 \Prb(Q_T<-c)\le e^{-2c^2/(m_0^2T)}.
\]
Take $a=b=c=c_1T/104$. Outside these three exceptional events,
\[
 A_T+Q_T\ge\frac{c_1T}{26}
              -3\frac{c_1T}{104}
          =\frac{c_1T}{104}>0.
\]
Thus the summed constraint fails. The first and third exponents equal
$c_1^2T/(5408m_0^2)\ge c_0r$.
The second is
$169c_1^2T^2/(10816m_0^2r)\ge c_0r$, since $T\ge r/2$.
Consequently
\[
 \Prb_{F,x}(A_T+Q_T\le0\mid G_F=H)\le3e^{-c_0r}.
\]

\par\medskip\noindent\textbf{3. Convert the count to running time.}
For a prefix $\alpha$, let $N(\alpha)$ count suffix assignments satisfying
its sum, and put $W(F)=\sum_{\alpha\in\bits^T}N(\alpha)$.
This definition includes prefixes rejected early, and so bounds the
number of completions actually enumerated. Since $x$ is uniform,
\[
 \E[W(F)\mid G_F=H]
 =2^r\Prb_{F,x}(A_T+Q_T\le0\mid G_F=H)
 \le3\cdot2^r e^{-c_0r}.
\]
If $H$ is empty, the algorithm finishes in polynomial time; define
$W(F)=1$ in this case, so the same bound holds with $r=0$.
Since $\kappa\le m_0^2$, we have $c_0\le1/216320<\ln2$.
The quantity $2^r e^{-c_0r}$ is therefore increasing in $r$.
Using $r\le n$ and averaging over $G_F$ gives
\[
 \E W(F)\le3\cdot2^n e^{-c_0n}.
\]
Lemma~\ref{lem:enumeration} bounds the search work by
$\Ostar(2^T+W(F))$. The deterministic inequality
$2^T\le2^{\lceil n/2\rceil}$ now proves the expected-time assertion.
Markov's inequality applied to the unconditional expectation also gives
\[
 \Prb_F\left(W(F)>2^n e^{-c_0n/2}\right)
 \le3e^{-c_0n/2}.
\]
The prefix cost $2^T$ is smaller than the claimed high-probability time
bound, up to a constant factor, because $c_0<\ln2$.
This completes the proof.
\end{proof}

\subsection{Supporting lemmas}

We first justify the enumeration cost used in Step~3. A small number of
satisfying completions does not by itself bound a search over arbitrary
constraints. For the single inequality~\eqref{eq:budget}, however, every
retained partial assignment has a completion obtained by setting its
remaining oriented variables to zero.

\begin{lemma}[Enumerating the summed constraint]\label{lem:enumeration}
Suppose an inequality $\sum_{i=1}^k w_i z_i\le B$ with nonnegative rational
weights has $N$ satisfying Boolean assignments. Depth-first enumeration
that rejects a partial assignment as soon as its assigned contribution
exceeds $B$ lists all these assignments in polynomial space and
$\operatorname{poly}(n)(N+1)$ time, including polynomial arithmetic costs.
Consequently, if learning is stopped after the first $T$ variables in a
fixed order, and $N(\alpha)$ counts the suffix assignments satisfying the
sum at prefix assignment $\alpha$, the total enumeration time is
\begin{equation}
 \Ostar(2^T+W(F)),\qquad
 W(F)=\sum_{\alpha\in\bits^T}N(\alpha).
 \label{eq:enumerationcost}
\end{equation}
\end{lemma}
\begin{proof}
If $B<0$, even the minimum left-hand side, zero, exceeds $B$, so the root
is rejected. Otherwise, every retained partial assignment extends to a
solution by setting all its unassigned $z_i$ to zero. At a fixed depth,
distinct retained partial assignments have disjoint sets of completions.
There are therefore at most $N$ retained nodes at that depth. Each has at
most two children, so rejected children increase the count by at most a
constant factor. Summing over the $k+1$ depths gives the first bound.
The depth-first stack has polynomial size.

There are $2^T$ prefix assignments. Constructing their sums costs polynomial
time per prefix; applying the first bound to each and adding the costs
gives the second bound. Prefixes rejected earlier only reduce the work.
\end{proof}

The algorithm does not know which pairs were available for sampling. The
next lemma justifies analyzing it after fixing the graph it actually
observes. Crucially, whether a pair is empty does not distinguish between
the literal polarities.

\begin{lemma}[Conditioning on the observed pairs]\label{lem:observed_pairs}
Suppose $F$ is generated according to Definition~\ref{def:invariant}.
Conditional on any observed graph $G_F=H$ of positive probability, the
clause multisets on its edges are independent. Each has its original
distribution conditioned on being nonempty, and satisfies all three
conditions of that definition with the same $m_0$ and $\kappa$.
\end{lemma}
\begin{proof}
The event $G_F=H$ requires $M_e\ne\emptyset$ for each edge of $H$ and
$M_e=\emptyset$ for every other edge of the generating graph. These events
concern separate independent draws. Conditioning on them therefore
preserves independence, and conditions each retained pair only on its
own nonemptiness. The clause-count bound remains $m_0$.

Complementing either endpoint does not change whether a multiset is
empty. Its conditional distribution is consequently still invariant
under either complementation. Finally, put
$Z_{u,e}=\sum_{c\in M_e}s_u(c)$ and
$p_e=\Prb(M_e\ne\emptyset)$. An empty multiset contributes zero to
$Z_{u,e}^2$, so for every retained edge
\[
 \E[Z_{u,e}^2\mid M_e\ne\emptyset]
 =\frac{\E Z_{u,e}^2}{p_e}\ge\kappa.
\]
The unit clauses were fixed independently of the pair draws and remain
fixed under this conditioning. This proves the claim.
\end{proof}

For the probability estimates below, fix a nonempty observed graph
$G_F=H=(V_H,E_H)$ and put $r=|V_H|$. All probabilities and expectations
refer to the conditional model in Lemma~\ref{lem:observed_pairs}.
We identify the variables with their positions $1,\ldots,r$ in $U$
when writing edges and neighbors.

The next lemma identifies the randomness available when a constraint is
created and when a later variable is assigned. It is the reason that
overlapping constraints can be analyzed without treating them as
independent. We use the original sampled clauses in this analysis, since
preprocessing preserves all gains.

Let $b_t$ be the number of original units $(x_t)$ minus the number of
units $(\neg x_t)$. For an edge $\{t,v\}$, define its contribution to
the gain at $t$ by
\[
 G_{tv}(x_v)=
 \sum_{\substack{c\in M_{\{t,v\}}\\
                 \text{the literal on }v\text{ is false}}}s_t(c).
\]
A clause contributes its sign at $t$ exactly when its partner literal
is false. Thus
\begin{align*}
 \mathrm{gain}_t(x)&=b_t+\sum_{v\sim t}G_{tv}(x_v),\\
 \Delta_t&=b_t+\sum_{\substack{v<t\\v\sim t}}G_{tv}(x_v)
       +\frac12\sum_{\substack{v>t\\v\sim t}}\sum_{c\in M_{\{t,v\}}}s_t(c).
\end{align*}
Write $d_t=\deg_H(\mathbf u_t)$.

\begin{lemma}[Revealing the endpoint signs]\label{lem:random_exposure}
The random formula and the independent uniform assignment can be exposed
in variable order so that, just before the signs at variable $t$ are
revealed,
\[
 \Delta_t=b_t+\sum_{v\sim t}\rho_{t,\{t,v\}}a_{tv},
 \qquad Y_t=\sum_{v\sim t}a_{tv}^2\le m_0^2d_t.
\]
The coefficients $a_{tv}$ are already known, and the signs
$\rho_{t,\{t,v\}}$ are independent fair signs. After they are revealed,
$x_t$ is still a fair bit independent of $\Delta_t$.
For each edge $\{t,v\}$ with $t<v$, the later deviation has the form
\[
 G_{tv}(x_v)-\frac12\sum_{c\in M_{\{t,v\}}}s_t(c)
   =\rho_{t,\{t,v\}}\eta_{tv}\theta_{v,\{t,v\}},
 \qquad |\eta_{tv}|\le m_0/2.
\]
Here $\eta_{tv}$ is known from the start. The signs
$\theta_{v,\{t,v\}}$ are fresh independent fair signs when endpoint $v$
is reached, conditional on all earlier exposed information.
\end{lemma}
\begin{proof}
For each edge $e$, first draw an independent base multiset $K_e$ with
the prescribed law of $M_e$. Write $\sigma_u(c)$ for an endpoint sign
in this base multiset. Independently draw a fair sign $\rho_{u,e}$ for
each endpoint of each edge. Obtain $M_e$ by replacing its endpoint signs
by $\rho_{u,e}\sigma_u(c)$. Complementation invariance implies that the
result has the required distribution, independently on different edges.

Draw the uniform assignment independently and put
$\theta_{u,e}=\rho_{u,e}(2x_u-1)$. Conditional on any assignment $x$,
the transformation from the $\rho$ signs to the $\theta$ signs preserves
their independent fair distribution. Consequently all the $\theta$ signs
are mutually independent and independent of $x$ and of the base multisets.
A partner literal on $v$ is false precisely when
$\sigma_v(c)\theta_{v,e}=-1$.

Reveal all $K_e$ first. At step $t$, reveal the signs $\rho_{t,e}$ on
edges incident to $t$, followed by $x_t$. Let $\mathcal H_t$ be the
information after this step. For $e=\{t,v\}$ set
\[
 a_{tv}=
 \begin{cases}
 \displaystyle\sum_{\substack{c\in K_e\\
                    \sigma_v(c)\theta_{v,e}=-1}}\sigma_t(c),&v<t,\\[6pt]
 \displaystyle\frac12\sum_{c\in K_e}\sigma_t(c),&v>t.
 \end{cases}
\]
These coefficients are determined by $\mathcal H_{t-1}$. Multiplying
them by $\rho_{t,e}$ gives precisely the earlier-pair contributions or
the average later-pair contributions to $\Delta_t$. This proves its
representation. The conditional variance of the signed sum is $Y_t$,
because the new signs are independent and have mean zero. Also
$|a_{tv}|\le m_0$, proving the upper bound on $Y_t$.

For $t<v$ define
\[
 \eta_{tv}=
 \sum_{\substack{c\in K_e\\\sigma_v(c)=-1}}\sigma_t(c)
       -\frac12\sum_{c\in K_e}\sigma_t(c).
\]
This is half the difference of the signed sums over the two partner
polarities, so $|\eta_{tv}|\le |K_e|/2\le m_0/2$.
If $\theta_{v,e}=1$, the first sum selects the false partner literals;
if $\theta_{v,e}=-1$, their complementary set is selected.
Subtracting the average gives the asserted deviation with opposite signs
in these two cases. The $\theta$ signs at $v$ are jointly independent
and fair conditional on $\mathcal H_{v-1}$, as required.
\end{proof}

\begin{samepage}
We will apply the following probability estimate both to sequential
choices and to functions of the independent edge data. Its interval
form gives the constants used in Step~2 of the main proof.

\begin{lemma}[Conditional ranges and bounded differences]\label{lem:random_ranges}
Let $X_1,\ldots,X_k$ be exposed in order, and let $\mathcal H_j$ denote
the information after step $j$. Conditional on the earlier information,
suppose $X_j$ lies in an interval of length $\ell_j$,
where the interval is known before $X_j$ is revealed. If
$\sum_j\ell_j^2\le B$ on every path for a deterministic $B>0$, then,
for $a>0$,
\[
 \Prb\left(\sum_{j=1}^k
       \bigl(X_j-\E[X_j\mid\mathcal H_{j-1}]\bigr)<-a\right)
 \le e^{-2a^2/B}.
\]
In particular, if $f$ is a function of independent inputs and changing
input $j$ changes $f$ by at most $c_j$, with $\sum_jc_j^2>0$, then
\[
\Prb(f<\E f-a)\le e^{-2a^2/\sum_jc_j^2}.
\]
If all $c_j$ are zero, $f$ is constant and this probability is zero.
\end{lemma}
\end{samepage}
\begin{proof}
A centered random variable $Z$ supported on an interval of length $\ell$
satisfies
$\E e^{\lambda Z}\le e^{\lambda^2\ell^2/8}$.
To see this, let $\psi(\lambda)=\log\E e^{\lambda Z}$.
Then $\psi(0)=\psi'(0)=0$, and $\psi''(\lambda)$ is the variance under
the distribution that reweights each outcome by $e^{\lambda Z}$.
This distribution has the same support. Every random variable in an
interval of length $\ell$ has variance at most $\ell^2/4$: subtract the
midpoint, bound the second moment by $\ell^2/4$, and use that variance
is no larger than this second moment. Integrating the bound on
$\psi''$ twice proves the exponential estimate.

Apply it conditionally at each step, with either sign of $\lambda$.
Successive conditional expectations and the bound on the sum of squared
interval lengths give
\[
 \Prb\left(\sum_j
       (X_j-\E[X_j\mid\mathcal H_{j-1}])<-a\right)
 \le \exp(-\lambda a+\lambda^2B/8).
\]
Taking $\lambda=4a/B$ proves the first assertion.
For the second assertion, reveal the independent
inputs one at a time and consider the successive conditional expectations
of $f$. Their increments have conditional mean zero and conditional
range at most $c_j$. Indeed, two possible values of input $j$ change
$f$ by at most $c_j$ for every fixed choice of the remaining inputs,
and hence also after averaging those inputs. The first assertion applies.
\end{proof}

An average gain may include an arbitrary unit-clause contribution.
The next lemma shows that such a shift cannot remove the useful absolute
fluctuations of a symmetric signed sum. Its second assertion converts a
variance estimate into the average gap needed below.

\begin{lemma}[Absolute gains and second moments]\label{lem:random_moments}
Let $Z=\sum_i a_i\rho_i$ for independent fair signs $\rho_i$, and put
$Y=\sum_i a_i^2$. For every real $b$ and every $K>0$,
\[
 \E\min\{|b+Z|,K\}
 \ge\frac12\left(\sqrt{Y/3}-\frac{Y}{4K}\right).
\]
For any nonnegative random variable $U$ with $\E U>0$ and finite second
moment,
\[
 \E\sqrt U\ge\frac{(\E U)^{3/2}}{(\E U^2)^{1/2}}.
\]
\end{lemma}
\begin{proof}
Independence and zero means give $\E Z^2=Y$ and
\[
 \E Z^4=\sum_i a_i^4+6\sum_{i<j}a_i^2a_j^2\le3Y^2.
\]
The H\"older inequality, with exponents $3/2$ and $3$, gives
\[
 \E|Z|^2
 =\E\bigl(|Z|^{2/3}|Z|^{4/3}\bigr)
 \le(\E|Z|)^{2/3}(\E|Z|^4)^{1/3}.
\]
Rearranging yields $\E|Z|\ge\sqrt{Y/3}$ when $Y>0$;
the claim is immediate when $Y=0$.
Write $r_+=\max\{r,0\}$. Then $(z-K)_+\le z^2/(4K)$ for $z\ge0$,
because for $z>K$
the inequality is equivalent to $(z-2K)^2\ge0$.
It follows that
\[
 \E\min\{|Z|,K\}\ge\sqrt{Y/3}-Y/(4K).
\]
The triangle inequality gives $|b+z|+|b-z|\ge2|z|$.
After truncation this implies
$\min\{|b+z|,K\}+\min\{|b-z|,K\}\ge\min\{|z|,K\}$.
Symmetry of $Z$ makes the two expectations on the left equal,
proving the first assertion.

For the second assertion, apply the same H\"older inequality to
$(\sqrt U)^{2/3}U^{2/3}$:
\[
 \E U\le(\E\sqrt U)^{2/3}(\E U^2)^{1/3}.
\]
Rearranging proves the claim.
\end{proof}

We can now bound the accumulated requirements $A_T$ in the main proof.
There are two sources of variation: the requirement at an individual
branching step, and the amount of conditional variance available across
the selected variables. The two terms in the following bound account
for these separately.

\begin{lemma}[Accumulated requirements]\label{lem:random_requirements}
With $A_T,c_1,T,r$ as in the proof of Theorem~\ref{thm:weighted},
for all $a,b>0$,
\[
 \Prb\left(A_T<\frac{c_1T}{26}-a-b\right)
 \le e^{-2a^2/(Tm_0^2)}+e^{-169b^2/(m_0^2r)}.
\]
\end{lemma}
\begin{proof}
Use the exposure of Lemma~\ref{lem:random_exposure}.

\par\medskip\noindent\textbf{1. A lower bound at each step.}
Put
\[
 \widetilde X_t=[t\in I]\,\omega_t
            \min\{|\Delta_t|,m_0\sqrt{d_t}\}.
\]
Here $[t\in I]$ is one when a constraint is created at step $t$,
and zero otherwise. Thus $0\le\widetilde X_t\le m_0$ and
$A_T\ge\sum_{t\le T}\widetilde X_t$.
Conditional on $\mathcal H_{t-1}$, the gain is the fixed shift $b_t$
plus a symmetric signed sum of variance $Y_t$.
The bit $x_t$ is independent of this gain and chooses the nonpreferred
value with probability one half. Lemma~\ref{lem:random_moments} gives
\begin{align*}
 \E[\widetilde X_t\mid\mathcal H_{t-1}]
 &\ge\frac{\omega_t}{4}
       \left(\sqrt{Y_t/3}-\frac{Y_t}{4m_0\sqrt{d_t}}\right)\\
 &\ge\frac{\omega_t\sqrt{Y_t}}4
             \left(\frac1{\sqrt3}-\frac14\right)
 \ge \mu_t:=\frac{\omega_t\sqrt{Y_t}}{13}.
\end{align*}
The middle inequality uses $Y_t\le m_0^2d_t$.
Lemma~\ref{lem:random_ranges}, with interval lengths $m_0$, therefore gives
\[
 \Prb\left(A_T<\sum_{t\le T}\mu_t-a\right)
 \le e^{-2a^2/(Tm_0^2)}.
\]

\par\medskip\noindent\textbf{2. The average available variance.}
For an earlier neighbor $v<t$, condition on the base multiset and let
$q_+$ and $q_-$ be its signed sums at $t$ over the two polarities at $v$.
The coefficient $a_{tv}$ takes these two values with equal probabilities.
Thus
\[
 \E[a_{tv}^2\mid K_{\{t,v\}}]
 =\frac{q_+^2+q_-^2}{2}
 \ge\frac{(q_++q_-)^2}{4}.
\]
Averaging and applying the model's third condition gives
$\E a_{tv}^2\ge\kappa/4$. For a later neighbor, $a_{tv}$ is half
the endpoint signed sum, giving the same lower bound directly.
The terms of $Y_t=\sum_{v\sim t}a_{tv}^2$ depend on independent edge
data, and each lies in $[0,m_0^2]$. Consequently
\[
 \E Y_t\ge\frac{\kappa d_t}{4},\qquad
 \operatorname{Var}(Y_t)\le\frac{m_0^4d_t}{4}.
\]
By Lemma~\ref{lem:random_moments},
\begin{align*}
 \E\sqrt{Y_t}
 &\ge\frac{\sqrt{\E Y_t}}
              {\sqrt{1+\operatorname{Var}(Y_t)/(\E Y_t)^2}}\\
 &\ge
 \frac{\sqrt{\kappa d_t}/2}
      {\sqrt{1+4m_0^4/(\kappa^2d_t)}}
 \ge c_1\sqrt{d_t}.
\end{align*}
We used $d_t\ge1$ in the last step.
The lower weight bound~\eqref{eq:random_weight_bounds} now implies
$\E\sum_{t\le T}\mu_t\ge c_1T/26$.

\par\medskip\noindent\textbf{3. Concentration over the edge data.}
Regard each base multiset and its two $\theta$ signs as one independent
edge input. Changing the input on $\{u,v\}$ affects only the terms
$\mu_u,\mu_v$ that belong to the prefix. In each affected square root,
only one squared coefficient changes, with absolute coefficient in
$[0,m_0]$. For $R\ge0$ and $a,b\in[0,m_0]$,
$|\sqrt{R+a^2}-\sqrt{R+b^2}|\le|a-b|\le m_0$.
Thus an edge changes the sum by at most
$m_0(\omega_u+\omega_v)/13$.
The sum of squared change bounds is at most
\[
 \frac{m_0^2}{169}\sum_{\{u,v\}\in E_H}(\omega_u+\omega_v)^2
 \le\frac{2m_0^2}{169}\sum_v\omega_v^2d_v
 \le\frac{2m_0^2r}{169}.
\]
Lemma~\ref{lem:random_ranges} gives
\[
 \Prb\left(\sum_{t\le T}\mu_t<c_1T/26-b\right)
 \le e^{-169b^2/(m_0^2r)}.
\]
Combining this with the bound from the first step proves the lemma.
\end{proof}

The last estimate controls how much the later values can help satisfy
the accumulated constraints. A shared variable may affect many constraints,
but the relevant endpoint signs on its different edges are still
independent when that variable is reached. This is the estimate for
$Q_T$ used in Step~2 of the main proof.

\begin{lemma}[Favorable contributions from later values]\label{lem:random_compensation}
For every $c>0$,
\[
 \Prb(Q_T<-c)\le e^{-2c^2/(m_0^2T)}.
\]
\end{lemma}
\begin{proof}
By Lemma~\ref{lem:random_exposure}, the contribution from $\{t,v\}$,
where $t\le T$ and $t<v$, is
\[
 [t\in I]\,\omega_t(2p_t-1)
      \rho_{t,\{t,v\}}\eta_{tv}\theta_{v,\{t,v\}}.
\]
Order these terms by their later endpoint $v$.
Their coefficients are known before step $v$, while the $\theta$ signs
at $v$ are conditionally independent fair signs.
We may reveal these signs one at a time and then reveal $x_v$;
the corresponding $\rho_{v,e}$ are then determined.
This is equivalent to the earlier exposure and gives a conditional
mean-zero increment for each displayed term.

Its conditional range has length at most $\omega_t m_0$.
The sum of squared range lengths is bounded by
\[
 m_0^2\sum_{t\le T}\omega_t^2
        |\{v>t:\{t,v\}\in E_H\}|
 \le m_0^2\sum_{t\le T}\omega_t^2d_t
 \le m_0^2T.
\]
Lemma~\ref{lem:random_ranges} proves the claim. This argument does not
require independence between $A_T$ and $Q_T$.
\end{proof}

\subsection{Independent clauses on an arbitrary support}

The general pair-distribution condition has a particularly simple
consequence for independent clauses. It allows any prescribed set of
variable pairs; neither a density condition nor regularity is needed.

\begin{corollary}[Arbitrary supports]\label{cor:support}
Fix $q\in(0,1)$ independently of $n$. Fix any set of variable pairs and
include each of the four clauses on each pair independently with
probability $q$. Add unit clauses fixed independently of the binary draw.
An optimal assignment can be computed by a deterministic polynomial-space
algorithm with expected running time
$\Ostar(2^n e^{-c_0(q)n})$, where $c_0(q)>0$ depends only on $q$.
The algorithm is exact on every input.
\end{corollary}
\begin{proof}
Apply Theorem~\ref{thm:weighted} with $m_0=4$ and
$\kappa=4q(1-q)$. The theorem already gives an algorithm whose only input
is $F$. Its prefix term is absorbed into $\Ostar(2^n e^{-c_0(q)n})$
because $c_0(q)<(\ln2)/2$. The same bound holds if the selected pair set
is itself random and chosen independently of the clause signs: condition
on that set, apply the uniform bound of the theorem, and then average.
\end{proof}

The saving is uniform over supports and their degree sequences, but
not over $q$ tending to zero or one. This fixed-$q$ result needs no
logarithmic degree threshold. Unit clauses chosen after observing the
binary draw are not covered by its exposure argument. The extension
below permits that stronger choice under a density assumption.
Complementary pairs alone have endpoint signed sum zero and fail the
third condition of Definition~\ref{def:invariant}; the Max-Cut results
in Section~\ref{sec:maxcut} treat such inputs by a different analysis.


\subsection{An extension using assignment counts}\label{sec:countalgorithm}

The preceding argument relies on fresh symmetric endpoint signs.
We next allow a systematic preference for equal or opposite literal
polarities on a pair, and unit clauses chosen after the binary formula
is known. To handle the systematic part, we guess one integer: the total
number of variables set to one. There are only $n+1$ possibilities.
Within each guess we again form comparison constraints and enumerate
completions satisfying their sum.

We introduce the additional notation here because the weighted search
above does not need it. Expanding the truth indicator of each clause
expresses the number of satisfied clauses as
\begin{equation}
 F(x)=c+\sum_v a_vx_v+\sum_{u<v}J_{uv}x_ux_v,
 \qquad x\in\bits^n.
 \label{eq:quadratic}
\end{equation}
The matrix $J$ is symmetric with zero diagonal.
A binary clause with equal literal polarities contributes $-1$ to
$J_{uv}$; one with opposite polarities contributes $+1$.
Unit clauses change only the constant and linear terms.
The gain from changing $x_t$ from zero to one is therefore
\[
 \mathrm{gain}_t(x)=a_t+\sum_{w\ne t}J_{tw}x_w.
\]
These quantities are computed by scanning clauses.

Fix a rational number $\beta$ representing the common part of the pair
coefficients. For a guess $h\in\{0,\ldots,n\}$ define
\begin{equation}
 H_t^{h,\beta}(x)=\mathrm{gain}_t(x)
          +\beta\left(h-\sum_{w\ne t}x_w\right)
       =a_t+\beta h+\sum_{w\ne t}(J_{tw}-\beta)x_w.
 \label{eq:countgain}
\end{equation}
The second expression explains the adjustment: once $h$ is fixed, the
common pair contribution $\beta$ is replaced by the known constant
$\beta h$. The remaining coefficients are the deviations $J_{tw}-\beta$.
When $h=\sum_wx_w$ is the correct guess,
$H_t^{h,\beta}(x)=\mathrm{gain}_t(x)+\beta x_t$.
We account for this last term by allowing a slack of $|\beta|$ in a
comparison.

\paragraph{The algorithm.}
Choose a prefix of $T$ variables in a fixed order and partition it into
ordered, nonempty rounds $B_1,\ldots,B_\ell$.
For the random formulas below each round contains one variable.
We allow larger rounds because Section~\ref{sec:maxcut} uses the same
procedure with two rounds.
Within round $B_j$, choose each preference using only the assigned values
in earlier rounds, whose union is denoted by $B_{<j}$.
The values of all other variables, including those in the current round,
are averaged uniformly. This average is computed exactly by substituting
$1/2$ in~\eqref{eq:countgain}; no sampling is involved.

If this average is nonnegative, the preferred value $p_t$ is one;
otherwise it is zero. When the prefix takes the nonpreferred value,
add the comparison
\[
 (2p_t-1)H_t^{h,\beta}(x)\le|\beta|.
\]
The code stores the sum of left-hand sides minus right-hand sides in
$L(x)$, so the single inequality for enumeration is $L(x)\le0$.
The global best score and assignment again are $M,\alpha^\star$.

\begin{algorithm}[H]
\caption{\textnormal{\textsc{CountComparisons}}$(F,\beta,B_1,\ldots,B_\ell)$}
\label{alg:counts}
\begin{algorithmic}[1]
\STATE $M\gets-\infty$; $\alpha^\star\gets\emptyset$.
\STATE $(F,s_0)\gets\textnormal{\textsc{Preprocess}}(F)$.
\STATE Set $\mathcal H\gets\{0,\ldots,n\}$ if $\beta\ne0$, and $\mathcal H\gets\{0\}$ otherwise.
\COMMENT{One run suffices when the comparison does not depend on $h$.}
\FOR{each $h\in\mathcal H$}
  \FOR{each prefix assignment $\alpha$ on $B_1\cup\cdots\cup B_\ell$}
    \STATE $L(x)\gets0$.
    \FOR{each round $B_j$, and each variable $t\in B_j$}
      \STATE Compute the average $\Delta_t$ of $H_t^{h,\beta}$ with $x_{B_{<j}}=\alpha_{B_{<j}}$.
      \STATE Set $p_t\gets1$ if $\Delta_t\ge0$, and $p_t\gets0$ otherwise.
      \IF{$\alpha_t\ne p_t$}
        \STATE $L(x)\gets L(x)+(2p_t-1)H_t^{h,\beta}(x)-|\beta|$.
      \ENDIF
    \ENDFOR
    \FOR{each suffix $y$ with $L(\alpha,y)\le0$, by Lemma~\ref{lem:enumeration}}
      \STATE Let $q\gets s_0+F(\alpha,y)$.
      \STATE If $q>M$, set $M\gets q$ and $\alpha^\star\gets\alpha\cup y$.
    \ENDFOR
  \ENDFOR
\ENDFOR
\STATE \textbf{return} $(M,\alpha^\star)$.
\end{algorithmic}
\end{algorithm}

Each guess uses a separate set of comparisons, and guesses are processed
consecutively. Enumeration need not enforce that a candidate has exactly
$h$ ones: every candidate is evaluated in the actual formula.
For the biased model below, use singleton rounds through
$T=\lfloor n/2\rfloor$ and
\[
 \beta=\binom n2^{-1}\sum_{u<v}J_{uv}.
\]
This is the empirical mean of the signed pair coefficients.
For the zero-mean models we may instead set $\beta=0$.
Neither choice requires an eigenvalue or eigenvector calculation.

\paragraph{Correctness.}
The count guess changes the comparisons, so we must check that a correct
guess preserves an optimum. The following lemma accounts for the added
slack. This argument is independent of the distribution of $F$ and also
justifies the round version used later.

\begin{lemma}[Correctness of count comparisons]\label{lem:count_correctness}
For every $F$, every rational $\beta$, and every ordered prefix partition,
\textnormal{\textsc{CountComparisons}} returns an optimum score and an
assignment attaining it.
\end{lemma}
\begin{proof}
Fix an optimal complete assignment $x^*$ and its correct count
$h=\sum_wx_w^*$. If $x_t^*\ne p_t$, changing just this variable to $p_t$
cannot increase the score, so
$(2p_t-1)\mathrm{gain}_t(x^*)\le0$.
For the correct count,~\eqref{eq:countgain} therefore gives
\[
 (2p_t-1)H_t^{h,\beta}(x^*)
 =(2p_t-1)\mathrm{gain}_t(x^*)+(2p_t-1)\beta x_t^*
 \le|\beta|.
\]
Every comparison on its prefix holds, and hence so does their sum.
The corresponding suffix is enumerated in this run.
When $\beta=0$, the comparisons do not depend on $h$, so the single run
also retains $x^*$. Preprocessing preserves all gains and records its
constant score change. Scoring candidates with that constant included
proves the claim.
\end{proof}

\subsection{Random clauses with weaker symmetry assumptions}

The next result applies the count adjustment to two explicit distributions.
It permits probabilities tending to zero and unit clauses chosen after
seeing the binary formula. Its density threshold supplies enough random
variation for the accumulated comparisons to remain effective.
These additional freedoms distinguish it from the fixed-$q$ result above.

\begin{corollary}[Random formulas with several clauses per pair]
\label{thm:randomclauses}
There are absolute constants $C,\gamma_0>0$ such that, for sufficiently
large $n$, each of the following models with $d\ge C\log n$ admits a
deterministic polynomial-space exact algorithm taking time
$\Ostar(2^{(1-\gamma_0)n})$ with probability at least $1-1/n$.
\begin{enumerate}
\item \textbf{Independent clauses.} On every variable pair, include each
equal-polarity clause with probability $q(1+\epsilon)$ and each
mixed-polarity clause with probability $q(1-\epsilon)$, independently across
all clauses. Here $\epsilon\in[-1,1]$,
$q(1+|\epsilon|)\le1/2$, and $d=4q(n-1)$.
\item \textbf{A random set of pairs.} Select each pair for
$S\subseteq\binom{[n]}2$ independently with probability $p$. On each selected
pair, include each of the four clauses independently with probability $1/4$;
include none on the other pairs. Here $d=p(n-1)$.
\end{enumerate}
Both models allow arbitrary additional unit clauses, even chosen after the
binary clauses are known. The algorithm is exact on every input; the
probability qualifies its running time.
\end{corollary}

Here $d=2\E m_2/n$, so the expected number of binary clauses per variable
is $d/2$, and $\log$ is natural. In the first model with $\epsilon=0$,
and in the second model, the mean pair coefficient is zero and one run
with $\beta=0$ suffices. The original strict pruning constraints may also
be used in these zero-mean cases: they retain an optimum by the earlier
correctness argument and are stronger than the comparisons with zero slack.

\paragraph{How the running-time argument applies.}
Appendix~\ref{app:membership} provides an independent-pair criterion
(Proposition~\ref{prop:pairmembership}) for the accumulation theorem in
Appendix~\ref{app:accumulation}. It reduces the task here to checking the
mean, variance, and boundedness of the pair coefficients $J_{uv}$.
We perform those checks explicitly below. The criterion's probability
event depends only on $J$, while its running-time conclusion is uniform
over the linear coefficients. This is why even unit clauses chosen after
the random draw are allowed.

\begin{proof}[Proof of Corollary~\ref{thm:randomclauses}]
Use \textnormal{\textsc{CountComparisons}} with singleton rounds through
$T=\lfloor n/2\rfloor$. Its correctness follows from
Lemma~\ref{lem:count_correctness}.

\par\medskip\noindent\textbf{1. State the condition to be checked.}
Set $\tau=1/2$ and $C_0=10^9/\tau$, as in
Proposition~\ref{prop:pairmembership}. All pairs have the same distribution,
so its variance-ratio parameter is $\chi=1$.
Write $\mu=\E J_{uv}$ and $\nu=\operatorname{Var}(J_{uv})$.
Let $B$ bound $|J_{uv}-\mu|$, and put $\sigma_0=|\mu|$.
It suffices to verify independence across pairs and
\[
 n\nu\ge C_0(B^2\log n+\sigma_0^2).
\]
For sufficiently large $n$, the proposition then supplies all the
conditions needed by the sequential case of Theorem~\ref{res:counts}.

\par\medskip\noindent\textbf{2. Independent clause inclusions.}
Let $B_{++},B_{--},B_{+-},B_{-+}$ be the four independent inclusion
indicators on one pair. By~\eqref{eq:quadratic},
\[
 J=B_{+-}+B_{-+}-B_{++}-B_{--}.
\]
Taking expectations and summing the variances of independent indicators gives
\[
 \mu=-4q\epsilon,\qquad
 \nu=4q-4q^2(1+\epsilon^2)\ge2q.
\]
The last inequality follows because every inclusion probability is at
most $1/2$, so each indicator has variance at least half its mean.
Since $|J|\le2$ and $|\mu|\le2$, take $B=4$ and $\sigma_0\le2$.
Also $n\nu\ge d/2$.
For sufficiently large $n$, the bound $d\ge34C_0\log n$ implies
\[
 n\nu\ge17C_0\log n\ge C_0(16\log n+4),
\]
as required. Use the empirical mean $\beta$; when $\epsilon=0$,
the known zero mean $\beta=0$ is also allowed.

\par\medskip\noindent\textbf{3. A random support.}
On one pair write $J=EZ$, where $E$ is its Bernoulli-$p$ selection
indicator, and $Z$ is the difference between the sums of two positive
and two negative independent Bernoulli-$1/4$ indicators.
The variables $E$ and $Z$ are independent. Therefore
\[
 \E J=0,\qquad
 \nu=\E J^2=p\,\E Z^2=\frac{3p}{4},\qquad
 |J|\le2.
\]
The pair coefficients are independent, even though the four clause
indicators on a pair share its selection indicator.
We may take $B=2$, $\sigma_0=0$, and $\beta=0$.
Since $n\nu\ge3d/4$, the condition $d\ge6C_0\log n$ is sufficient.

\par\medskip\noindent\textbf{4. Obtain a fixed exponential saving.}
In both cases Proposition~\ref{prop:pairmembership} gives its event
with probability at least $1-1/n$, with $K=9$ and
$g=\tau^{3/2}/12$. In the notation of Theorem~\ref{res:counts},
the sequential case then has
\[
 \alpha=\frac{81}{\sqrt\tau},\qquad
 c_*=\frac{\tau^3}{144\cdot64\cdot81^2\cdot9}.
\]
That theorem bounds the time by
$\Ostar((n+1)(2^{\tau n}+2^ne^{-c_*n}))$.
The factor $n+1$ accounts for count guesses and is polynomial;
it is absent when $\beta=0$. At $\tau=1/2$,
$\gamma_0=10^{-10}$ is smaller than both $1-\tau$ and $c_*/\ln2$,
so it gives the claimed binary-exponent saving.
The common threshold $C=34C_0$ suffices for the two models.
The event concerns only the binary pair coefficients, so the conclusion
holds for every choice of additional units on that event.
\end{proof}


\section{Max-Cut families}\label{sec:maxcut}

Two familiar graph families illustrate accumulation without an initial
literal gap: regular graphs with small nontrivial eigenvalues, and graphs
drawn from the stochastic block model. Encode an edge $\{u,v\}$ by
$(u\lor v)$ and $(\neg u\lor\neg v)$. It contributes
$1+[x_u\ne x_v]$, so maximizing the resulting formula is exactly Max-Cut
up to the constant number of edges. Every variable has equal positive and
negative binary literal degrees. Additional units are arbitrary linear
fields, and both results below allow them even when they depend on the graph.

The regular-graph algorithm uses only gains, degrees and a guessed count;
its spectral hypothesis is used in the analysis. For hidden communities,
the algorithm also computes approximate spectral directions. This is an
explicit preprocessing step, described after the class statement.

\subsection{Max-Cut on graphs with small nontrivial eigenvalues}

For a $d$-regular simple graph $G$, let
$\lambda(G)=\max_{i\ge2}|\lambda_i(A_G)|$, retaining multiplicities.

\begin{corollary}[Max-Cut with external fields]\label{cor:maxcut}
Fix $\lambda_0\ge1$ and $0<\varepsilon\le1/2$. There is
$c=c(\lambda_0,\varepsilon)<2$ such that every $d$-regular graph satisfying
\[
 d\le(1-\varepsilon)n,\qquad \lambda(G)\le\lambda_0\sqrt d
\]
admits a deterministic polynomial-space Max-Cut algorithm of time
$\Ostar(c^n)$. The same bound holds after adding any multiset of unit
clauses to the Max 2-SAT encoding with $(u\lor v)$ and
$(\neg u\lor\neg v)$ on each edge.
\end{corollary}

For $d\ge64\lambda_0^2$, use Algorithm~\ref{alg:counts} with
$\beta=-2d/(n-1)$, $T=\lfloor n/2\rfloor$, and two prefix rounds
of sizes $\lfloor n/4\rfloor$ and $T-\lfloor n/4\rfloor$.
For smaller $d$, branch outside a greedy independent set and optimize
its vertices separately. The solver computes no spectral quantities;
the eigenvalue condition is used only in the following analysis.

\begin{proof}
If $d<64\lambda_0^2$, greedily find an independent set of size at least
$n/(d+1)$, branch on its complement, and optimize its variables independently.
This already gives a uniform fixed saving in this range. Assume henceforth
$d\ge64\lambda_0^2$ and $n$ is sufficiently large in terms of the fixed
parameters; finitely many smaller sizes can be absorbed in the constant.

Apply the round version of Theorem~\ref{res:counts} with $\tau=1/2$,
two prefix rounds of sizes $\lfloor n/4\rfloor$ and
$\lfloor n/2\rfloor-\lfloor n/4\rfloor$ in any order, and
\[
 J=-2A_G,\quad P=\beta\ind\ind^{\mathsf T},\quad
 \beta=-2d/(n-1),\quad N=-2A_G-\beta(\ind\ind^{\mathsf T}-I).
\]
Use exact counts $0,\ldots,n$, so $s_t=|\beta|$.
The matrix $N$ annihilates $\ind$ and has other eigenvalues
$-2\lambda_i(A_G)+\beta$. Since $|\beta|\le2$,
\[
 \|N\|\le3\lambda_0\sqrt d,\qquad
 D_N=4d(1-d/(n-1))\ge2\varepsilon d
\]
for $n\ge2/\varepsilon$. Thus $K=9\lambda_0^2/(2\varepsilon)$ suffices.

An edge between the two prefix rounds contributes at least
$\varepsilon^2$ to $W_t$. The expander mixing inequality gives, for $n\ge80$,
\[
 e(B_1,B_2)\ge \frac{d|B_1||B_2|}{n}
           -\lambda(G)\sqrt{|B_1||B_2|}
 \ge\frac{dn}{32}-\frac d2\ge\frac{dn}{40}.
\]
Indeed $|B_1||B_2|\ge(n/4-1)^2$,
$\sqrt{|B_1||B_2|}\le n/4$, and $\lambda(G)\le d/8$.
The mixing inequality itself follows by writing each indicator vector
as its constant component plus a vector perpendicular to $\ind$ and
applying the spectral bound to the latter components. Consequently
\[
 \sum_{t\in B_2}\sqrt{W_t}
 \ge\varepsilon\sum_{t\in B_2}\sqrt{\deg_{B_1}(t)}
 \ge\frac{\varepsilon e(B_1,B_2)}{\sqrt d}
 \ge\frac{\varepsilon n\sqrt d}{40}.
\]
Since $\sqrt{D_N}\le2\sqrt d$, choose $g=\varepsilon/(320\sqrt3)$.
Finally $\sum_{t\le T}s_t\le2d=o(n\sqrt d)$ uniformly for $d\le n$,
so the slack condition holds for all sufficiently large $n$.
There are $n+1$ runs and all theorem parameters are fixed.
\end{proof}

\paragraph{Concrete graph families.}
Paley graphs on $q\equiv1\pmod4$ vertices, for prime powers $q$, have
degree $(q-1)/2$ and nontrivial eigenvalues $(-1\pm\sqrt q)/2$
\cite[Section~9.1.2]{BH12}. Thus $\lambda_0=1$, $\varepsilon=1/2$
suffice for all sufficiently large $q$. Connected nonbipartite Ramanujan
graphs satisfy $\lambda(G)\le2\sqrt{d-1}$, so they give another family
when the density restriction holds. Bipartite Ramanujan graphs have the
additional eigenvalue $-d$ and do not satisfy this hypothesis at growing
degree; ordinary Max-Cut on bipartite graphs is of course immediate.
For each fixed $a>0$, uniform random $d$-regular graphs with
$n^a\le d\le n/2$ satisfy $\lambda(G)\le C_a\sqrt d$ with probability
at least $1-1/n$ by Tikhomirov and Youssef~\cite{TY19}. The corollary
therefore supplies a fixed saving on that random graph family as well.
These are applications of one deterministic class theorem, not separate
algorithm designs. No priority claim is made here.


\subsection{The stochastic block model}\label{sec:sbm}

Fix a partition $V=V_1\cup\cdots\cup V_k$ and a symmetric probability
matrix $\Pi=(\pi_{ab})_{a,b\le k}$. Include each unordered pair of distinct
vertices independently, with probability $\pi_{ab}$ when its endpoints
belong to $V_a,V_b$. This is the stochastic block model~\cite{Abbe18}.
The partition can be arbitrary and need not be balanced. Alternatively,
it can be drawn first; the result below holds uniformly conditional on it.
The number $k$ is fixed, while the probabilities may depend on $n$.

Define the smallest and largest edge variances by
\[
 v_{\min}=\min_{a,b}\pi_{ab}(1-\pi_{ab}),\qquad
 v_{\max}=\max_{a,b}\pi_{ab}(1-\pi_{ab}).
\]
These conditions measure the amount of random variation, not how well
the communities can be distinguished.

\begin{theorem}[Exact Max-Cut on the stochastic block model]\label{thm:sbm}
Fix $k\ge1$ and $\chi\ge1$. There are constants $C=C(k,\chi)$ and
$\gamma=\gamma(k,\chi)>0$ such that, whenever
\begin{equation}
 0<v_{\min},\qquad v_{\max}\le\chi v_{\min},\qquad
 n v_{\min}\ge C\log n,                       \label{eq:sbmvariance}
\end{equation}
Max-Cut on a graph from this model can be solved by a deterministic
polynomial-space exact algorithm taking
$\Ostar(2^{(1-\gamma)n})$ time with probability $1-o(1)$.
The algorithm is given the fixed class parameters $k,\chi$, but neither
the partition nor the probabilities $\pi_{ab}$. The same statement holds
with arbitrary additional unit clauses, even chosen after the graph is
drawn. Exactness holds on every input; only the time bound is probabilistic.
\end{theorem}

\begin{proof}
We verify the auxiliary discovery lemma in
Appendix~\ref{app:discovery}. In the Max 2-SAT encoding,
$J_{uv}=-2A_{uv}$ for the adjacency indicators $A_{uv}$. Conditional on
the partition, these coefficients are independent, with
\[
 \E J_{uv}=-2\pi_{ab},\qquad
 \operatorname{Var}J_{uv}=4\pi_{ab}(1-\pi_{ab}),\qquad
 |J_{uv}-\E J_{uv}|\le2.
\]
Let $Z_{ua}=[u\in V_a]$. The matrix
$\bar A=-2Z\Pi Z^{\mathsf T}$ agrees with the means off the diagonal,
has rank at most $k$, has diagonal entries of absolute value at most $2$,
and satisfies $\|\bar A\|\le2n$ by its maximum absolute row sum.
In Lemma~\ref{cor:planted}, take
\[
 r=k,\quad \nu=4v_{\min},\quad B=B_0=\mu=2.
\]
Its threshold becomes
$4n v_{\min}\ge C_0(4\log n+8)$, which follows from
\eqref{eq:sbmvariance} after increasing $C$.
The lemma's algorithm does not use $\bar A$ or its probabilistic
parameters. Its good event depends only on the pair coefficients and its
bound is uniform over linear fields, proving the assertion about units.
\end{proof}

\paragraph{The usual within/between model.}
When $\pi_{aa}=p$ and $\pi_{ab}=q$ for $a\ne b$, the assumptions reduce to
\[
 \max\{p(1-p),q(1-q)\}\le\chi\min\{p(1-p),q(1-q)\},\qquad
 n\min\{p(1-p),q(1-q)\}\ge C\log n.
\]
No separation between $p$ and $q$ is required; in particular $p=q$ is
allowed. Concrete regimes are:
\begin{itemize}
\item \emph{Dense blocks:} all $\pi_{ab}\in[\delta,1-\delta]$ for a
fixed $0<\delta\le1/2$. The hypotheses hold for sufficiently large $n$,
with constants depending only on $k,\delta$.
\item \emph{Sparse blocks:} $\pi_{ab}=c_{ab}\rho$ with fixed
$0<a\le c_{ab}\le b$, $\rho=o(1)$, and $n\rho\ge C'\log n$.
\item \emph{Almost complete blocks:} $\pi_{ab}=1-c_{ab}\rho$ with the
same bounds on $c_{ab}$ and $\rho$.
\end{itemize}
For the last two regimes, eventually
$a\rho/2\le\pi_{ab}(1-\pi_{ab})\le b\rho$; one can take
$\chi=2b/a$ and increase $C'$. The saving is fixed throughout each regime.

\paragraph{How the hidden-community algorithm uses the graph.}
The group probabilities create a systematic part of the comparisons.
The algorithm accounts for it without recovering the community labels:
\begin{enumerate}
\item Compute a rational approximate diagonalization of $J=-2A_G$ and
retain at most $k$ directions with the largest absolute diagonal coefficients.
This takes polynomial time; the precise construction is in
Appendix~\ref{app:discovery}.
\item Enumerate polynomially many guesses for the projections of an
assignment onto those directions. Each guess defines adjusted gains and
an explicit error allowance.
\item In each run use two prefix rounds, add hard comparisons in the
nonpreferred branches, and enumerate the suffixes satisfying their sum.
Evaluate the original objective and retain the best candidate over all runs.
\end{enumerate}
At least one run preserves an optimum. Constraints are local to their run.
Thus the spectral step changes how gains are adjusted; the search is still
pruned by accumulated hard comparisons. Lemma~\ref{cor:planted} and the
generator in Appendix~\ref{app:accumulation} specify the arithmetic,
grid sizes and all-input fallback exactly.

\paragraph{When the communities are supplied.}
If the partition and rational probabilities of polynomial encoding length
are given, their role has an elementary description. Guess the number
$h_a$ of vertices assigned $1$ in each $V_a$: there are at most $(n+1)^k$
tuples. For $t\in V_a$, use the adjusted gain
\[
 H_t^{\boldsymbol h}(x)=G_t(x)
 -2\sum_{b=1}^k\pi_{ab}
       \left(h_b-\sum_{w\in V_b\setminus\{t\}}x_w\right).
\]
On an assignment with these counts it equals
$G_t(x)-2\pi_{aa}x_t$, so the comparison allowance is $2\pi_{aa}$.
The same two-round accumulation proof applies, with total allowance at
most $2n$. The centered independent entries satisfy the required norm and
cross-round mass bounds under~\eqref{eq:sbmvariance}; this is the proof of
Lemma~\ref{cor:planted} with the mean supplied and no approximation error.
Consequently the same type of fixed exponential saving follows using
only group counts and gains. The spectral preprocessing in the theorem
removes the need to supply this information.

\paragraph{Deterministic endpoints and the role of the assumptions.}
The variance hypotheses are sufficient, not necessary. If $p=q=1$,
the graph is a clique and ordinary Max-Cut has value $\lfloor n^2/4\rfloor$.
With units, sort the linear fields: for each count $h$, take the $h$
largest fields and maximize $h(n-h)$ plus their sum. This is polynomial
time. A clique has $\Theta(2^n/\sqrt n)$ optimal Boolean assignments,
so enumerating every optimum would give no fixed exponential saving;
count compression avoids that enumeration. More generally, when a supplied
fixed-block model has only probabilities $0$ or $1$, the objective depends
on the group counts, and sorting fields within each group gives an
$n^{O(k)}$ algorithm. These endpoints use repetition, whereas
Theorem~\ref{thm:sbm} uses random variation. The theorem does not assert
that arbitrary intermediate or highly unequal variances are covered.

\section{Repeated interactions across a cut}\label{sec:structure}

\subsection{Several repeated interaction profiles}

An interaction profile toward $B$ is the vector $(J_{uv}:v\in B)$.
Suppose $A=A_1\mathbin{\dot\cup}\cdots\mathbin{\dot\cup}A_t$ and all vertices
of $A_i$ have the same profile toward $B$. Put
\[
 P=\prod_{i=1}^t(|A_i|+1),\qquad
 \Gamma=|A|-\log_2 P.
\]

\begin{theorem}[Compression by group counts]\label{thm:profiles}
Given $A,B$ and their profile classes, the instance can be solved in
polynomial space and time
\[
 \Ostar\left(2^{n-|B|}+2^{n-\Gamma}\right).
\]
\end{theorem}
\begin{proof}
Fix the outside variables. The effect on $B$ depends on an assignment to
$A$ only through its vector of group counts. For each such vector, find the
best assignment to $A$ with those counts, then enumerate $B$.
Count-restricted generation over all vectors visits exactly $2^{|A|}$
assignments to $A$; enumeration on $B$ costs $P2^{|B|}$.
Stream the vectors and keep one representative at a time, so space stays
polynomial even when $P$ is exponential. Multiply by the outside cost.
\end{proof}

For example, $t=o(n)$ and $|A|,|B|=\Omega(n)$ imply
$\log_2P\le t\log_2(1+|A|/t)=o(n)$ and hence a fixed saving.
This is a statement about repeated profiles across a cut; it does not
require literal cliques within either side.

\subsection{Literal bicliques and a linear number of defects}

Let $A,B$ be disjoint sets, with one prescribed literal for each variable.
A good cross pair contains exactly the prescribed clause. A missing clause,
a different clause, or an additional binary clause makes the pair defective.
All clauses outside the cross pairs and all units are arbitrary.

\begin{corollary}[An exact literal biclique]\label{thm:biclique}
Given an exact literal biclique with $|A|=a$ and $|B|=b$, Max 2-SAT can be
solved in polynomial space and time $\Ostar(2^{n-\min(a,b)})$.
\end{corollary}
\begin{proof}
Complement variables so that the prescribed literals are positive.
All rows of $J_{A,B}$ then agree, so there is one profile class.
Theorem~\ref{thm:profiles} gives
$\Ostar(2^{n-b}+(a+1)2^{n-a})=\Ostar(2^{n-\min(a,b)})$.
\end{proof}

\begin{lemma}[Extracting a core from average defect]\label{lem:core}
If the $A$--$B$ cross pairs contain $q$ defects, there is an exact literal
biclique with both sides of size at least
\[
 k=\left\lfloor\min\left(a,\frac{ab}{a+q}\right)\right\rfloor.
\]
When $A,B$ and their literals are supplied, it can be found in polynomial time.
\end{lemma}
\begin{proof}
Take the $k$ vertices of $A$ with smallest defect counts. Their total number
of defects is at most $kq/a$. Deleting their defective partners leaves at
least $b-kq/a\ge k$ vertices of $B$.
\end{proof}

Thus $a,b=\Omega(n)$ and $q=O(n)$ suffice. The condition allows a linear
number of defective pairs, not an arbitrary fixed fraction of the
$\Theta(n^2)$ cross pairs.

\begin{proposition}[Detection can be paid for]\label{prop:detect}
Fix $0<\delta\le1/8$ and put $k=\lfloor\delta n\rfloor$. The existence of
an exact literal biclique with both sides at least $k$ can be detected in
time $\Ostar(\binom nk2^k)$ and polynomial space. Upon detection, the
instance can be solved in time
\[
 \Ostar\left(2^{(H_2(\delta)+\delta)n}+2^{n-k}\right).
\]
Both exponents have a fixed saving for fixed $\delta>0$.
\end{proposition}
\begin{proof}
Enumerate a set of $k$ variables and its literal choices. Intersect their
good literal neighborhoods and check whether it contains literals on at
least $k$ distinct outside variables. A promised biclique is eventually
found. Run Corollary~\ref{thm:biclique} once, after detection. The usual binomial
entropy bound gives the displayed cost; $H_2(1/8)+1/8<0.669$.
\end{proof}

\subsection{Why bias and cross structure differ}\label{sec:separation}

Let $n=2s$, with equal sides $A,B$. Put positive-positive clauses on every
cross pair and negative-negative clauses on every within-side pair. Every
variable has positive degree $s$ and negative degree $s-1$, so its gap is
only one, but Corollary~\ref{thm:biclique} applies.

Conversely, on the complete graph, choose each incidence to be positive
independently with probability $3/4$. With high probability, every variable
has a positive polarity gap linear in $n$. Yet there is no prescribed
literal biclique with linear sides and only $O(n)$ defects, with high
probability. For any fixed sides and literal choices, each cross pair is
good with probability at most $9/16$, independently across pairs. The
probability that all but $O(n)$ of $\Theta(n^2)$ pairs are good is
$e^{-\Omega(n^2)}$. There are only $e^{O(n)}$ choices of sides and literals,
so a union bound proves the claim.

\section{Collective redundancy in bounded blocks}\label{sec:blocks}

For a block $B$ and two different assignments $a,b\in\bits^B$, the gain
from replacing $a$ by $b$ is
\begin{equation}
 F(b,y)-F(a,y)=\kappa_B(a,b)+\sum_{v\notin B}d_vy_v,
 \qquad d_v=\sum_{u\in B}J_{uv}(b_u-a_u).       \label{eq:blockgain}
\end{equation}
Thus equality of outside effects concerns sums of interaction rows.
Pairwise similar neighborhoods are only one source of such equality.

\begin{proposition}[An unconditional block clause]\label{prop:block}
If $\kappa_B(a,b)+\sum_{v\notin B}\min(0,d_v)>0$, pattern $a$ can be
forbidden. The weak inequality also suffices if the common tie convention
prefers $b$ to $a$. If all $d_v=0$, one of $a,b$ can always be forbidden.
\end{proposition}
\begin{proof}
The expression is the minimum gain over all assignments outside the block.
Fix one order on complete assignments to break objective ties, for example
lexicographic order in a fixed variable order. Among the optima choose the
first assignment in this order. If it contains pattern $a$ and replacing
$a$ by $b$ always increases the score, this contradicts optimality. If the
gain is only nonnegative but the fixed order prefers the replacement for
every outside assignment, either the score increases or an earlier optimum
is obtained. Both are contradictions. Thus the selected optimum avoids
$a$. The same reasoning applies simultaneously to all such exclusions
using that one order. If all $d_v=0$, the comparison is independent of the
outside assignment, so the smaller score, or the later tied pattern, can
be excluded.
\end{proof}

For instance, if four vertices satisfy
$J_{u_1v}+J_{u_2v}=J_{u_3v}+J_{u_4v}$ for every outside $v$, the patterns
$1100$ and $0011$ have identical outside effects. Their worse pattern is
excluded by a clause of length four. No two of the four rows need agree.

\begin{theorem}[Blocks with exceptional outside variables]\label{thm:blocks}
Suppose $B_1,\ldots,B_q$ are disjoint nonempty blocks of size at most $b$.
For each block there are two patterns for which the discrepancy set
$D_i=\{v\notin B_i:d_v\ne0\}$ has size at most $r$.
Given these certificates, the instance can be solved in polynomial space
and time
\[
 \Ostar\left(2^n(1-2^{-b})^{q/(2r+1)}\right).
\]
For constant $b,r$ and $q=\Omega(n)$ this is a fixed saving, at any density.
\end{theorem}
\begin{proof}
Make a graph on the blocks, joining $i,j$ if $D_i$ meets $B_j$ or $D_j$
meets $B_i$. Orient each such witness from its discrepancy source; there
are at most $r$ outgoing witnesses per block. Every induced graph has at
most $r$ times its number of vertices in edges, so it is $2r$-degenerate.
A greedy coloring with $2r+1$ colors gives an independent set $I$ of size
at least $q/(2r+1)$.

Assign all variables outside $\bigcup_{i\in I}B_i$. Every discrepancy
variable of a retained block is now fixed. Its replacement gain is a
constant independent of the other retained blocks, so one of its two
patterns can be excluded. Choose ties consistently. The surviving states
number
\[
 2^{n-\sum_{i\in I}|B_i|}\prod_{i\in I}(2^{|B_i|}-1)
 \le2^n(1-2^{-b})^{|I|}.
\]
Direct enumeration over these block states has this cost and polynomial
space, and preserves an optimum for each outside assignment.
\end{proof}

For fixed $b$, all candidate blocks and pairs of patterns can be checked
in polynomial time. If a disjoint family of size $q$ exists, a maximal
greedy disjoint selection has size at least $q/b$: each chosen block
intersects at most $b$ blocks of the promised disjoint family. Thus the
existential linear-family version also has a polynomial-time certificate
search, with a constant-factor loss in the saving.

This gives a robust structural exit. It does not justify excluding blocks
with an unbounded number of exceptions at the same cost: the displayed
exponent improvement can tend to zero when $r$ grows.

\section{Comparing regions and retaining information}\label{sec:cache}

A region of complete assignments is \emph{covered} when the best score
already recorded is at least the score of every assignment in that region.
Exhaustive exploration establishes coverage; a sound domination certificate
can establish it as well. In the following rule, coverage is an explicit
hypothesis about the entire stated region.

The preceding arguments motivate a broader learning rule. Fix the same
assigned variable set in two prefixes $\alpha,\beta$, and let their residual
objectives be $F_\alpha(y),F_\beta(y)$. Suppose the region $R_\alpha(y)$
under $\alpha$ has already been covered.

\begin{proposition}[Conditional comparison]\label{prop:conditional}
Within the $\beta$-branch it is sound to impose
\begin{equation}
 \neg R_\alpha(y)\ \vee\ [F_\beta(y)>F_\alpha(y)].
                                                        \label{eq:conditional}
\end{equation}
The difference $F_\alpha-F_\beta$ is affine, because their remaining
quadratic terms agree.
\end{proposition}
\begin{proof}
Every excluded completion lies in a covered region under $\alpha$ and has
score at most its counterpart there. It cannot improve the incumbent.
Restricting~\eqref{eq:quadratic} proves the affine-difference claim.
\end{proof}

For equal residual objectives up to a favorable additive constant, this
retains $R_\beta\wedge\neg R_\alpha$. If
$R_\alpha=C\wedge P$ and $R_\beta=C\wedge Q$, it simplifies to
$C\wedge Q\wedge\neg P$. For example, if $P$ is the clause
$y_1\vee y_2\vee y_3$, the second branch needs only $y_1=y_2=y_3=0$.
If $P$ is a conjunction of $s$ ordinary clauses each on at least $\ell$
distinct remaining variables, its complement contains at most
$s2^{|U|-\ell}$ assignments. Enumerating the subcubes that falsify the
clauses gives the same upper bound, allowing duplication. The cost of
covering $\alpha$ must still be included.

The negation of a conjunction is not generally an ordinary clause.
To learn a clause, find a partial suffix assignment $\tau$ for which the
current constraints certify $R_\alpha$ and $F_\alpha-F_\beta\ge0$.
Then forbid $\tau$ within $\beta$. If retained across backtracking, the
clause must include the guard specifying the $\beta$-prefix. Sufficient
certificates can use syntactic clause witnesses and the exact minimum of
an affine function over a Boolean subcube; a general implication oracle
is not assumed.

\subsection{A bounded cache for exact repeated effects}

Normalize a residual objective as $c_\alpha+G_\alpha(y)$ with
$G_\alpha(0)=0$. A completed node can store its $G_\alpha$ and constant
$c_\alpha$. If the whole subtree is covered, a later node with the same
$G$ and a no-larger constant can be discarded regardless of its inherited
auxiliary constraints. Equality of soft objectives is not an assertion
that the constrained subproblems are identical; whole-region coverage is
what makes this rule sound.

For a fixed variable order, the quadratic terms on unassigned variables at a given depth are
fixed, so a cache key needs only the residual linear coefficients. Retain
at most $n+1$ records at each depth, evicting the least recently accessed.
Only completed or already certified nodes are inserted. Eviction loses a
possible cut but cannot invalidate one.

\begin{proposition}[Bounded caching on a repeated-profile cut]
Suppose all rows of $J_{A,B}$ agree up to sign, and use the order
$V\setminus(A\cup B),A,B$. The preceding cache, together with any
coverage-preserving cuts in this fixed search, gives running time
$\Ostar(2^{n-\min(|A|,|B|)})$ and polynomial space.
\end{proposition}
\begin{proof}
For each fixed outside assignment, the residual field on $B$ depends on
$A$ only through $\sum_{u\in A}s_ux_u$, with $s_u\in\{-1,1\}$.
There are at most $|A|+1$ field vectors. At depth $|V\setminus B|$, all
records touched during this phase fit in the cache, so a record touched
in the phase cannot be evicted during it. A field vector is explored again
only when its constant strictly improves. Since
$c_\alpha=F(\alpha,0)\in\{0,\ldots,m\}$, there are at most $m+1$ such
explorations per vector. The phase costs
$\Ostar(2^{|A|}+(|A|+1)(m+1)2^{|B|})$. Multiply by the outside cost.
\end{proof}

The explicit unweighted input is relevant to the polynomial bound on record
improvements. Binary-encoded arbitrary weights require a different argument.
This proposition is a structured application of caching, not a claim that
a polynomial cache always finds all useful similar states.

\section{Toward one branching algorithm}\label{sec:unified}

A concrete candidate uses the same search operations throughout:
\begin{enumerate}
\item Maintain the residual objective, the scoped comparison constraints,
and a bounded cache of covered residuals.
\item Reject a node if its constraints have an infeasible linear relaxation,
or a cached covered residual dominates it.
\item Select the next variable in a prescribed order. Try first the value
preferred when remaining variables are set to $1/2$.
\item After that child is covered, impose its affine comparison in the other
child and recurse. Record coverage on returning.
\end{enumerate}
The linear relaxation is an explicit polynomial-time test, not an oracle
for Boolean satisfiability. Single-variable comparisons and cache cuts are
sound together by induction on the order in which regions become covered.

\begin{proposition}[Why no enumeration switch is needed for the basic bounds]
For a fixed order, this engine inherits the bounds of
Theorems~\ref{thm:bias} and~\ref{thm:weighted}, the sharper
Corollary~\ref{thm:random}, and the unbiased cases of
Corollary~\ref{thm:randomclauses}. With the indicated repeated-profile order,
it also inherits the preceding caching bound.
\end{proposition}
\begin{proof}
The bias analysis applies with no forced steps when only binary branching
and the sum test are used. Linear infeasibility detects every rejection
made by that sum test, so the same time bound holds for this engine. For the random-sign argument, fix any
prefix at the analysis depth $T$, and freeze the sum of constraints already
created there. Every later node retained by the relaxation has a fractional
completion satisfying this one affine inequality. The minimum of one affine
function over a Boolean subcube is attained at a Boolean vertex, so this
node also has a Boolean completion satisfying the frozen sum. If that sum
has $N$ completions, there are at most $N$ retained descendants at any suffix
depth. Thus the cost is bounded by~\eqref{eq:enumerationcost}, even though
the algorithm continues learning and never switches modes.
Positive weights cause no difficulty: the linear relaxation implies
every positive weighted sum of its comparisons, including that of
Theorem~\ref{thm:weighted}. For the unbiased models of
Corollary~\ref{thm:randomclauses}, take $P=0$.
The preferred values then agree with those of this engine, and each strict
second-child comparison implies the corresponding weak comparison used
in~\eqref{eq:counts-sum}. The same frozen-sum argument therefore bounds
its descendants by the surviving mass in Theorem~\ref{res:counts}.
The cache only removes nodes of this fixed-order, fixed-preference search. The structured
bound was proved directly above.
\end{proof}

\paragraph{Ordering still has a cost.}
The repeated-profile bound is not asserted for every order. For a fixed
$0<\delta\le1/8$, one explicit policy enumerates $k=\lfloor\delta n\rfloor$
element sets $B$, groups outside rows on $B$ up to sign, and searches for a
group $A$ of size $k$. A successful test supplies an order with the
structured bound; if none succeeds, use a fixed default order.
The search costs $\Ostar(\binom nk)$. On every draw from a covered unbiased
random model, either the
test succeeds and gives its unconditional structured bound, or the default
order is used and the exceptional event is contained in the exceptional
event for that fixed order. Bias is order-independent.
The resulting bound includes the detection cost, the possible successful
case $\Ostar(2^{(1-\delta)n})$, and the default-order bound. Thus it retains
a fixed saving, though not necessarily the same numerical saving as the
default-order theorem. For a promised biclique family, choose $\delta$
below its guaranteed side fraction (or the extracted core fraction).
This is a uniform solver for literal bias, the unbiased random models,
and the indicated structured cuts, with an explicit ordering search.
A purely local adaptive ordering rule with the same guarantees remains open.

The count guesses used for biased random formulas and regular-graph
Max-Cut change the gains used to generate comparisons. Hidden communities
use the same principle with computed projection directions. The round
version also fixes preferences earlier. These operations are
handled by the single generator in Theorem~\ref{res:counts}, with fresh
constraints for every run. One can interleave this generator (using the
explicit constructions for the desired classes) and the preceding engine,
stopping when an exact solver finishes. A fixed number of polynomial-space
solvers costs only a constant factor over the fastest one, up to simulation
overhead. This gives a uniform solver for the combined class guarantees.
It does not require transferring cache records between different auxiliary
objectives. A local adaptive policy achieving all these guarantees without
this interleaving remains open.
Nor do we claim that arbitrary
pre-branching clauses can be added without reconciling their representative
order with its coverage invariant. These are separate integration tasks.

\section{Open directions: exploiting redundancy}
\label{sec:open}

\subsection{A structural explanation of persistent correlation}

The most immediate target is to weaken the cancellation assumptions of the
balanced class. The bias proof already permits arbitrary correlations.
In the balanced argument, coherent comparison coefficients can make the
remaining uncertainty large. Some coherent families are resolved by group
counts or by the bounded-block theorem. This suggests analyzing the search
only after applying those available comparisons.

\begin{question}[From poor cancellation to exploitable comparisons]
When a large set of generated comparisons fails to accumulate useful
pruning, can one efficiently find either a family of bounded block
replacements with few exceptions, or a larger group of assignments whose
outside effects admit affordable compression?
\end{question}

Absence of an exact short relation gives only a nonzero signed sum of
neighborhood vectors. It does not give the quantitative cancellation bound
needed in the present proofs. Theorem~\ref{thm:blocks} handles a robust
version with a constant number of exceptions. Extending this to a growing
number of exceptions without losing a fixed exponent saving is a concrete
intermediate problem. It may require using the learned conditional
comparisons rather than assigning every exception independently.

\subsection{Similarity of regions, not just equality of formulas}

\begin{question}[Affordable regional coverage]
Can the algorithm maintain polynomially many conditional comparison
certificates that cover most repeated optimization behavior, even when no
two residual objectives are identical?
\end{question}

At a fixed depth, residual linear vectors provide a direct representation
for comparisons. Differences supported on a small set give local tests;
common learned clauses cancel from the regional comparison. The relevant
cost is finding and retaining the certificates as well as exploiting them.
The affine-difference observation avoids general formula isomorphism, but
does not make arbitrary Boolean implication tests polynomial-time.

\subsection{An amortized measure for overlapping constraints}

Deleting a satisfied clause and shortening an unsatisfied clause have
different effects on the surviving completion space. A useful measure should
account for both, and should charge only genuinely new restrictions.
For a constraint family $C$ on $U$, the ideal information measure is
\[
 |U|-\log_2\#\{y\in\bits^U:y\models C\}.
\]
It is not generally efficiently computable and does not by itself pay for
enumeration. A tractable bound based on sums, selected subfamilies, or
conditional coverage could support an adaptive branching rule. The measure
should also charge the discovery and storage of comparisons.

\subsection{Randomized exploration with exact answers}

Las Vegas randomization may help choose comparison blocks or enumerate
survivors. It must preserve the exact optimum on every run. If a capped
attempt costs at most $B$ and succeeds with probability at least $p$,
independent restarts cost at most $B/p$ in expectation. A small number of
satisfying full assignments does not supply this success probability by
itself. A useful result must bound the work spent on partial assignments
without a surviving completion.

\subsection{The long-term target}

\begin{target}[Growing density with a uniform saving]\label{target:main}
Find a deterministic polynomial-space exact algorithm and a constant
$\gamma>0$ for all canonical unweighted Max 2-SAT formulas without repeated
binary clauses, allowing complementary pairs, running in
$\Ostar(2^{(1-\gamma)n})$ time.
An intermediate target is the same guarantee for every formula up to a
specified density bound $\rho(n)\to\infty$, with $\gamma$ independent of $n$.
\end{target}

This is a desired result, not a consequence or a conjecture established by
the present arguments. The current paper's contribution should be assessed
through its explicit classes. A structural dichotomy is valuable if both
outcomes are algorithmic and their costs combine into a fixed saving; merely
showing that some useful relation exists is insufficient.

\appendix
\section{A sharper constant for almost-regular random signs}\label{app:regular}

Fix a simple graph $G$ and a vertex order before drawing signs. Each edge
contributes one clause, and its two literal signs are independent and uniform;
all incidence signs are mutually independent. There are no initial unit
clauses. The graph itself is arbitrary. Write
$d=\min_v\deg(v)$ and $D=\max_v\deg(v)$.

\begin{corollary}[A fixed saving for random signs]\label{thm:random}
For $d\ge6$, there is a deterministic polynomial-space algorithm that is
exact on every signing and, with probability at least
$1-3e^{-nd/(1250D)}$, takes time
\[
 \Ostar\left((2e^{-d/(1250D)})^n\right).
\]
Its expected running time over the signs is
$\Ostar((2e^{-d/(625D)})^n)$.
\end{corollary}

In particular, every regular graph gives the base
$2e^{-1/1250}\simeq1.99840064$, independently of its degree. This includes
any sequence of regular graphs of degree $\omega(1)$, and the complete graph.
Randomness belongs to the instance distribution; the solver neither guesses
an answer nor has a probability of returning an incorrect optimum.

\begin{proof}
Learn~\eqref{eq:cardcut} in the fixed order through depth
$T=\lceil n/2\rceil$, and then enumerate the sum as in
Lemma~\ref{lem:enumeration}. Let $x$ be a uniform complete assignment,
independent of the signs, and retain even constraints of already cut prefixes
for this upper-bound analysis. Define
\[
 S=\sum_{i\text{ creator},\ i\le T}(|g_i|/2+1),\qquad
 Q=\sum_{i\text{ creator},\ i\le T}(L_i(x)-D_i^+/2),
\]
where $D_i^+$ counts later neighbors. An assignment passes its prefix's
summed inequality exactly when $S+Q\le0$. Therefore
\begin{equation}
 \E_F W(F)=2^n\Prb_{F,x}(S+Q\le0).              \label{eq:randommass}
\end{equation}

Reveal the incidence signs of vertices in order, together with their assigned
values. Just before vertex $i$, let $r_i$ count earlier neighbors whose
partner literal in their clause with $i$ is false. Conditional on the history,
$g_i$ is a sum of independent signs of weights one and two, of variance
$Y_i=D_i^++4r_i$. The fourth-moment bound for a Rademacher sum gives
$\E|g_i|\ge\sqrt{Y_i/3}$. The contribution $X_i$ to $S$ thus has conditional
mean at least
$\mu_i=\sqrt{Y_i}/(4\sqrt3)+1/2$ and second moment at most $D$, for $D\ge6$.
The same nonnegative-increment estimate as above gives
\begin{equation}
 \Prb\left(S<\sum_{i\le T}\mu_i-a\right)
 \le e^{-a^2/(2TD)}.                            \label{eq:randomStail}
\end{equation}

The earlier endpoint-falsity indicators are independent fair coins, including
after averaging over $x$. Hence $r_i$ has distribution
$\operatorname{Bin}(D_i^-,1/2)$, with
$\E Y_i=D_i^++2D_i^-\ge d$ and $\operatorname{Var}(Y_i)\le2\E Y_i$.
For nonnegative $Y$ of mean $\mu$, the pointwise inequality
$\sqrt z\ge1+(z-1)/2-(z-1)^2/2$ yields
$\E\sqrt Y\ge\sqrt\mu-\operatorname{Var}(Y)/(2\mu^{3/2})$.
It follows that
$\E\sum_i\mu_i\ge T\sqrt d/(4\sqrt3)$.
Changing one earlier-endpoint coin changes this sum by at most
$1/(2\sqrt3)$. There are at most $TD/2$ relevant internal edges, so
bounded differences gives
\begin{equation}
 \Prb\left(\sum_i\mu_i<T\sqrt d/(4\sqrt3)-b\right)
 \le e^{-48b^2/(TD)}.                           \label{eq:randommutail}
\end{equation}

Finally reveal all later endpoints, without creating constraints after $T$.
For each newly revealed vertex, the bad-value indicators from earlier
creators use distinct unseen incidence signs. Conditional on the history,
they are independent fair bits. There are at most $TD$ relevant edge
contributions in total. The conditional $\cosh$ bound therefore gives
\begin{equation}
 \Prb(Q<-c)\le e^{-2c^2/(TD)}.                  \label{eq:randomQtail}
\end{equation}
No independence between $S$ and $Q$ is needed.

Take $(a,b,c)=(2/25,1/100,1/20)T\sqrt d$. Their sum is less than
$T\sqrt d/(4\sqrt3)$. Combining the three tails with~\eqref{eq:randommass}
gives $\E_F W(F)\le3\cdot2^n e^{-2Td/(625D)}$.
Markov's inequality implies
\[
 \Prb_F\left(W(F)>2^n e^{-Td/(625D)}\right)
 \le3e^{-Td/(625D)}.
\]
Use $T=\lceil n/2\rceil$ and~\eqref{eq:enumerationcost}. The prefix cost
$2^T$ is smaller than the asserted bounds, up to constant factors.
\end{proof}

The argument counts total remaining work, not merely the fraction of prefixes
with a conflict. A satisfiable summed constraint can already justify a large
saving. Adaptive variable selection based on unrevealed signs is not covered
by this exposure proof.

\section{Independent-pair membership}\label{app:membership}

For this appendix, $P_{uv}=\sum_k\beta_kt_k(u)t_k(v)$ with ternary
$t_k$, $N=\operatorname{offdiag}(J-P)$, $\sigma=\sum_k|\beta_k|$,
$T=\lfloor\tau n\rfloor$, $R_{tw}=N_{tw}[w>t]$ for $t\le T$, and
$W_t=\sum_{j<t}N_{tj}^2$. Exact count grids have $s_t\le\sigma$.
Thus the following imply the sequential case of Theorem~\ref{res:counts}:
\begin{equation}
 \|R\|^2\le KD_N,\qquad
 \sum_{t\le T}\sqrt{W_t}\ge4\sqrt3\,gn\sqrt{D_N},\qquad
 \sigma\le g\sqrt{D_N}/(8\tau).\label{eq:remainder}
\end{equation}

The following is the independent-pair criterion of Proposition~12 in
\cite{NoteCounts}. We write $\chi$ for the variance ratio to distinguish
it from clause density. The probability event concerns $J$ alone, which is
why the linear coefficients and unit clauses can be arbitrary.

\begin{proposition}[Independent-pair membership]\label{prop:pairmembership}
Fix $0<\tau<1$, $\chi\ge1$, $C_0=10^9\chi/\tau$, and
$n\ge10^9\chi^2/\tau^2$. Suppose the coefficients $J_{uv}$, $u<v$, are
independent, with
\[
 |J_{uv}-\E J_{uv}|\le B,\qquad
 0<\nu\le\operatorname{Var}J_{uv}\le\chi\nu.
\]
Choose either of the following decompositions:
\begin{enumerate}
\item All means equal $\mu$, and $A_{uv}=\beta$, where $\beta$ is the
empirical mean over pairs. Put $\sigma_0=|\mu|$.
\item A fixed supplied aligned decomposition has $A_{uv}=\E J_{uv}$.
Put $\sigma_0=\sum_k|\beta_k|$. In particular, if all means vanish,
take $A=0$ and $\sigma_0=0$.
\end{enumerate}
If $n\nu\ge C_0(B^2\log n+\chi\sigma_0^2)$, then with probability at least
$1-1/n$, conditions~\eqref{eq:remainder} hold with
$K=9\chi$ and $g=\tau^{3/2}/(12\sqrt\chi)$.
The lower bound on $n$ in Theorem~\ref{res:counts} also holds.
\end{proposition}

\begin{proof}
Put $X_{uv}=J_{uv}-\E J_{uv}$, and let $\delta=\beta-\mu$ in the first
case and $\delta=0$ in the second. Then $N_{uv}=X_{uv}-\delta$ off the
diagonal. Set $\eta=1/100$.

We use scalar Bernstein concentration~\cite{BLM13}, Talagrand's convex
concentration inequality~\cite{Tal95}, and the rectangular norm estimate
of Bandeira and van Handel~\cite[Theorem~3.1 and Corollary~3.2]{BvH16}.
The last estimate, for independent random signs $\xi_{tw}$ and a fixed
coefficient array $a$, is
\[
 \E\|(\xi_{tw}a_{tw})\|
 \le(1+\vartheta)\left(\sigma_1+\sigma_2+
       \frac{5\max_{t,w}|a_{tw}|\sqrt{\log n}}
            {\sqrt{\log(1+\vartheta)}}\right),\quad
 0<\vartheta\le1/2,
\]
where $\sigma_1,\sigma_2$ are the largest row and column Euclidean norms.
For a convex $1$-Lipschitz function of independent coordinates in $[-B,B]$,
we use the tail
$\Prb(f>\E f+8\sqrt\pi B+r)\le4e^{-r^2/(16B^2)}$.

\paragraph{Centering and total variance.}
Bernstein's inequality gives, in the empirical case,
$|\delta|\le\delta_0=6\sqrt{\chi\nu\log n}/n$ except with probability
at most $n^{-3}$. For example, applying it to
$\sum_{u<v}X_{uv}$ gives a two-sided exponent at least
$4.5(1-1/n)\log n$. The bound on $n$ implies
$\delta_0\le10^{-3}\sqrt\nu$ and
$n\delta_0\le10^{-3}\sqrt{n\chi\nu}$.
Bernstein applied to $\sum X_{uv}^2$ also gives
\[
 (1-\eta)n(n-1)\nu\le\sum_{u\ne v}X_{uv}^2
           \le(1+\eta)n(n-1)\chi\nu.
\]
In the empirical case,
$D_N=n^{-1}\sum_{u\ne v}X_{uv}^2-(n-1)\delta^2$; in the expected case,
$D_N=n^{-1}\sum_{u\ne v}X_{uv}^2$. Hence, outside another event of
probability at most $n^{-3}$,
\begin{equation}
 (1-2\eta)(n-1)\nu\le D_N\le(1+\eta)(n-1)\chi\nu. \label{eq:random-D}
\end{equation}

\paragraph{The operator norm.}
Let $R_X$ have entries $X_{tw}$ for $t\le T$, $w>t$, and zero otherwise.
The entries in this array are independent. Symmetrize with an independent
copy $X'$ and put $\widehat X=X-X'$. Then
$\E\|R_X\|\le\E\|R_{\widehat X}\|$.
Condition on $\widehat X$ and introduce independent signs. Since its entries
are symmetric, the displayed norm estimate with $\vartheta=\eta$ gives
\[
 \E\|R_X\|\le
 (1+\eta)(\E\sigma_1(\widehat X)+\E\sigma_2(\widehat X))
       +102B\sqrt{\log n}.
\]
Each squared entry is at most $4B^2$ with expectation at most $2\chi\nu$.
Bernstein and a union bound show that every row-square sum is at most
$2(1+\eta)n\chi\nu$, and every column-square sum at most
$2(\tau+\eta)n\chi\nu$, except with probability
$2n e^{-\eta^2n\chi\nu/(5B^2)}\le n^{-3}$.
Bounding each exceptional sum by $4nB^2$ gives
\[
 \E\sigma_1(\widehat X)^2\le2(1+2\eta)n\chi\nu,\qquad
 \E\sigma_2(\widehat X)^2\le2(\tau+2\eta)n\chi\nu.
\]
Thus $\E\|R_X\|\le2.89\sqrt{n\chi\nu}$, using
$102/\sqrt{C_0}\le0.004$. Convex concentration with
$r=\eta\sqrt{n\chi\nu}$, together with
$8\sqrt\pi B\le10^{-3}\sqrt{n\chi\nu}$, now gives
\[
 \|R\|\le\|R_X\|+n|\delta|\le2.902\sqrt{n\chi\nu},
 \qquad \|R\|^2\le8.43n\chi\nu\le9\chi D_N
\]
outside an event of probability at most $n^{-3}$.

\paragraph{The earlier-neighbor mass.}
Put $t_0=\lceil\tau n/10\rceil$. For $t>t_0$, scalar Bernstein gives
$V_t=\sum_{j<t}X_{tj}^2\ge(1-\eta)(t-1)\nu$ with failure probability at most
$e^{-\eta^2\tau n\nu/(30B^2\chi)}\le n^{-3}$.
Consequently $\sqrt{W_t}\ge\sqrt{V_t}-|\delta|\sqrt{t-1}
\ge0.993\sqrt{(t-1)\nu}$. Summing and using $T\ge\tau n-1$ yields
\[
 \sum_{t\le T}\sqrt{W_t}
 \ge0.635\,\tau^{3/2}n^{3/2}\sqrt\nu
 \ge4\sqrt3\,g n\sqrt{D_N};
\]
the last expression is at most
$0.581\,\tau^{3/2}n^{3/2}\sqrt\nu$ by~\eqref{eq:random-D}.

\paragraph{The allowed comparison slack.}
Equation~\eqref{eq:random-D} gives
$g\sqrt{D_N}/(8\tau)\ge0.0102\sqrt{\tau n\nu/\chi}$.
In the expected case,
$\sigma=\sigma_0\le\sqrt{n\nu/(C_0\chi)}
\le10^{-4}\sqrt{\tau n\nu/\chi}$.
In the empirical case,
$\sigma=|\beta|\le\sigma_0+\delta_0
\le0.002\sqrt{\tau n\nu/\chi}$.
This establishes the third condition.

There are at most $n+4$ failure events, each of probability at most
$n^{-3}$, so their union has probability at most $1/n$.
Finally
$256\pi\tau K/g^2=256\pi\cdot1296\,\chi^2/\tau^2
<1.1\cdot10^6\chi^2/\tau^2\le n$.
\end{proof}

\section{Analysis of accumulated comparisons}\label{sec:counts}\label{app:accumulation}

This appendix supplies the common analysis for the random-formula and
Max-Cut families in Sections~\ref{sec:random} and~\ref{sec:maxcut}.
It also specifies the projection comparison routine used for hidden
communities. A few aggregate guesses account for systematic interactions;
the remaining comparisons accumulate enough gap while a norm bound
controls their combined fluctuation. The matrix formulation is proof
machinery for these input families, rather than an additional class
definition required of the reader.

\subsection{One comparison algorithm}

Choose a rational representation
$P=\sum_{k=1}^r\lambda_kv_kv_k^{\mathsf T}$ of polynomial encoding length,
and put $N=\operatorname{offdiag}(J-P)$. For each $k$, a finite rational
grid $\mathcal Y_k$ must approximate every $\langle v_k,x\rangle$ to
within a stated rational error $\delta_k\ge0$. Define
\[
 s_t=\sum_k|\lambda_k|\,|v_k(t)|(\delta_k+|v_k(t)|),
 \qquad \mathcal Y=\prod_k\mathcal Y_k.
\]
The grid is streamed; its description and each point have polynomial size.
In a run indexed by $\hat y\in\mathcal Y$, use
\[
 a^{\hat y}=a+\sum_k\lambda_k\hat y_kv_k,
 \qquad F'_{\hat y}(x)=c+\langle a^{\hat y},x\rangle+
                    \sum_{u<v}N_{uv}x_ux_v.
\]
Partition a prefix of $T=\lfloor\tau n\rfloor$ variables into consecutive
rounds $B_1,\ldots,B_\ell$; the suffix is round $\ell+1$.
Let $\rho(t)$ be the round of $t$. Throughout round $j$, prefer
$p_t=[\Delta_t\ge0]$, where
\[
 \Delta_t=a^{\hat y}_t+\sum_{w:\rho(w)<j}N_{tw}x_w
             +\frac12\sum_{w\ne t:\rho(w)\ge j}N_{tw}.
\]
The preference uses only earlier rounds. Call $t$ a creator when
$x_t\ne p_t$, and let $\gamma_t(x)$ be the gain under $F'_{\hat y}$ of
flipping it to $p_t$. At each depth-$T$ prefix enumerate the suffixes satisfying
\begin{equation}
 \sum_{t\text{ creator}}(\gamma_t(x)-s_t)\le0. \label{eq:counts-sum}
\end{equation}
This is one affine inequality on the suffix, so Lemma~\ref{lem:enumeration}
applies. Evaluate the original objective on every candidate. Candidates
need not have projections near the guessed grid point.

\begin{lemma}[Sound comparisons with approximate projections]\label{lem:counts-exact}
Every optimum is enumerated in a run that approximates its projections.
The algorithm is exact for every supplied representation and round partition.
\end{lemma}
\begin{proof}
Write $y_k=\langle v_k,x\rangle$. The difference between the original
$0\to1$ gain at $t$ and the auxiliary gain is
\[
 \sum_k\lambda_kv_k(t)(y_k-\hat y_k-v_k(t)x_t).
\]
If $|y_k-\hat y_k|\le\delta_k$, its absolute value is at most $s_t$.
Every flip at an optimum has original gain at most zero, so each creator
has $\gamma_t\le s_t$ in an approximating run. Their sum satisfies
\eqref{eq:counts-sum}. Weak inequalities retain ties.
\end{proof}

The accumulation proof below uses only nonnegative comparison allowances,
their total size, and the fact that an optimum survives in at least one
run. Thus a smaller allowance may replace $s_t$ if its validity is proved
directly. The same proof allows a joint finite list of aggregate guesses,
rather than a Cartesian product of independent grids. This includes the
supplied-community counts and diagonal allowances in
Section~\ref{sec:sbm}; the running bound uses the number of joint guesses.

\subsection{A general accumulation theorem}

Let $D_N=n^{-1}\sum_{u\ne v}N_{uv}^2>0$ and
$W_t=\sum_{w:\rho(w)<\rho(t)}N_{tw}^2$. The shared assumptions are
\begin{equation}
 \sum_{t\le T}\sqrt{W_t}\ge4\sqrt3\,g n\sqrt{D_N},\qquad
 \sum_{t\le T}s_t\le\tfrac18g n\sqrt{D_N}.       \label{eq:profile-slack}
\end{equation}
All constants below are independent of $n$.

\begin{theorem}[Accumulation after projection guessing]\label{res:counts}
Assume~\eqref{eq:profile-slack}, with $0<\tau<1$ and $g,K>0$.
Either of the following conditions gives a deterministic exact algorithm
in polynomial space and time
\[
 \Ostar\bigl(|\mathcal Y|(2^{\tau n}+2^ne^{-c_*n})\bigr).
\]
\begin{enumerate}
\item \textbf{Sequential comparisons.} Each prefix round is a singleton.
Let $R_{tw}=N_{tw}[w>t]$ for $t\le T$, and assume $\|R\|^2\le KD_N$.
Set $K_L=K$ and
$c_*=g^2/(64\tau\alpha^2K)$.
\item \textbf{A fixed number of rounds.} There are $\ell\ge2$ prefix
rounds, with $\ell$ fixed, and $\|N\|^2\le KD_N$.
Set $K_L=(\ell-1)K$ and
\[
 c_*=\min\left\{
 \frac{g^2}{32\tau\alpha^2K_L},\,
 \frac{g^2}{512\ell^2\tau K},\,
 \frac{c_{\rm HW}g^2}{256\ell^2\tau K},\,
 \frac{c_{\rm HW}g}{16\ell\sqrt K}\right\}.
\]
Here $c_{\rm HW}>0$ is an absolute Hanson--Wright constant.
\end{enumerate}
In both cases, $\alpha=\max(2,9\tau\sqrt{K_L}/(4g))$ and
$n\ge256\pi\tau K_L/g^2$. Linear coefficients and units are arbitrary.
For polynomially many grid points this is a fixed base below two.
\end{theorem}

Exact count guessing is a special case: for
$v_k\in\{-1,0,1\}^n$, use all integer counts and $\delta_k=0$.
Then $|\mathcal Y|\le(2n+1)^r$ and $s_t\le\sum_k|\lambda_k|$.
Empirical mean centering takes $r=1$, $v_1=\ind$,
$\lambda_1=\binom n2^{-1}\sum_{u<v}J_{uv}$, and only $n+1$ counts.
If $P=0$, there is one run and no comparison slack.

\begin{lemma}[A common gap estimate]\label{lem:counts-profile}
Let $L$ be the $T\times T$ matrix with $L_{tw}=N_{tw}$ when
$\rho(w)<\rho(t)$ and zero otherwise. Suppose $\|L\|^2\le K_LD_N$ and
\eqref{eq:profile-slack}'s gap bound holds. With
$M=\alpha\sqrt{K_LD_N}$ and the theorem's lower bound on $n$,
\[
 \Prb\left(\sum_{t\le T}\min\{|\Delta_t|,M\}
               <\tfrac34gn\sqrt{D_N}\right)
 \le4e^{-g^2n/(16\tau K_L)}.
\]
The probability is over a uniform Boolean assignment, with the formula fixed.
\end{lemma}
\begin{proof}
Put $u=2x-\ind$ and $H_t=a^{\hat y}_t+\frac12\sum_{w\ne t}N_{tw}$.
Write $Lu$ for $L(u_1,\ldots,u_T)^{\mathsf T}$, so
$\Delta_t=H_t+\frac12(Lu)_t$.
Partition the indices into $B_0=\{t:|H_t|\le M/2\}$ and its complement
$B_1$. Since $\|Lu\|^2\le K_LD_NT$, clipping the gaps at $M$ loses at most
$K_LD_NT/(8M)$ from $\sum_{B_0}|\Delta_t|$; use
$(a-b)_+\le a^2/(4b)$ with $a=|(Lu)_t|/2$, $b=M/2$.
For $t\in B_1$, the clipped gap is at least $M/4$ whenever
$|(Lu)_t|\le M/2$, and at most $4T/\alpha^2$ indices fail this condition.
Thus, for $G(u)=\sum_{B_0}|\Delta_t|$,
\[
 \sum_{t\le T}\min\{|\Delta_t|,M\}
 \ge G(u)+\frac M4|B_1|-\frac{9T\sqrt{K_LD_N}}{8\alpha}.
\]
A symmetric Rademacher sum of variance $W_t$ has first absolute moment
at least $\sqrt{W_t/3}$. Hence
$\E|\Delta_t|\ge\frac12\sqrt{W_t/3}$, uniformly in $H_t$.
Also $W_t\le K_LD_N$ and $M/4\ge\frac12\sqrt{W_t/3}$. Consequently
$\E G+M|B_1|/4\ge2gn\sqrt{D_N}$, and the displayed loss is at most
$gn\sqrt{D_N}/2$.
The convex function $G$ is Lipschitz with constant at most
$h=\frac12\sqrt{TK_LD_N}$. Talagrand's inequality~\cite{Tal95} gives
\[
 \Prb(G\le\E G-a-8\sqrt\pi h)\le4e^{-a^2/(16h^2)}.
\]
The mean-to-median shift follows by integrating the median tail.
The bound on $n$ makes $8\sqrt\pi h\le gn\sqrt{D_N}/4$;
set $a=gn\sqrt{D_N}/2$.
\end{proof}

\begin{proof}[Proof of Theorem~\ref{res:counts}]
Correctness is Lemma~\ref{lem:counts-exact}. Fix a run and a uniform $x$,
and write $u=2x-\ind$ and $I$ for its creators. In the sequential case,
$L$ is the transpose of the leading square block of $R$, so $\|L\|\le\|R\|$.
In the round case, each row block is $N_{B_j\times B_{<j}}$, whence
$\|L\|^2\le(\ell-1)\|N\|^2$. The common gap estimate applies in both cases.

\paragraph{Accumulated gap.}
Given the earlier variables, each creator indicator is a fair coin and
$\Delta_t$ is fixed; within a round all preferences were fixed before it
started. For $Y_t=[t\in I]\min(|\Delta_t|,M)$, Azuma--Hoeffding gives
\[
 \Prb\left(\sum_tY_t<\tfrac12\sum_{t\le T}\min(|\Delta_t|,M)-a\right)
 \le e^{-2a^2/(TM^2)}.
\]
Use $a=gn\sqrt{D_N}/8$. The gap lemma and the slack bound imply
$S'=\sum_{t\in I}(|\Delta_t|-s_t)\ge gn\sqrt{D_N}/8$
outside these two exceptional events.

\paragraph{Sequential comparisons.}
Here the sum tested by the algorithm is
\begin{equation}
 S'+Q,\qquad Q=-\tfrac12\sum_wu_w\sum_{t\in I,t<w}u_tN_{tw}.
                                                        \label{eq:counts-race}
\end{equation}
The inner coefficient is predictable, and its squared sum is at most
$\|R^{\mathsf T}(u_t[t\in I])_{t\le T}\|^2\le KD_NT$.
The conditional $\cosh$ bound gives
$\Prb(Q\le-b)\le e^{-2b^2/(KD_NT)}$.
For $b=gn\sqrt{D_N}/16$, a passing assignment belongs to an exceptional
event. Their total probability is at most $6e^{-c_*n}$ with the stated
sequential constant (which conservatively retains the earlier bound).

\paragraph{Comparisons within rounds.}
The tested sum is $S'+Q_{\rm later}+\sum_{j\le\ell}Q_j$, where
\[
 Q_{\rm later}=-\tfrac12\sum_wu_w
        \sum_{t\in I:\rho(t)<\rho(w)}u_tN_{tw},\qquad
 Q_j=-\tfrac12\sum_{t\in I\cap B_j}\sum_{w\in B_j\setminus\{t\}}u_tN_{tw}u_w.
\]
For a later round the coefficients are fixed from earlier rounds.
Each block of coefficients has squared norm at most $KD_NT$, and there
are at most $\ell$ such blocks, including the suffix. Thus
\[
 \Prb(Q_{\rm later}\le-b)\le e^{-2b^2/(\ell KD_NT)}.
\]
Take $b=gn\sqrt{D_N}/32$.

Condition on the rounds preceding $B_j$. Put $c_t=2p_t-1$ and
$\beta_t=-c_tu_t$. These are independent fair signs, and
$t\in I$ precisely when $\beta_t=1$. The matrix
$M^{(j)}_{tw}=c_tc_wN_{tw}$ on $B_j$ has zero diagonal and norm at most
$\sqrt{KD_N}$. Direct expansion gives
\[
 Q_j=-\tfrac14\beta^{\mathsf T}M^{(j)}\beta
          -\tfrac14\langle M^{(j)}\ind,\beta\rangle.
\]
Both $\|M^{(j)}\ind\|^2$ and $\|M^{(j)}\|_F^2$ are at most $|B_j|KD_N$.
For $e=gn\sqrt{D_N}/(64\ell)$, the linear term is below $-e$ with
probability at most $e^{-8e^2/(|B_j|KD_N)}$. Hanson--Wright~\cite{RV13}
bounds the corresponding quadratic probability by
\[
 2\exp\left[-c_{\rm HW}\min\left\{
       \frac{16e^2}{|B_j|KD_N},\frac{4e}{\sqrt{KD_N}}\right\}\right].
\]
A union bound therefore gives $\sum_jQ_j\ge-gn\sqrt{D_N}/32$ outside
these events. On the remaining assignments the tested sum is at least
$(1/8-1/32-1/32)gn\sqrt{D_N}>0$.
Using $|B_j|\le T\le\tau n$, the total passing probability is at most
$(6+3\ell)e^{-c_*n}$. Lemma~\ref{lem:enumeration} finishes both time
bounds. Runs and candidates are streamed, so space is polynomial.
\end{proof}

\subsection{Projection grids and discovering the low-rank part}

\begin{corollary}[Polynomial projection grids]\label{cor:spectral}
Suppose $\|v_k\|_2\le1$ and use
$\mathcal Y_k=\delta\mathbb Z\cap[-\|v_k\|_1-\delta,\|v_k\|_1+\delta]$,
where $\delta>0$ is rational. Then
\[
 |\mathcal Y|\le(2\sqrt n/\delta+3)^r,\qquad
 \sum_{t\le T}s_t\le\sum_k|\lambda_k|(\delta\sqrt n+1).
\]
Thus Theorem~\ref{res:counts} applies whenever its norm and gap conditions
hold and the last expression is at most $gn\sqrt{D_N}/8$.
\end{corollary}
\begin{proof}
Cauchy--Schwarz gives $\|v_k\|_1\le\sqrt n$. Also
$\sum_ts_t=\sum_k|\lambda_k|(\delta\|v_k\|_1+\|v_k\|_2^2)$.
\end{proof}

The vectors need not be exact eigenvectors. One deterministic construction
uses a rational, approximately diagonalizing orthogonal matrix, takes its
$r$ largest diagonal entries in absolute value, and uses the corresponding
columns. Appendix~\ref{app:discovery} gives a polynomial-bit construction
and a random-model guarantee. All pruning comparisons use the rational
vectors actually computed, so numerical approximation affects the running
bound, not exactness.

\section{Spectral preprocessing for hidden communities}\label{app:discovery}

\subsection{A discovery lemma for the stochastic block model}

\begin{lemma}[Discovering a bounded-rank mean]\label{cor:planted}
Fix a rank bound $r$ and variance ratio $\chi\ge1$. Suppose the integer
coefficients $J_{uv}$ are independent over $u<v$,
\[
 |J_{uv}-\E J_{uv}|\le B,\quad
 0<\nu\le\operatorname{Var}J_{uv}\le\chi\nu,
\]
where we take $B\ge1$. Suppose their means are the off-diagonal entries
of a symmetric matrix $\bar A$ with
$\rank\bar A\le r$, $|\bar A_{uu}|\le B_0$ and $\|\bar A\|\le\mu n$.
For a sufficiently large $C=C(r,\chi)$, if
\[
 n\nu\ge C(B^2\log n+B_0^2+\mu^2),
\]
there is a deterministic exact polynomial-space algorithm, given the fixed
class parameters $r,\chi$ but not $\bar A$ or $\nu,B,B_0,\mu$, taking
$\Ostar(2^{(1-\gamma)n})$ time with probability $1-o(1)$, for a fixed
$\gamma>0$. Arbitrary units, even chosen after the binary draw, are allowed.
\end{lemma}


We give the computational detail separately so that the general theorem
requires only rational comparisons, with no exact-real arithmetic oracle.

\paragraph{Rational approximate diagonalization.}
For a rational symmetric matrix $J$ and rational $\zeta>0$, one can compute
a rational orthogonal $Q$ and rational diagonal $D$ with
$\|J-QDQ^{\mathsf T}\|_F\le\zeta$ in time polynomial in the input bit
length, $n$, and $\log(1/\zeta)$. Here a slow Jacobi implementation suffices.
For completeness, let $A=Q^{\mathsf T}JQ$ and let
$E=\sum_{i\ne j}A_{ij}^2$. An exact plane rotation eliminating a largest
entry $A_{ij}$ reduces $E$ by $2A_{ij}^2\ge2E/[n(n-1)]$.
Approximate that rotation by a rational orthogonal plane rotation using
\[
 \cos\theta=\frac{1-t^2}{1+t^2},\qquad
 \sin\theta=\frac{2t}{1+t^2},\qquad t\in\mathbb Q.
\]
The eliminating angle can be chosen in $[-\pi/4,\pi/4]$, so its parameter
$t=\tan(\theta/2)$ is approximated on a bounded interval by elementary
root approximation. With rotation error
$O(\zeta^2/(n^2(1+\|J\|_F)^2))$, the error in the decrease of $E$ is at
most $\zeta^2/[2n(n-1)]$. Therefore while $E>\zeta^2$, it contracts by
at least a fixed multiple of $1/n^2$ per step. After
$O(n^2\log(2+\|J\|_F/\zeta))$ steps, set $D=\operatorname{diag}(A)$.
Every rotation is exactly orthogonal and rational; products and matrix
updates are performed exactly. The bit lengths grow at most polynomially
in the number of rotations and their polynomial bit lengths. This proves
the claimed computational bound without an eigenvalue-separation assumption.

Let $P$ retain the $r$ entries of $D$ of largest absolute value and the
corresponding columns of $Q$. It has rank at most $r$ and a rational
orthonormal representation. Since $QDQ^{\mathsf T}$ is within $\zeta$
of $J$, the variational characterization of singular values gives
\begin{equation}
 \|J-P\|\le\sigma_{r+1}(J)+2\zeta.\label{eq:spectral-tail}
\end{equation}
This is sufficient; individual eigenvectors need not be close to any
particular exact eigenbasis.

\begin{proof}[Proof of Lemma~\ref{cor:planted}]
Put $X=\operatorname{offdiag}(J-\bar A)$ and fix two prefix rounds of
sizes $\lfloor n/4\rfloor$ and $\lfloor n/2\rfloor-\lfloor n/4\rfloor$.
Take $\zeta=n^{-4}$ in the preceding construction. Write
$E_0=\|X\|+B_0$, so $J=\bar A+X-\operatorname{diag}(\bar A)$ and
$\sigma_{r+1}(J)\le E_0$. Thus
\[
 \|J-P\|\le E_0+2\zeta,\qquad
 \|N\|\le2E_0+4\zeta,
 \quad N=\operatorname{offdiag}(J-P).
\]
For $H=\bar A-P$, we have $\rank H\le2r$ and
$\|H\|\le2E_0+2\zeta$, hence
$\|H\|_F\le\sqrt{2r}(2E_0+2\zeta)$.

The symmetric version of the random-matrix bound of
Bandeira and van Handel~\cite{BvH16}, symmetrization and convex
concentration as in Appendix~\ref{app:membership}, imply
$\|X\|\le3\sqrt{\chi n\nu}$ with probability $1-o(1)$ when
$n\nu\ge CB^2\log n$ with sufficiently large $C$.
The leading symmetrized norm constant is $2\sqrt2<3$; the
$O(B\sqrt{\log n})$ term is absorbed by increasing $C$.
Bernstein's inequality and a union bound also give, on the same event,
\[
 .99n(n-1)\nu\le\|X\|_F^2\le1.01\chi n(n-1)\nu,
 \qquad
 \sum_{s\in B_1}X_{ts}^2\ge.99|B_1|\nu\quad(t\in B_2).
\]
Here $B_0^2\le n\nu/C$, and $B\ge1$ ensures that the numerical error
$\zeta$ is negligible on this scale. Because
$N=X+\operatorname{offdiag}H$, these bounds and the triangle inequality
show, for fixed $r,\chi$ and sufficiently large $n$,
\[
 .9n\nu\le D_N\le1.1\chi n\nu,\qquad \|N\|^2\le100\chi D_N.
\]
Furthermore the correction to the sum of cross-round row norms is at most
\[
 \sqrt{|B_2|}\,\|H\|_F=O(n\sqrt{r\chi\nu})
      =o(n^{3/2}\sqrt\nu).
\]
The uncorrected sum is at least
$|B_2|\sqrt{.99|B_1|\nu}$. Consequently the gap condition holds with
$g=1/(256\sqrt\chi)$, with room to spare, and $K=100\chi$.

Let $\Lambda=\sum_{k\le r}|\lambda_k|$ for the retained diagonal entries.
They are Rayleigh quotients of $J$, so
$\Lambda\le r\|J\|\le r(\mu n+E_0)$.
Increasing $C=C(r,\chi)$ and then $n$ gives
$\Lambda\le gn\sqrt{D_N}/64$.
If $\Lambda=0$, take $P=0$ and no grid. Otherwise choose a dyadic rational
$\delta$ within a factor two below
$g\sqrt{nD_N}/(64\Lambda)$, and use Corollary~\ref{cor:spectral}.
The comparison slack is at most $gn\sqrt{D_N}/32$ and
\[
 |\mathcal Y|\le
 \left(\frac{256\Lambda}{g\sqrt{D_N}}+3\right)^r=n^{O(r)}.
\]
This proves the promised bound by the round case of
Theorem~\ref{res:counts}. Constants $g,K$ can be replaced by slightly
weaker rational bounds without affecting any conclusion.

The algorithm uses only $J$, the fixed class parameters $r,\chi$, and
its computed $D_N,\Lambda$. To define it on all inputs, use exhaustive
search when $D_N=0$; otherwise cap $\delta$ from below by $n^{-2}$.
The cap keeps the number of runs polynomial on arbitrary inputs and is
inactive on the probability event above for sufficiently large $n$.
Every resulting grid still covers all projections, so these choices
preserve exactness also outside that event. No planted mean, variance
parameter, or spectral cutoff has to be supplied.
\end{proof}

\section{Development record and proof priorities}\label{app:record}

The following record keeps the first draft auditable while the research
continues. It is intended to be removed or shortened before submission.
\begin{enumerate}
\item The bias theorem is consolidated from
\path{max2sat_biased_and_direction.tex}; the elementary random-sign theorem
from \path{max2sat_early_enumeration.tex}. Both are proved in this draft.
The bias proof counts occurrences with multiplicity and includes the exact
pair reductions and complementary-pair extension discussed in
\path{max2sat_results_and_ideas.tex} and
\path{max2sat_explicit_class.tex}. Corollary~\ref{cor:lineargap} states the
dense frequency-gap consequence directly for the original input.
Its simple-input statement also permits unit clauses, using the effective
gap; the same scale-based proof covers the anchored partitioning example
in Section~\ref{sec:applications}. The representative-optimum argument now
explicitly separates preferred values from actual child traversal order.
\item The group-count theorem comes from
\path{max2sat_redundancy_and_group_counts.tex}. The average-defect extraction
and explicit detection charge incorporate the subsequent discussion.
\item The general block certificate builds on
\path{max2sat_prebranch_constraints.tex}. The bounded-exception block
theorem generalizes the conflict-graph argument of
\path{max2sat_near_twin_pairs.tex}; its proof is included here.
\item Theorem~\ref{thm:weighted} and the round case of
Theorem~\ref{res:counts} incorporate \path{max2sat_weighted_rounds.tex}.
The former uses rational inverse-square-root weights, costing a factor
four in its conservative exponent constant. The latter and the sequential
count argument share one algorithm and one gap lemma. Max-Cut and the random
models are corollaries. The Max-Cut proof includes integer round sizes and
the bounded-degree range. Projection discovery uses rational orthogonal
rotations and a fixed rank bound rather than a cutoff requiring unknown
sampling parameters. The diagonal mean bound is charged explicitly.
The main statement is now Theorem~\ref{thm:sbm} for the stochastic block
model, including sparse and almost complete regimes; the more general
bounded-rank assertion is a supporting lemma in
Appendix~\ref{app:discovery}. The elementary specialization sets
$\nu=4v_{\min}$ and $B=B_0=\mu=2$.
\item Proposition~\ref{prop:pairmembership} retains Proposition~12 of
\path{max2sat_count_guessing.tex} in an appendix. Its stronger unit-clause
adversary and vanishing sampling probabilities distinguish it from the
arbitrary-support weighted corollary. The previously sharper almost-regular
constant is retained separately in Appendix~\ref{app:regular}.
\item The regional comparison rule and bounded-cache proposition are recent
additions. Their coverage invariant, retention policy, unweighted score
range, and variable order are explicit. The literal-gap proof now includes
an explicit count of assignments made by \textnormal{\textsc{Simplify}}, as well as the
loss of their contributions to the branching-gap sum. Additional
simplifications in other algorithms still require their own analysis.
\item The immediate research priority is a quantitative bridge from poor
cancellation to a recoverable family of approximate block comparisons or
aggregate counts. Excluding exact short relations alone is not that bridge.
\item The literature review must compare class guarantees and memory bounds,
including earlier exact algorithms on random inputs, before making any
priority or novelty claim. Constants should be optimized only after the
class statements and algorithmic model are settled.
\end{enumerate}

\paragraph{Finite checks.}
The exact-arithmetic scripts retain the following reproducible checks.
\path{verify_dense_max2sat_blocks.py}: 36 integral quadratic instances,
638 conditional block certificates, and 1,280 affine and regional
comparisons, checked against exhaustive optimization. The objectives
admit unweighted Max 2-SAT encodings with repeated clauses and units.
Earlier enumeration checks cover 44 instances, 1,408 prefix counts and
65,024 assignment identities.
\path{verify_dense_max2sat_random_models.py}: 19 pair distributions,
36 formulas with units, 39,936 count-shift identities, 50,688 comparisons,
and optimum preservation for the original and adjusted gains.
\path{verify_dense_max2sat_rounds.py}: 18 instances, 2,688 round-sum
identities, 17,280 projection-error bounds, optimum preservation, and
36 rational Jacobi rotations.
\path{verify_dense_max2sat_sbm.py}: 34 pair moments, 18 block-mean ranks,
8,640 supplied-group gain identities over 1,344 assignments,
35 optimum-survival cases, and seven clique optimizations with fields.
Its models include unequal groups and dense, sparse, almost complete,
and deterministic probabilities. Separately, the floating-point script
\path{explore_weighted_rounds.py} loses no optimum on 200 weighted
instances and 120 instances each with count and spectral grids.
These finite consistency checks do not establish the asymptotic bounds.

\begin{thebibliography}{99}
\small\raggedright
\setlength{\itemsep}{3pt}\setlength{\parsep}{0pt}
\bibitem{Williams}
R. Williams. A new algorithm for optimal 2-constraint satisfaction and its
implications. \emph{Theoretical Computer Science}, 348(2--3):357--365, 2005.
\href{https://www.cs.cmu.edu/~ryanw/2-csp-final.pdf}{Author manuscript}.

\bibitem{KK}
A. S. Kulikov and K. Kutzkov. New bounds for MAX-SAT by clause learning.
\emph{Computer Science--Theory and Applications (CSR 2007)},
LNCS 4649, pp. 194--204, 2007.
\href{https://www.pdmi.ras.ru/pdmi/publications/2007}{Institutional publication record}.

\bibitem{GK}
A. Golovnev and K. Kutzkov. New exact algorithms for the 2-constraint
satisfaction problem. \emph{Theoretical Computer Science}, 526:18--27, 2014.
\href{https://www.golovnev.org/papers/max2csp.pdf}{Author manuscript}.

\bibitem{Hastings}
M. B. Hastings. Local maxima and improved exact algorithm for MAX-2-SAT.
2016. \href{https://arxiv.org/abs/1610.07100}{arXiv:1610.07100}.

\bibitem{BIPS}
P. Beame, R. Impagliazzo, T. Pitassi and N. Segerlind.
Memoization and DPLL: formula caching proof systems.
\emph{Proceedings of the 18th IEEE Conference on Computational Complexity},
pp. 225--236, 2003.
\href{https://homes.cs.washington.edu/~beame/papers/caching.pdf}{Author manuscript}.

\bibitem{CBGV12}
N. Cesa-Bianchi, C. Gentile, F. Vitale and G. Zappella.
A correlation clustering approach to link classification in signed networks.
\emph{Proceedings of the 25th Annual Conference on Learning Theory},
PMLR 23, pp. 34.1--34.20, 2012.
\href{https://proceedings.mlr.press/v23/cesa-bianchi12.html}{Conference article}.

\bibitem{MM03}
R. Miyashiro and T. Matsui.
Semidefinite programming based approaches to the break minimization problem.
Technical Report METR 2003--28, University of Tokyo, 2003.
\href{https://www.keisu.t.u-tokyo.ac.jp/data/2003/METR03-28.pdf}{Technical report}.

\bibitem{EHS02}
T. Erlebach, A. Hall and T. Schank.
Classifying customer-provider relationships in the Internet.
TIK Report 145, ETH Z\"urich, 2002.
\href{https://doi.org/10.3929/ethz-a-004403617}{Institutional report}.

\bibitem{BvH16}
A. S. Bandeira and R. van Handel.
Sharp nonasymptotic bounds on the norm of random matrices with independent
entries. \emph{Annals of Probability}, 44(4):2479--2506, 2016.
\href{https://arxiv.org/abs/1408.6185}{Author manuscript}.

\bibitem{BLM13}
S. Boucheron, G. Lugosi and P. Massart.
\emph{Concentration Inequalities: A Nonasymptotic Theory of Independence}.
Oxford University Press, 2013.

\bibitem{Tal95}
M. Talagrand.
Concentration of measure and isoperimetric inequalities in product spaces.
\emph{Publications Math\'ematiques de l'IH\'ES}, 81:73--205, 1995.
\href{https://pmihes.centre-mersenne.org/articles/10.1007/BF02699376/}{Journal article}.

\bibitem{RV13}
M. Rudelson and R. Vershynin. Hanson--Wright inequality and sub-gaussian
concentration. \emph{Electronic Communications in Probability}, 18(82):1--9, 2013.
\href{https://arxiv.org/abs/1306.2872}{Author manuscript}.

\bibitem{BH12}
A. E. Brouwer and W. H. Haemers. \emph{Spectra of Graphs}. Springer, 2012.
\href{https://homepages.cwi.nl/~aeb/math/ipm/}{Authors' book page}.

\bibitem{TY19}
K. Tikhomirov and P. Youssef. The spectral gap of dense random regular graphs.
\emph{Annals of Probability}, 47(1):362--419, 2019.
\href{https://arxiv.org/abs/1610.01765}{Author manuscript}.

\bibitem{Abbe18}
E. Abbe. Community detection and stochastic block models: recent developments.
\emph{Journal of Machine Learning Research}, 18(177):1--86, 2018.
\href{https://jmlr.org/papers/v18/16-480.html}{Journal article}.

\bibitem{NoteCounts}
\emph{Count guessing: aligned structure plus a pseudo-random remainder}.
Companion working note, October 2026.
Local source: \path{max2sat_count_guessing.tex}.
\end{thebibliography}
\end{document}
